% Gerado por tools/tei2latex.py a partir de entrada.xml. Regenerar, não editar.

%%>>> parte-1

%%>>> cap-4

%%>>> t122

8) Divisão da razão: --- A \emph{razão} divide-se em \emph{theoretica} e \emph{pratica}. --- A \emph{theoretica}, ou \emph{especulativa}, limita-se á percepção \pg{170}  da verdade. A \emph{pratica} applica á operação a verdade percebida. --- A razão \emph{theoretica} e á \emph{pratica} não são duas faculdades distinctas \pg{171} , mas duas fuzegoes da mesma faculdade; pois que o objecto da razão \emph{theoretica} é o mesmo que o da \emph{pratica}, e é \pg{172}  apenas \emph{accidental} que a razão \emph{pratica} applique á operação a verdade, que a razão \emph{theoretica} considera em si mesma \footnote{\emph{Kant} (\emph{Critica da razão pura}) ensinou erradamente que á razão \emph{theoretica} e a. \emph{pratica} são duas faculdades distinctas e diversas.\par \emph{A} razão \emph{pratica}, applicando ás operações os principios universaes, forma juizos particulares ácerca de certos e determinados actos. Para isso, emprega um syllogismo. \emph{A} premissa \emph{maior} d'este syllogismo contém um principio universal; a premissa \emph{menor} indica que um certo e determinado sujeito está contido no sujeito do principio universal; a \emph{conclusão} infere que ao sujeito particular compete a propriedade, que se attribue ao sujeito universal. Exemplo: \emph{Os filhos devem respeitar seus paes}; \emph{Pedro é filho de Paulo}; \emph{logo}: \emph{Pedro deve respeitar Paulo} (\emph{Sum}. \emph{Theol}., p. I, q.. 79, a. 11).\par Com a razão \emph{pratica} está intimamente ligada a \emph{synderesis}. --- A \emph{synderesis} é o \emph{habito natural dos principios moraes} ou \emph{axiomas praticos}. Assim como a razão \emph{theoretica} discorre ácerca de assumptos \emph{especulativos}, assim a razão \emph{pratica} discorre ácerca de assumptos \emph{praticos}. Ora o discurso, como todo e qualquer movimento, deve ter o seu começo e o seu termo num ponto immoyel; e este ponto immovel deve ser constituido pelos primeiros principios. Deve, portanto, \emph{o} homem possuir os primeiros principios tanto da ordem \emph{especulativa} (todo o effeito exige uma causa), como da \emph{pratica} (deve fazer-se o bem; deve evitar-se o mal) Ora os primeiros principios da ordem \emph{especulativa} não pertencem a uma faculdade especial, mas derivam de uma disposição natural, que a intelligencia tem para formar juizos relativos aos principios \emph{theoreticos} ou \emph{especulativos}, \emph{e} que se chama \emph{habito dos principios}. Logo tambem os primeiros principios da ordem \emph{pratica} não derivam de uma faculdade especial, mas de uma disposição ou inclinação natural que a intelligencia tem para formar juizos relativos aos principios \emph{praticos}, e que se chama \emph{synderesis} (\emph{Sum}. \emph{Theol}., p. I, q. 79, a. 12). --- O habito dos primeiros principios \emph{praticos}, como tambem o dos primeiros principios \emph{especulativos}, diz-se \emph{natural}, emquanto se adquire sem esforço e por uma simplicissima abstracção da. intelligencia agente, \emph{e} não se corrompe nem \$e destroe. (\emph{Sun}. \emph{Theol}., 1-2, q. 58, a. 1).\par \emph{A razão} costuma tambem dividir-se em \emph{superior} e \emph{inferior}. A \emph{superior} julga das coisas, segundo as verdades eternas, necessarias e immutayeis da intelligencia divina, que é o exemplar de todos os entes creados. A \emph{inferior} julga das coisas, segundo as verdades temporaneas, accidentaes e mutayeis, que se conhecem pela consideração das coisas creadas. Assim a razão \emph{superior} afirma que se deve obedecer ao rei, porque \textsc{DEUS} o \emph{quer}; e á \emph{inferior} diz que se deve obedecer ao rei, porque a \emph{felicidade da nação o exige}.\par Como se vê, a divisão da \emph{razão} em \emph{superior} e \emph{inferior} basêa-se na qualidade do seu \emph{objecto}, que é a verdade; assim como a divisio da mesma faculdade em \emph{theoretica} e \emph{pratica} se funda no \emph{fim}, para que consideramos a verdade.\par \emph{A} razão \emph{inferior} e á \emph{superior} não são duas faculdades, mas duas funcções\pgn{173}da mesma faculdade. Por quanto, as verdades eternas, necessarias e immutaveis, que são o objecto da razão \emph{superior}, e as verdades temporaneas, contingentes e mutaveis, que são o objecto da razão \emph{inferior}, embora sejam \emph{materialmente} muitas e diversas, todavia constituem um só objecto \emph{formal}; pois todas são attingidas emquanto \emph{émomateriaes} ou \emph{universaes}, e a unidade ou multiplicidade das faculdades depende, como dissemos, da unidade ou multiplicidade do seu objecto \emph{formal} (\emph{Sum}. \emph{Theol}. p. l, q. 79, a. 9).\par A razão \emph{superior}, ainda que seja mais nobre e mais elevada do que a \emph{inferior} (pois aquella julga das coisas por criterios mais altos do que esta), comtudo, na sua operação, é \emph{posterior} á razão \emph{inferior}; pois não podemos elevar-nos \& consideração do eterno e do immutayvel, senão pela consideração do temporaneo e do mutavel.\par \emph{Cousin} affirma que a razão \emph{superior}, sendo o meio pelo qual a razão divina communica com a razão humana, existe fôra de nós. --- Os \emph{ontologistas} dizem que a razão \emph{superior} é a propria razão \emph{absoluta} --- a divina.\par Estas opiniões são falsas. A razão \emph{superior} e a \emph{inferior} não são duas faculdades distinctas, mas indicam apenas duas funcções distinctas, embora não diversas \emph{formalmente}, de uma e a mesma faculdade. Sendo assim, é claro que a razão \emph{superior} não está acima de nós, mas está dentro de nós, é \emph{pessoal} e \emph{propria} de cada um de nós, como \emph{pessoal} e \emph{propria} de cada um de nós é a razão \emph{inferior}. :\par Objectam: O homem percebe o immutavel. Ora o immutayel não póde ser percebido pela razão humana, porque esta é mutavel. Logo deve ser percebido por uma \emph{rasdo superior} á nossa, e por isso a \emph{rasdo superior} esta \emph{fóra} ou \emph{acima} de nos.\par Respondemos que não só \emph{não repugna}, mas até é \emph{necessario} que a nossa alma perceba o \emph{immutavel} por meio de uma faculdade \emph{mautavel}. --- \emph{Não repugna}; pois que a nossa alma, na sua percepção, não \emph{cria} as propriedades ou relações das coisas, limita-se a investigal-as e descobril-as. --- \emph{E necessario}; visto que o meio, pelo qual conhecemos o immutavel, deve ser \emph{intrinseco} á nossa alma, pois o conhecimento é uma operação \emph{immanente}, e tudo o que está em a nossa alma é mutavel.}.

b) \emph{A} razão e a intelligencia não são duas faculdades distinctas, mas são dois nomes da mesma faculdade.

a) \emph{Não são duas faculdades distinctas}. A \emph{vazão} e a \emph{intelligencia} seriam duas faculdades distinctas, se o objecto e o acto de uma fossem \emph{formalmente} distinctos do objecto e do acto da \pg{173}  outra. Ora os objectos e os actos da \emph{razão} e da \emph{intelligencia} não são \emph{formalmente} distinctos. --- Não são distinctos os \emph{objectos} ; pois tanto a \emph{razão} como a \emph{intelligencia} têm por objecto o \emph{immaterial} e o \emph{universal}. --- Não são distinctos os \emph{actos}; porque o acto da \emph{intelligencia} é a \emph{intuição} e o da \emph{razão} é o \emph{raciocinio}, ora a intuição e o raciocinio não são operações \emph{especificamente} diversas. Logo a \emph{intelligencia} e a \emph{razão} não são duas faculdades distinctas \footnote{O argumento, tirado da identidade do \emph{objecto} da \emph{intelligencia} e da \emph{razão}, póde expor-se d'este outro modo. --- A \emph{intelligencia} considera os \emph{principios} dos raciocinios; e a \emph{razão} tira as \emph{conclusões}, contidas naquelles principios. Ora a \emph{razão} não poderia tirar as \emph{conclusões}, contidas em certos \emph{principios}, se não conhecesse tambem os proprios \emph{principios}. Se a \emph{razão}, além das conclusões, conhece tambem os \emph{principios} dos raciocinios, o objecto da \emph{razão} é o mesmo que o da \emph{intelligencia}. Se o objecto é o mesmo, a \emph{razão} e a \emph{intelligencia} não são duas faculdades distinctas.}.

\emph{b}) \emph{Mas são dois nomes da mesma faculdade}. --- E' um corollario da proposição” antecedente. Se a \emph{intelligencia} e a \emph{razão} não são duas faculdades distinctas, é necessario concluir que são apenas dois nomes da mesma faculdade, isto é, que uma e a mesma faculdade intellectual, --- emquanto percebe ?\emph{mmediatamente} a verdade, chama-se \emph{intelligencia}, --- e, emquanto de uma verdade conhecida deduz outra verdade não conhecida, diz-se \emph{razão}. Logo são dois nomes da mesma faculdade \footnote{\emph{Kant} opinou que a \emph{intelligencia} e a \emph{razão} são duas faculdades \emph{especificamente} distinctas. Outros escriptores e naturalistas modernos fazem a mesma distinegão, attribuindo a \emph{intelligencia} aos animaes superiores e reservando a \emph{razão} para o homem. \emph{Flourens} diz: « A \emph{intelligencia} propriamente dicta começa com os animaes superiores; a \emph{razão} é propriedade exclusiva do homem » (\emph{De la vie et de l'intellig}., p. I, ch. VII)\par A opinião de \emph{Kant}, como a dos outros, é falsa. Por quanto, como provámos, a faculdade \emph{intellectual} é \& que percebe o \emph{immaterial} ou \emph{universal}. Ora o \emph{immaterial} ou \emph{universal} póde perceber-se --- ou \emph{directa} e \emph{immediatamente}, isto é, em si mesmo, --- ou \emph{indirecta} e \emph{mediatamente}, isto é, por meio de uma verdade já conhecida. Emquanto percebe o immaterial \emph{directa} ou \emph{immediatamente}, a faculdade \emph{intellectual} chama-se \emph{intelligencia} ; emquanto percebe o immaterial \emph{indirecta} ou \emph{mediatamente}, diz-se \emph{razão}.\par A \emph{intelligencia}, pois, e a \emph{razão} não denotam uma distincção de faculdades, mas apenas uma differença no modo de conhecer. A mesma faculdade perceptiva, emquanto conhece por simples \emph{intuição}, diz-se \emph{intelligencia}, e, emquanto conhece pelo \emph{discursa} ou \emph{raciocinio}, chama-se \emph{razão}. O \emph{raciocinio} está para a \emph{intuição}, como o \emph{movimento} para o \emph{repouso}. Ora é claro que o \emph{movimento} e o \emph{repouso} não são actos de duas potencias ou faculdades distinctas, mas são dois actos de uma e a mesma potencia; porque é sempre em virtude da mesma natureza que um ente se move para um logar e nelle descança. Logo a intelligencia e á razão não são duas faculdades. --- Os Anjos apprehendem, simplesmente e sem discurso, a verdade das coisas, e por isso são chamados \emph{intelligencias}; os homens chegam ao conhecimento da verdade pelo discurso ou raciocinio, e por isso são chamados \emph{raciondes}.\par Mas, como dissemos, o movimento procede sempre de um ponto immovel e num ponto immoyel acaba. Por isso, a \emph{razão}, quando começa o seu discurso de invenção ou inquisição, parte de verdades simples ou intuitivamente percebidas, que são os \emph{primeiros principios}, e depois de ter chegado ao termo, analysa a conclusão e resolve-a nos mesmos principios, e assim julga da rectidão do discurso. D'este modo, a razão refere-se á intelligencia, como ao seu principio e ao seu termo. (\emph{Sum}. \emph{Theol}, p. I, q. 719, a. 8; --- QQ. DD. \emph{De verit}., q. 15, a. 1).\par S. Thomaz nota tambem que \emph{éntellecto} e \emph{intelligencia}, se se tomam em sentido \emph{lato}, significam e sao uma e a mesma coisa, --- mas, se se tomam em sentido \emph{rigoroso}, distinguem-se como á \emph{potencia} e o \emph{acto}, porque á \emph{intelligencia} é o \emph{acto} do \emph{intellecto} (\emph{Sum}. \emph{Theol}., p. L, q. 79, a. 10).}.

%%>>> t123

\tese{123}{Consciencia.}
\emph{Consciencia} é a propria intelligencia, que, pela reflexão, percebe a \emph{existencia actual} das suas operações \pg{169} \pg{174} \pg{175}  \emph{espirituaes}. --- Que á intelligencia, por ser uma faculdade \emph{immaterial} ou \emph{espiritual}, possa, pela reflexão, perceber a actual existencia das proprias operações \emph{intellectuaes}, é uma verdade demonstrada pela razão e confirmada pela experiencia \footnote{É um facto innegavel que a nossa intelligencia não só percebe os objectos externos, nias, depois de ter conhecido os objectos externos, é tambem capaz de reflectir sobre si mesma, sobre os seus proprios actos, de modo que, ficando sempre uma e a mesma faculdade, se torna, simultaneamente, \emph{sujeito intelligente} e \emph{objecto entendido}. E' a \emph{reflexão perfeita}, dote exclusivo de uma faculdade \emph{immaterial}, que, por ser simplicissima, pode dobrar-se sobre si mesma, e como que compenetrar-se, havendo assim \emph{identidade} entre o sujeito cognoscente e o objecto conhecido. --- Esta reflexão diz-se \emph{subjectiva}, e distingue-se da \emph{objectiva}, a qual consiste em repetir os actos de attenção sobre um objecto externo, anteriormente percebido. --- A percepção da nossa vida espiritual, ou o conhecimento das nossas operações espirituaes, por meio da \emph{reflexão perfeita} e \emph{subjectiva}, é o acto \emph{proprio} da \emph{consciencia intellectual}, ou, simplesmente, \emph{consciencia}. Diz S. Thomaz: « Est alius intellectus, scilicet humanus, qui nec est suum intelligere, nec sui intelligere est objecium primum ipsa ejus essentia, sed aliquid \emph{extrinsecum}, scilicet \emph{natura materialis rei}. Et ideo id quod \emph{primo} cognoscitur ab intellectu humano est hujusmodi objectum; et \emph{secundario} cognoscitur \emph{ipse actus}, quo cognoscitur objectum, et per actum cognoscitur ipse intellectus, cuius est perfectio ipsum intelligere. Et ideo Philosophus dicit (\emph{De anima}, lib. II, text. 58) quod objecta praecoguoscuntur actibus, et actus potentiis » (\emph{Sum}. \emph{Theol}, p. I, q. 87, a. 9).\par D'onde se deduzem as differengas entre a \emph{consciencia sensitiva}, que se chama propriamente \emph{sentido commun}: e convém tambem aos animaes, e a \emph{consciencia intellectual}, que é exclusivamente propria do homem. A consciencia \emph{sensitiva} é uma faculdade \emph{organica}; \& \emph{intellectual} é uma faculdade \emph{inorganica}. \emph{A sensitiva} tem por objecto as \emph{operações sensitivas} e o \emph{sujeito sensitivo agente}; a \emph{intellectual} tem por objecto as \emph{operações espirituaes} e o proprio \emph{espirito agente}. \emph{A sensitiva} não é capaz de \emph{reflexão perfeita}, porque é uma faculdade \emph{extensa}, e por isso se uma sua parte pode dobrar-se sobre outra, não pode comtudo dobrar-se sobre si mesma, e muito menos o todo pode reflectir sobre o todo (oppõe-se a lei da impenetrabilidade dos corpos); ao passo que a \emph{intellectual}, sendo uma faculdade simplicissima, pode duplicar-se em sujeito e objecto, é capaz de uma perfeita reflexão. Diz S. Thomaz: u Nullus sensus seipsum cognoscit, ac suam operationem; visus enim non videt seipsum, nec videt se videre. Intellectus autem cognoscit seipsum, et cognoscit se intelligere » (\emph{Cont}. \emph{gent}., l. II, c. 66; --- C£. \emph{Sum}. \emph{Theol}, p. I, Eq. 14, a. 2 ad 1).\par Alguns philosophos negaram a possibilidade da \emph{reflexão subjectiva}, e por isso a existencia da \emph{consciencia intelectual}, --- \emph{Locke} escreveu que « a nossa -intelligencia é semelhante aos olhos, que vêm os objectos externos, mas não se vêm a si mesmos ». --- \emph{Comte} disse que a doutrina relativa á \emph{consciencia} - \emph{intellectual} se funda num \emph{dualismo absurdo}, porque é absurdo que o nosso espirito possa ser, ao mesmo tempo, \emph{sujeito intelligente} e \emph{objecto entendido}. --- Outros escriptores identificam a \emph{consciencia} com o \emph{conhecimento intellectual},\pgn{176}pelo facto de haver em ambos estes actos o \emph{dualismo}, isto é, a distincção entre o \emph{sujeito}, que conhece, e o \emph{objecto}, que é conhecido.\par Estas opiniões são falsas. E' falsa a opinião de \emph{Locke}. Por quanto, esse escriptor identifica a nossa intelligencia com as faculdades sensitivas; o que, como vimos, é completamente falso. --- E' falsa a opinião de \emph{Comte}. Desde o momento que se admitte, como não se pode deixar de admittir, que a intelligencia é uma faculdade espiritual, isto é, não ligada \& materia, o tal \emph{dualismo} não só não repugna, mas é muito natural. Para a \emph{distincção} nao é necessaria a \emph{separação}, nem a \emph{diversidade}. --- E' falsa a opinião dos que confundem a \emph{consciencia} com o \emph{conhecimento intellectual}. Na verdade, o \emph{conhecimento} tem por termo os objectos diversos do \emph{eu}, ou o \emph{eu} considerado emquanto é uma \emph{essencia} ou \emph{natureza}; ao passo que a \emph{consciencia} tem por termo os \emph{actos} do \emph{eu}, e o proprio \emph{eu} mas considerado emquanto é o \emph{sujeito agente}. Uma coisa é \emph{conhecer}; outra é \emph{conhecer que se conhece}. Uma e a mesma é a faculdade que opera, mas os actos são diversos. Diz S. Thomaz: « Ipsum intelligere humanum non est actus et perfectio naturae intellectae, ut sic possit uno acto intelligi natura rei materialis ot ipsum intelligere, sicut uno acto intelligitur res cum sua perfectione. Unde \emph{alius est actus quo intellectus intelligit lapidem}, et \emph{alius} est \emph{actus quo intelligit se intelligere lapidem} » (\emph{Sum}. \emph{Theol}, p. I, p. 81, a. 3 ad 2).\par Do que deixamos dicto deduz-se que a \emph{consciencia}, tomada em sentido rigoroso, é uma faculdade exclusivamente propria dos seres intelligentes. Os animaes irracionaes poderão mostrar, nos seus actos, uma certa apparencia de reflexão' perfeita; mas, se analysarmos esses actos, veremos que reflexão perfeita não ha, nem pode haver. A sua alma, ligada com a materia para existir e operar, não se pode desligar da materia para \emph{reflectir} sobre si mesma e os seus actos. --- Os seus phantasmas representam as sensações e o proprio organismo, mas não se representam a si mesmos, nem o principio, de que derivam. Por isso o animal não pode attribuir a si mesmo o que percebe pelas suas faculdades. Quem não se pode dobrar sobre si mesmo, não pode ter consciencia do que nelle se realiza; quem não diz \emph{eu}, não pode dizer \emph{meu}. O animal, pelo acto da imaginação, percebe as suas sensações, mas não as percebe como \emph{suas}; percebe-as como coisas \emph{intimus}, mas não as reconhece como \emph{proprias}. Os actos da percepção são seguidos pelos actos do appetite; mas esse nexo o animal nem o conhece nem o determina livremente, Concluimos que, entre todas as creaturas visiveis, \emph{s} o homem tem \emph{consciencia} dos seus actos, da sua identidade, da sua personalidade. Admiravel prerogativa esta, que imprime na nossa vida intellectual uma semelhança, ainda que remotissima, da propria vida de Deus! Assim como a intelligencia do Pae, pelo conhecimento da essencia divina, gera o Verbo, que é imagem substancial do seu Principio, assim tambem a intelligencia humana, quando, pela reflexão, se percebe a si mesma, produz uma imagem, um verbo, que traz gravada a semelhança do principio, de que procede e que é a propria intelligencia. (Cf. S. Thom., QQ. DD., q. II \emph{de Pot}., a. 1).}.

Vamos indicar o seu \emph{objecto}, as suas \emph{especies}, e provar a sua \emph{identidade real} com a intelligencia.

a) Objecto da consciencia, --- O objecto da \emph{consciencia é directo} e \emph{indirecto}. --- O objecto \emph{directo} é constituido pelas nossas \emph{operações} \pg{177} \emph{ espirituaes}, \emph{actualmente} presentes. Na verdade, o objecto \emph{directo} de uma faculdade é o que ella \emph{primeiramente} \textbf{percebe}. Ora a consciencia \emph{primeiramente} percebe as nossas \emph{operações espirituaes}, que por si e immediatamente estão ao alcance da reflexão. --- O objecto \emph{indirecto} é o proprio \emph{sujeito agente}. Por quanto, á alma, por meio da \emph{consciencia}, percebe as operações espirituaes não só como presentes, mas tambem como \emph{proprias}, Ora não poderia conhecer as operações como \emph{proprias}, se não se conhecesse, ao mesmo tempo, como principio das mesmas operações, e por isso não attingisse o proprio \emph{sujeito agente} \footnote{As operações \emph{jatellectuaes}, de que temos consciencia, não são apenas os actos da \emph{intelligencia}, mas tambem os da \emph{vontade}, pois que tambem \textbf{estes}, por serem espirituaes, são intelligiveis (\emph{Sum}. \emph{Theol}, p. I, q. \emph{Vt}, a. 4)..\par Dissemos --- as operações \emph{intellectuaes}; pois que as operações \emph{sensitivas}, sendo materiaes, não são percebidas pela \emph{consciencia}, que é uma faculdade immaterial, mas pelo \emph{sentido commum}. ,\par Devemos notar: --- 19) pela \emph{consciencia}, conhecemos o proprio \emph{sujeito agente}, não na sua \emph{intima natureza}, isto é, não como uma substancia espiritual, mas como \emph{principio} e \emph{causa} das operações ?\emph{nfellectuaes} ; --- 2.9) o objecto \emph{proprio} da consciencia é o acto \emph{interior}, não emquanto este é \emph{objectivo}, isto é, não emquanto este representa ou pode representar um \emph{objecto}, mas emquanto é \emph{subjectivo}, isto é, emquanto é um \emph{phenomeno interior}, ou uma \emph{modificação} da faculdade intellectual, porque a percepção do \emph{objecto}, representado pelo acto intellectual, pertence á \emph{intelligencia}.\par \emph{Vacherot}, com os \emph{Cartesianos}, opinou que, por meio da \emph{consciencia}, conhecemos a propria \emph{natureza} da nossa alma. --- Esta opinião é falsa. Por quanto á consciencia só attesta o que se passa na alma; mas isto não basta para conhecermos a \emph{natureza} da alma, pois que, para isso, é necessario fazer um raciocinio, que do conhecimento da natureza das operações da alma se \textbf{eleve} ao conhecimento da \emph{natureza} da propria alma. --- De mais, todos têm a consciencia de si mesmos, mas nem todos conhecem a \emph{natureza} da sua alma.\par \emph{Hamilton} disse que, por meio da consciencia, percebemos o mundo \textbf{exterior}. --- O erro é manifesto. O mundo exterior só pode ser e é percebido pelos sentidos \emph{externos}, que estão ao contacto immediato com as coisas visiveis; mas o que se passa no intimo do nosso espirito só pode ser percebido por uma faculdade espiritual, como é a consciencia.}.

b) Divisão da consciencia, --- A \emph{consciencia} divide-se em \emph{habitual} e \emph{actual}. --- Consciencia \emph{habitual} é a que \emph{actualmente} não exerce o seu acto, mas tem grande facilidade em o exercer. --- Qonscieneia \emph{actual} é o proprio acto da consciencia, com que a alma percebe as suas operações e o sujeito pensante \footnote{Alguns escriptores negam á divisão da \emph{consciencia} em \emph{habitual} e \emph{actual}, dizendo que temos \emph{continuamente} uma consciencia \emph{actual} das \textbf{operações} \emph{intellectuaes}. --- Esta opinião é falsa. Para que o exercicio da consciencia fosse sempre \emph{actual}, seria necessario que o acto do pensamento fosse continuo, ou que tivessemos constantemente consciencia do nosso \textbf{pensamento}. Ora é falso que o acto do pensamento seja \emph{continuo}; porque á \textbf{nossa intel}-\par 12 \textbar{} Vor. II.\pgn{178}ligencia, como provámos, é uma faculdade \emph{potencial}, que sempre \emph{yóde} entender, mas que nem sempre exerce o seu acto; --- como tambem é falso que tenhamos constantemente consciencia do nosso pensamento, porque muitas vezes pensamos, sem que todavia reflictamos sobre o nosso pensamento.\par A consciencia costuma dividir-se tambem em \emph{espontanea} e \emph{reflexa}, \emph{clara} e \emph{obscura}, \emph{distincta} e \emph{confusa}, \emph{individual} e \emph{social}. --- Digamos alguma cois\& \&ceroa d'estas especies.\par \emph{a}) \emph{Consciencia espontanea} e \emph{veflexa}. --- \emph{Espontanea} é a simples percepção dos phenomenos psychicos, sem que se distinga o sujeito, que percebe, da operação, que é percebida. --- \emph{Reflexa} é a percepção dos phenomenos psychicos, e tambem da nossa personalidade, emquanto é causa d'esses phenomenos e é distincta d'elles. --- A \emph{espontanea} chama-se tambem \emph{actual}, ou \emph{de primeiro grau}, ou simplesmente \emph{percepção}; e á \emph{reflexa} diz-se tambem \emph{substancial}, ou \emph{de segundo grau}, ou \emph{apercepção} (Leibnitz). --- Outros dizem que a \emph{consciencia} --- é \emph{espontanea}, quando se excita por si mesma ao acto, --- e é \emph{reflexa}, quando é provocada pelo appetite, sensitivo ou intellectual.\par \emph{b}) \emph{Consciencia clara e obscura}, --- Esta divisão, como tambem a seguinte, refere-se aos diversos graus de consciencia, emquanto os nossos actos internos, ou phenomenos psychicos, podem ser mais ou menos advertidos. --- E \emph{clara}, quando adverte plenamente um phenomeno interior, sem perceber a distinogüo entre os varios estados psychicos, \emph{presentes} e \emph{anteriores}. --- E' \emph{obscura}, quando nem a advertencia dos factos internos é plena. --- A \emph{obscura} costuma chamar-se tambem \emph{subconsciencia}, a qual tem, naturalmente, varios graus, desde o vertice, que se approxima da consciencia \emph{clara}, até ao infimo limite, que se perde quasi na \emph{inconsciencia}, isto é, na \emph{privação de consciencia}. --- Para alguns escriptores a \emph{subconsciencia} é uma disposição psychica, innata ou adquirida, pela qual a nossa consciencia é capaz de concretisar, \emph{e} muitas vezes coneretisa, num facto psychico certos principios, theoreticos ou praticos, ou certos sentimentos moraes, que estavam na consciencia num estado \emph{germinal} e \emph{latente}.\par \emph{c}) \emph{Consciencia distincta e confusa}. --- E) \emph{distincta}, quando, além de uma clara ou plena advertencia do estado subjectivo, se conhece tambem a differenga entre os varios estados psychicos, presentes e anteriores, --- E' \emph{confusa}, quando, embora se tenha uma percepção clara ou plena do estado subjectivo, todavia não se distingue a differenca entre o estado presente e os anteriores. Por isso a consciencia \emph{clara} pode ser \emph{confusa}, como pode ser \emph{confusa} uma idea \emph{clara}. --- Alguns escriptores dão á consciencia \emph{obscura} e \emph{confusa} o nome de \emph{subconsciencia}, reservando para a consciencia \emph{clara} e \emph{distincta} o \textbf{nome} de \emph{consciencia}. --- \emph{Leibnitz} (\emph{Nouveau essai}, lib. II, XXIX) distingue\pgn{179}\emph{quatro} graus nos estados da consciencia: \emph{claros} e \emph{distinctos}, \emph{claros} e \emph{confusos}, \emph{surdos}, \emph{mais do que surdos}.\par \emph{d}) \emph{Consciencia individual e social}. --- \emph{A individual} é a percepção que o homem ter das suas operações espirituaes, e da propria personalidade, emquanto é o sujeito substancia], immutavel e immanente dos mutaveis phenomenos psyehicos. Esta percepção diz-se tambem \emph{autoconsciencia}, \emph{suiconsciencia}, \emph{consciencia} do \emph{eu}. E' pela consciencia \emph{individual} que o homem, atravez dos mutaveis e diversos phenomenos psychicos, percebe a \emph{unidade} e a \emph{continuídade} da sua vida, e reconhece o proprio \emph{eu} substancial, sempre identico e permanente. Note-se que o \emph{eu} é todo o nosso ser, é á nossa pessoa, é o composto de alma e de corpo. --- A \emph{social}, dieta impropriamente \emph{consciencia}, é a maneria de pensar ou o estado psychico da sociedade, que se manifesta exteriormente nos actos, nas instituições, nas conversas, nos escriptos, etc. Diz-se tambem \emph{collectiva}, e segue as varias formas de sociedade: d'ahi á consciencia \emph{domestica}, \emph{civil}, \emph{religiosa}, etc. \emph{A}" \emph{social} reduz-se a consciencia \emph{scientifica}, propria das pessoas sabias, que se desenvolve sob as irradiações da sciencia e se distingue da \emph{vulgar}, propria do povo.}.

c) A consciencia não é realmente distincta da intelligencia, --- A \emph{consciencia} seria uma faculdade realmente distineta da ?\emph{ntelligencia}, se o objecto \emph{formal} d'aquella fosse distincto do d'esta. Mas não é assim. Na verdade, a \emph{consciencia} tem por objecto as operações espirituaes e o sujeito pensante. Ora as operações espirituaes e o sujeito pensante são coisas immateriaes, e tudo o que é immaterial é intelligivel, e tudo o que é intelligivel é objecto da intelligencia. --- De mais, a intelligencia, por ser espiritual, póde reflectir sobre si mesma, e perceber as proprias \pg{179}  percepções e por isso o proprio sujeito agente. Logo a \emph{consciencia} não é realmente distineta da \emph{intelligencia} \footnote{\emph{Reid}, \emph{Dugald-Stewart} e \emph{Balmes} pensaram que a \emph{consciencia} é uma faculdade realmente distincta da intelligencia. Repetimos que a \emph{consciencia} é á mesma energia intellectual, emquanto reflecte sobre si mesma, e se conhece por si mesma.\par Muitos escriptores modernos erram ácerca da faculdade, por meio da qual a alma percebe as proprias operações. Admittem uma faculdade, \emph{especificamente} distineta da intelligencia e dos sentidos, á qual dão o nome de \emph{sentido intimo} ou \emph{consciencia} e a que attribuem a percepção \emph{directa} de todas as operações humanas, não só intellectuaes, mas tambem sensitiyas.\par Repetimos --- 1.9) que, se a percepção das operações immateriaes e intellectaaes exige uma faculdade immaterial, como é a \emph{consciencia}, a percepção das operações sensitivas exige uma faculdade material que é o \emph{sentido commam}; porque objectos \emph{formal} e \emph{especificamente} diversos exigem faculdades \emph{formal} e \emph{especificamente} diversas: --- 2.9) que a \emph{consciencia} é uma faculdade \emph{realmente} distincta dos sentidos, mas não da \emph{intelligencia}.\par A consciencia, de que até agora temos fallado, não é a \emph{moral}, mas a \emph{psychologica}. --- Para melhor comprehender esta distincção, devemos advertir que a nossa alma póde considerar as suas operações espirituaes com o duplice fim --- ou de perceber a sua existencia, --- ou de conhecer se estão, ou não, conformes com as regras dos costumes. D'ahi a distineção da consciencia em \emph{psychologica} e \emph{moral}. \emph{A psychologica} refere-se á existencia das operações; a \emph{moral} versa ácerca da conveniencia das mesmas operações com as regras dos costumes. --- \emph{A} consciencia, de que se occupa a Metaphysica, é a \emph{psychologica}; da consciencia \emph{moral} trata a Philosophia moral.}.

%%>>> t124

\tese{124}{Memoria intellectual.}
\emph{Memoria intellectual} é a propria intelligencia, emquanto \emph{conserva} as \emph{proprias} percepções \emph{intel lectuaes passadas}, e as \emph{reconhece} como \emph{proprias} e como \emph{passadas}. --- A intelligencia, sendo dotada de uma natureza incorruptivel, póde \emph{conservar} as percepções \emph{Intellectuaes} passadas, e, reflectindo \pg{176} \pg{178} \pg{180}  sobre si mesma e sobre os seus actos, póde \emph{reconhecer} aquellas percepções como \emph{proprias} e como \emph{passadas}.

Diremos alguma coisa ácerca do \emph{objecto} da \emph{memoria intellectual}, e provaremos que, se ella é \emph{especificamente} distineta da \emph{memoria sensitiva}, não o é da \emph{intelligencia} \footnote{A \emph{memoria intellectual} não forma novas percepções, mas reproduz as que anteriormente se formaram. --- Esta disposição para reproduzir as percepções passadas --- não é mera \emph{potencia}, pois que \emph{a potencia} póde encontrar muitas dificuldades na realisação do \emph{acto}, ao passo quo á \emph{memoria intellectual} não encontra a minima difficuldade na reproducção das percepções; --- nem é \emph{acto}, pois que \& consciencia attesta que o exercicio \emph{actual} da \emph{memoria} não é continuo; --- mas é um \emph{Habito}, porque a \emph{memoria}, com a maxima facilidade, reproduz as percepções passadas. --- Os escholasticos dizem que a \emph{meanoria intellectual} é a faculdade que conserva as \emph{especies intelligiveis}, anteriormente formadas. E' a mesma coisa. Emtanto se conservam as percepções \emph{intellectuaes passadas}, emquanto se conservam as \emph{especies intelligiveis}, que são o principio d'essas percepções. (Cf. \emph{Sum}. \emph{Theol}, p. 1, q. 19, a. 6).}.

a) \textbf{Objecto} da memoria intellectual, --- O objecto \emph{proprio} da \emph{memoria intellectual} é triplice: \emph{directo}, \emph{indirecto}, e \emph{implicito}. --- O \emph{directo} é constituido pelas operações \emph{intellectuaes passadas},

\emph{sz} como taes; --- o \emph{indirecto} é formado pelos \emph{objectos} d'aquellas percepções; --- o \emph{implicito} é o nexo existente entre as percepções \emph{intellectuaes} e o sujeito, que as reconhece como \emph{proprias} \footnote{O objecto de uma faculdade \emph{espiritual} não póde deixar de ser \emph{immaterial}. Por isso, tanto a relação de \emph{tempo}, que distingue as percepções \emph{iztellectuaes}, como o \emph{nexo}, que existe entre essas operações e o seu sujeito, são coisas conhecidas por um modo \emph{immaterial}, como é proprio da intelligencia.\par Objectam: A \emph{memoria}, no seu proprio conceito, é a percepção do \emph{passado}. Ora o \emph{passado}, por ser uma coisa particular e concreta, é objecto de uma faculdade \emph{sensitiva}, e não de uma faculdade \emph{intellectual}. Logo o objecto da \emph{memoria intellectual} não póde ser o \emph{passado}.\par Resposta. Concedendo a premissa \emph{maior}, devemos distinguir a \emph{menor}. O \emph{passado} pode referir-se --- tanto ao «\emph{cto}, por que percebemos um ente, --- como ao proprio \emph{ente}, que é percebido. --- Se se referir ao \emph{ente}, que é percebido, o \emph{passado} é objecto \emph{proprio} da memoria \emph{sensitiva}, porque um ente \emph{passado} é um ente \emph{singular} e \emph{concreto}, determinado pela cireumstancia do \emph{tempo}, e por isso só pode ser percebido directamente por uma faculdade \emph{sensitiva}, --- e não pode ser objecto proprio da \emph{intelligencia}, porque esta considera a essencia das coisas, abstrahida de todas as circumstancias de \emph{tempo} e de \emph{logar}. --- Se, porém, se referir ao \emph{acto}, o \emph{passado} é objecto \emph{proprio} não só de uma faculdade \emph{sensitiva}, mas tambem o é da \emph{intelligencia}, porque o. acto da intelligencia, embora \emph{singular} ou \emph{particular}, produzido nesta ou naquella differenca de tempo, todavia é sempre \emph{immaterial} ou \emph{espiritual}; e por isso a intelligencia, assim como se percebe a si mesma, embora seja uma coisa singular, assim tambem percebe o sen acto, embora seja este singular e concreto (Sum. \emph{Theol}, p. I, q. 79, a. 6 ad 2).\par Dissemos que os objectos \emph{passados}, como \emph{tes}, não formam o objecto \emph{proprio}\pgn{181} da \emph{memoria intellectual}. Todavia isso não obsta a que esta faculdade perceba de algum modo os objectos \emph{passados}, como \emph{taes}; pois a experiencia attesta que fallamos e raciocinamos d'esses objectos. Devemos, pois, dizer que, assim como a intelligencia percebe os entes \emph{singulares} e \emph{materiaes}, reflectindo sobre os phantasmas, assim tambem a memoria \emph{intellectual} veconhece os objectos \emph{passados}, como \emph{taes}, reflectindo sobre a memoria \emph{sensitiva}, de que esses objectos formam o objecto \emph{proprio}.\par Insistem : Os actos \emph{passados} da intelligencia, pelo facto de serem \emph{passados}, são particulares e concretos, e por isso não podem, \emph{de} modo algum, ser objecto da \emph{intelligencia}, que só attinge o universal e o abstracto.\par Repetimos que a intelligencia não percebe o particular ou concreto, quando este é \emph{material}; mas percebe-o, quando é \emph{immaterial}. Ora os actos da intelligencia, posto que, pelo facto de se considerarem como \emph{passados}, sejam particulares e concretos, são comtudo \emph{immateriaes}. Logo a intelligencia, assim como se entende a si mesma, embora seja singular, assim tambem entende as proprias percepções, embora singulares ou existentes numa certa differença de tempo. Diz S. Thomaz: « \emph{Singulare} non repugnat intelligi inquantum est singulare, \emph{sed inquantum est materiale}, quia nihil intelligitur nisi immaterialiter. Et ideo si sit aliquid \emph{singulare} et \emph{immateriale}, sicut. est intellectus, hoc non repugnat intelligi » (\emph{Sum}. \emph{Theol}., p. I, q. 86, a. 1 ad 3).}.

b) A memoria intellectual é especificamente distincte, da \textbf{memoria} sensitiva, --- \textsc{Duas} faculdades são \emph{especificamente} distinctas, ,quando os seus respectivos objectos são \emph{especificamente} distinctos. Ora o objecto proprio da memoria \emph{sensitiva}, por ser \emph{sensivel} e \emph{material}, \emph{é especificamente} distincto do objecto da memoria \emph{intellectual}, que é \emph{immaterial} e \emph{espiritual}. Logo a memoria \emph{intellectual} é \emph{especificamente} distineta da memoria \emph{sensitiva}. \footnote{O objecto da memoria \emph{sensitiva é sensivel}; porque, como dissemos, tanto as percepções \emph{passadas}, que a \emph{memoria sensitiva} conserva e reproduz, como tambem os \emph{objectos} d'essas percepções, são coisas \emph{sensiveis} e \emph{materiaes}. O objecto da memoria \emph{intellectual é immaterial}; porque as percepções \emph{ntellectuaes passadas}, que a memoria \emph{intellectual} conserva e reproduz, os \emph{objectos} d'essas percepções e o \emph{nexo} entre as percepções e o seu sujeito, são coisas \emph{espirituaes} e \emph{àmmalteriaes}. --- Alem d'isso, \& memoria intellectual não só conserva as \emph{proprias} percepções \emph{passadas}, o que faz tambem a \emph{sensitiva}, --- mas tambem, por ser dotada de reflexão e de poder abstractivo, \emph{reconhece} essas percepções como \emph{proprias} e como \emph{passadas}, o que não faz a \emph{sensitiva}, a qual, não reflectindo sobre si mesma, não pode conhecer o eu como sujeito de sensações, nem conhecer o abstracto (como é o \emph{passado}, \emph{emquanto tal}).}.

c) A memoria intellectual não é especificamente distincta da inteligencia, mas é uma funogão especial d'esta faculdade,

a) \emph{Não é especificamente distincta da intelligencia}. Com effeito, duas faculdades não são \emph{especificamente} distinctas, quando os seus objectos não são \emph{especificamente} distinctos. Ora o objecto da memoria \emph{intellectual} e o da \emph{éntelligencia} não são \emph{especificamente} distinctos; pois a memoria \emph{intellectual conserva} e \emph{reconhece} \pg{182}  as proprias percepções \emph{intellectuaes}, que, por serem \emph{zmmateriaes}, estão eomprehendidas no objecto da \emph{intelligencia} \footnote{Muitos escriptores modernos afirmam que a memoria \emph{intellectual} é uma faculdade \emph{especificamente} distincta da \emph{ntelligencia}.\par Esta affirmação é falsa; pois, como está provado no texto, o objecto da memoria \emph{intellectual} não é \emph{especificamente} distincto do objecto da \emph{intelligencia}. --- De mais, a \emph{intelligencia} póde exercer as funcções da \emph{memoria intellectual}, que são conservar e \emph{reconhecer} as percepções \emph{intellectuaes passadas}. Póde \emph{conserval-as}; pois que a intelligencia, sendo dotada de uma natureza mais du1adoira e mais forte que a dos sentidos, póde, com maior facilidade que estes, reter os proprios conhecimentos. Póde \emph{reconAecel-os}; pois que a intelligencia, tendo o poder de reflectir sobre si mesma e sobre as \emph{proprias} operações, \emph{existentes numa certa differenca de tempo}, póde perceber não sô que \emph{presentemente} percebe, mas tambem que \emph{anteriormente} percebeu (\emph{Sum}. \emph{Theol}., p. T, q. V9, a. 7).\par Objectam: A memoria \emph{intellectual} attinge o intelligivel emquanto é \emph{passado}, e a intelligencia attinge-o emquanto é \emph{presente}. Logo, sendo diversos os objectos \emph{formaes} das duas faculdades, tambem estas devem ser diversas.\par Resposta. Concedemos o antecedente, mas negamos o consequente. A diversidade entre dois objectos só exige e suppõe faculdades diversas, quando é \emph{essencial} ou \emph{substancial}, \emph{e} não quando é apenas \emph{accidental}. Assim o som e a cór suppõem sentidos diversos, porque a \emph{cór é essencialmente} diversa do \emph{som}; mas a côr \emph{verde} e a côr \emph{vermelha} não se referem a duas faculdades diversas, porque o \emph{verde} e o \emph{vermelho} são \emph{accidentaes} á côr. Ora, se o \emph{intelligivel} e o \emph{sensivel} exigem duas faculdades diversas, porque o \emph{intelligivel é essencialmente} diverso do \emph{sensivel}, não exigem duas faculdades diversas o intelligivel \emph{presente} (objecto da \emph{intelligencia}) e o intelligivel \emph{passado} (objecto da \emph{memoria intellectual}), porque o \emph{presente} e o \emph{passado} são coisas \emph{accidentaes} para o \emph{intelligivel}, que consiste inteiramente na reprodue(áo ou expressão da \emph{essencia} das coisas, abstrahida de todas as condições de \emph{tempo} e de \emph{logar}. (S. Thom., \emph{De} Verit., q. 10, a. 8).\par Agora vemos a razão por que a \emph{memoria intellectual} não é uma faculdade distincta da \emph{intelligencia}, ao passo que a \emph{memoria sensitiva} é distincta, de todas as outras faculdades sensitivas. A faculdade \emph{perceptivo-éntellectuat}, por ser inorganica ou espiritual, não só póde \emph{perceber} as coisas, mas tambem conservar e \emph{reproduzir} essas \emph{percepções}; e por isso a faculdade \emph{intellectual} que \emph{percebe} não se distingue da faculdade que conserva e \emph{regrodus} as percepções. Pelo contrario, uma faculdade \emph{perceptivo-sensitiva}, por ser organica ou material, se \emph{percebe} as coisas, não póde \emph{conservar} nem \emph{reproduzir} as percepções; porque uma faculdade \emph{organica}, tendo uma só disposição, não póde exercer mais que uma funcção; e por isso a faculdade \emph{sensitiva} que \emph{percebe} deve distinguir-se da faculdade que \emph{conserva} e \emph{reproduz} as percepções. (S. Thom., \emph{De Verit}., q. 10, a. 2).}.

6) \emph{Mas é uma funceüo especial da intelligencia}. Porquanto, a funeção, por que a intelligencia \emph{conserva} e \emph{reconhece} as percepções \emph{passadas}, é distineta d'aquella, por que a mesma intelligencia percebe \emph{actua}? e \emph{presentemente} a essencia das coisas \footnote{Os principaes auxilios da \emph{memoria intellectual} são dois: a \emph{attenção} e a \emph{associação das ideas}. --- A \emph{attenção} é a intensa applicação da intelligencia\pgn{183}a um objecto. --- A \emph{associação das ideas} é o nexo que liga as ideas, de modo que a reproduecio de uma traz a reproduccüo de outra. Esta \emph{associação é natural} ou \emph{artificial}, segundo a relação entre as duas ideas depende da \emph{zatureza} (como a que existe entre a \emph{causa} e o \emph{effeito}), ou da arte(como a que existe entre a \emph{oliveira} e a \emph{paz}) --- A \emph{associação} das ideas obedece a certas leis, que S: Thomaz, explicando Aristoteles (\emph{De mem}. \emph{et rem}.), reduz a tres, e que são a lei de \emph{semelhança}, a lei de \emph{contrariedade}, á lei de \emph{conjunccüáo}, emquanto a lembrança de uma coisa traz naturalmente comsigo a lembranca de outra coisa, ou \emph{semelhante}, ou \emph{contraria}, ou \emph{unida}, de algum modo, á primeira. Nestas leis \emph{especiaes} está incluida uma lei \emph{geral}, chamada lei de \emph{attenção}, segundo a qual nos lembramos mais facilmente das coisas, que foram percebidas com maior esforço de \emph{attenção} (Cf. \emph{Sum}. \emph{Theol}., 2-2, q. 49, a. 1).}.

%%>>> art-18

\section*{Artigo II.\\ Vontade, sua natureza e liberdade.}
\addcontentsline{toc}{section}{Artigo II. --- Vontade, sua natureza e liberdade.}

%%>>> t125

\pg{181} \pg{183} \tese{125}{Vontade.}
\emph{Vontade}, ou \emph{appetite intellectual}, é a faculdade pela qual o homem tende para o bem \emph{immaterial}, conveniente á natureza \emph{racional} e percebido pela \emph{intelligencia}.

Esta faculdade existe no homem. Com effeito, as tendencias de um ente são sempre proporcionadas aos seus conhecimentos. Ora o homem, pela \emph{intelligencia}, conhece o bem que aperfeiçõa a sua natureza racional. Logo deve ter e tem uma tendencia para esse bem. Esta tendencia deriva da \emph{vontade} \footnote{A \emph{vontade}, pois, distingue-se do \emph{appetite natural}, do \emph{appetite sensitivo} e da \emph{intelligencia}. --- Distingue-se do \emph{appetite natural}; porque este nasce da propria \emph{natureza} do ente, sem percepção precedente, ao passo que a vontade suppõe uma previa percepção do bem. --- Distingue-se do \emph{appetite sensitivo}; porque este tende para um bem \emph{sensivel}, e a vontade para um bem \emph{immaterial}. --- Distingue-se da 2\emph{ntelligencia}, porque o objecto d'esta faculdade, embora seja \emph{materialmente} identico ao da \emph{vontade}, todavia \emph{formalmente} é diverso d'elle; visto que a intelligencia apprehende o ente emquanto é \emph{verdadeiro}, e a vontade tende para o ente, emquanto é \emph{bom} (Cf. S. Thom,, \emph{De} Ver., q. 22, a. 10).\par Pela mesma inclinação, pela qual tende para o bem, que a intelligencia conhece ser conveniente á natureza humana, a vontade foge do mal que a intelligencia percebe ser prejudicial á mesma natureza.\par Diz S. Thomaz: « Todas as coisas são inclinadas pelo appetite para o bem, mas nem todas do mesmo modo. --- Algumas são inclinadas para o bem unicamente pela natural tendencia, e sem conhecimento algum, como são as plantas e os corpos inanimados. Esta inclinação chama-se \emph{appetite}\pgn{184}\emph{natural}. --- Outras são inclinadas para o bem com algum conhecimento, não emquanto percebem a propria razão ou essencia de bem, mas emquanto conhecem um certo bem particular; assim os sentidos percebem o doce, o branco, etc. Esta inclinação chama-se \emph{appetite sensitivo}. --- Outras são inclinadas para o bem, e conhecem a propria essencia de bem. Esta inclinação é propria dos seres intelligentes, e é a mais perfeita de todas. Por quanto, \textbf{estes} seres, além de terem uma tendencia natural para o bem, como os que são destituidos de conhecimento, têm tambem uma inclinação, proveniente do conhecimento, e esta inclinação não só a têm para o bem \emph{particular}, mas tambem para o bem \emph{universal}. Esta inclinação chama-se \emph{vontade} » (\emph{Sum}. \emph{Theol}., p. I, q. 59, a. 1).}.

%%>>> t126

\tese{126}{Objecto da vontade.}
Objecto da \emph{vontade} é o \emph{bem immaterial}, \emph{conveniente á natureza racional}, e \emph{percebido pela intel}-

4 \emph{ligencia}. --- W' o \emph{bem immaterial}; pois a vontade, sendo, como veremos, uma faculdade \emph{espiritual}, só póde tender para um \pg{184}  objecto \emph{immaterial}: --- \emph{conveniente á natureza racional}; pois que, \textbf{para} a vontade tender para um objecto, não basta que seja um \emph{lem}, mas é mister que seja um bem \emph{conveniente}: --- \emph{percebido pela intelligencia}, visto que a vontade não póde desejar uma coisa, que não tenha sido anteriormente conhecida pela intelligencia \footnote{\emph{A} vontade tende para o bem, que lhe apresenta a intelligencia. Ora a inteligencia póde apontar um bem \emph{verdadeiro}, ou um bem \emph{apparente}, que na realidade é um mal. Por isso a vontade póde tender --- tanto para um bem \emph{verdadeiro}, --- como para um bem «\emph{apperente}. Mas, nesta, segunda hypothese, a vontade não tende \emph{directamente} para o \emph{mal} emquanto \emph{tal} (como erradamente affirmou Saisset, \emph{Dict}. \emph{phil}., art. \emph{Liberté}) mas tende para o mal, emquanto este é conhecido como bem.}.

%%>>> t127

\tese{127}{Divisão do objecto da vontade.}
O objecto da \emph{vontade} divide-se em \emph{adequado} e \emph{inadequado}.

a) O objecto \emph{adequado} (que eguala a capacidade da vontade, satisfazendo inteiramente as suas tendencias) é o bem \emph{universal}, que é só bem e que por isso forma a nossa felicidade; pois que só o bem \emph{universal} póde saciar completamente os desejos da alma. --- O bem \emph{universal} é o \emph{fim ultimo} da vontade \footnote{O bem \emph{universal} e \emph{absoluto}, que é objecto adequado da vontade, não é um bem \emph{concreto} e \emph{determinado}, mas é um bem \emph{abstracto} e \emph{indeterminado}, formado pela intelligencia, a qual exclue dos bens particulares todo o limite e toda a imperfeição. Todavia «sse bem \emph{universal}, \emph{abstracto} e \emph{indeterminado}, que a intelligencia forma e aponta á vontade como o objecto da felicidade, suppõe necessariamente a existencia de um bem \emph{universal}, mas \emph{concreto} e \emph{determinado}; porque os bens finitos suppõem o bem infinito, e porque não póde ser irrasoavel nem frustrada a tendencia irresistivel de todos os homens para a \emph{felicidade}. Este bem, possuido, sacia completamente a vontade.}.

b) O objecto \emph{inadequado} (que não eguala a capacidade da vontade, nem satisfaz perfeitamente as suas tendencias) é todo o bem \emph{particular}, que é uma participação \emph{limitada} do bem \emph{universal} e \emph{absoluto}; porque todos os bens \emph{particulares}, embora reunidos, por serem limitados, não podem satisfazer as infinitas \pg{185}  aspirações da nossa vontade. --- Os bens \emph{particulares} são \emph{eios} para a consecução do \emph{fim ultimo} \footnote{Os bens \emph{particulares} e \emph{relativos}, que formam o objecto \emph{inadequado} da vontade, podem serimmateriaes ou materiaes, pois a vontade póde quere-los, emquanto convém --- ou ás faculdades \emph{espirituaes} do homem, --- ou ás suas --- faculdades \emph{orgamicas}.--- A vontade muitas vezes tende para os bens \emph{particulares}, que, se convém ás faculdades \emph{organicas}, são prejudiciaes ás faculdades \emph{espirituaes}. D'ahi a desordem moral; porque o bem da razão, sendo o verdadeiro \emph{bem} do homem, deve prevalecer aos bens das faculdades materiaes.}.

A vontade tende para os bens \emph{particulares}, pelo facto de estes serem uma participação, mais ou menos limitada, do bem \emph{universal}, e, tendendo para aquelles, tende para este, pois os meios empregam-se para a consecução do fun. --- D'este modo o bem \emph{universal} é o principio e o fim, por que a nossa vontade tende para os bens \emph{particulares} \footnote{Com razão, pois, os philosophos dizem que o bem \emph{universal} e \emph{absoluto} € o objecto \emph{principal} da vontade, e que o bem \emph{particular} e \emph{relativo} constitue o seu objecto \emph{secundario}.}.

%%>>> t128

\tese{128}{A vontade é uma faculdade espiritual.}

2) Toda a faculdade, eujo objecto é universal e immaterial, é \emph{espiritual}; não só porque todo o ser tende para o bem, que lhe é semelhante, mas tambem porque o universal e o immaterial não podem impressionar uma faculdade material. Ora o objecto da \emph{vontade} é o bem \emph{universal} e \emph{immaterial}; pois que a vontade tende para o bem, que é só bem e exclue todo o mal, e para os bens \emph{Zmmateriaes}, que são a \emph{virtude}, \emph{á sciencia}, ete. Logo a vontade é uma faculdade \emph{espiritual}.

b) Se a \emph{vontade} fosse uma faculdade organica ou dependente da materia, não poderia levantar-se sobre a mesma materia e dominal-a; porque repugna que o instrumento domine o agente pelo qual é movido. Ora essa faculdade effectivamente rege e domina a materia, sujeitando-a ás dores, aos tormentos e até á morte, apezar de toda a resistencia do appetite sensitivo. Logo a \emph{vontade} não depende da materia, mas é uma faculdade \emph{espiritual}.

c) Uma faculdade, que reflecte sobre si mesma e sobre os seus actos, é \emph{espiritual}; pois uma faculdade material é composta de partes extensas, e uma parte extensa póde reflectir sobre uma outra, mas não sobre si mesma. Ora a \emph{vontade} reflecte sobre si mesma e sobre os seus actos; pois quer os seus desejos e ama as suas tendencias, de modo que é faculdade volitiva e objecto querido. Logo é faculdade \emph{espiritual} \footnote{Os \emph{actos} da vontade, como tambem os da intelligencia e das outras\pgn{186}faculdades, podem ser objecto d'esta faculdade volitiva; pois que podem ser percebidos como \emph{bens} e bens \emph{convenientes} á nossa natureza.\par Como se vê, o methodo para a demonstração da espiritualidade da \emph{vontade} é o mesmo que seguimos para a demonstração da espiritualidade da \emph{intelligencia}. E' o que se adopta em todas as sciencias, que do conhecimento dos phenomenos se elevam ao conhecimento das causas respectivas. Diz S. Thomaz: « Rei cujuslibet perfecta cognitio haberi non potest, nisi ejus operatio cognoscatur. Ex modo enim operationis et specie, mensura et qualitas virtutis pensatur. Virtus vero naturam rei monstrat; secundum hoc enim unumquodque natum est operari, qnod actu talem naturam sortitur n. (C. \emph{Gentes}, 1. II, c. 1).}.

%%>>> t129

\pg{186} \tese{129}{A vontade e a liberdade.}
\emph{Liberdade} é a propriedade da \emph{vontade}, pela qual esta é senhora dos seus actos; de modo que, proposto um objecto, póde quere-lo e deixar de o querer, ou querer um bem diverso e até opposto. --- A \emph{vontade}, emquanto dotada de \emph{liberdade}, chamar-se \emph{livre arbitrio} \footnote{\emph{Liberdade}, em geral, é a immunidade ou isenção de dependencia; pois que dizemos que um ser é livre, quando não depende de outro.\par E, como esta dependencia póde consistir numa \emph{coacção} ou \emph{violencia externa}, procedente de um agente externo, --- ou numa \emph{necessidade} ou \emph{determinação interna}, procedente da natureza, --- ou na \emph{lei}, procedente do superior, que tem a seu cuidado o bem social, d'ahi resulta uma triplice especie de liberdade: de \emph{coacção}, de \emph{necessidade} e de \emph{lei}.\par \emph{a}) Liberdade de \emph{concção} é a isenção de toda a violencia feita a um ser por um agente extrinseco e contra a inclinação do proprio ser. --- A vontade não póde soffrer \emph{coacção} nos seus actos \emph{elicitos,isto} é, nos actos que a vontade exerce \emph{immediatamente} por si; não só porque estes actos escapam á acção directa dos agentes externos, mas tambem porque, se elles derivam \emph{immediatamente} da vontade, não podem ser contrarios á inclinação d'ella. Mas, se não póde soffrer \emph{coacção} nos seus actos \emph{elícitos}, póde soffre-la nos seus actos \emph{émperados}, que a vontade exerce por meio de outras faculdades; pois que o agente exterior póde obstar a que uma faculdade \emph{organica} e \emph{exterior} execute o imperio da vontade.\par \emph{b}) Liberdade de \emph{necessidade} é a isenção de toda a \emph{determinação intrinseca}, de modo que o ente póde fazer ou deixar de fazer uma coisa e até fazer o contrario. O ente, pois, que possue esta liberdade, não é movido ou dirigido por outro, mas é elle que se dirige e se move a si mesmo. --- Esta liberdade chama-se tambem --- liberdade de \emph{eleição}, porque \emph{escolhe} uma coisa, deixando outra, --- e liberdade de \emph{indifferença}, porque por si póde determinar-se tanto para uma coisa como para outra, diversa ou contraria.\par c) Liberdade de \emph{lei} é a isenção de toda a \emph{determinação}, proveniente da vontade do legislador. Só \textsc{Deus} não está sujeito á lei, porque a sua vontade é a propria rectidão e santidade, e porque não tem superior. Os homens estão todos sujeitos á lei do seu soberano Creador e Senhor.\par A liberdade, que é um dote da nossa vontade, é a de \emph{necessidade}. D'esta especie devemos aqui occupar-nos com algum desenvolvimento.\par A liberdade de \emph{necessidade} chama-se tambem \emph{indifferença}. --- Como \emph{indifferente} ou \emph{contingente} se diz o ente que é \emph{intrinsecamente} indeterminado\pgn{187}, por isso á ?\emph{zdifferenca} é duplice: \emph{passiva} e \emph{activa}. A indifferenca \emph{passiva} denota a capacidade de receber de outro ente uma determinação; a \emph{activa} é a força, que o ente possue, de se determinar a si mesmo, ou de sahir do estado de indifferenga e de produzir o acto. --- \emph{A passiva} denota imperfeição, e encontra-se no appetite \emph{sensitivo}; a \emph{activa} diz \emph{perfeição}, e é propria do appetite \emph{intellectual}, que é a vontade.\par \emph{A} liberdade de \emph{necessidade}, ou a indifferenca \emph{activa}, divide-se em varias especies, conforme a vontade é \emph{intrinsecamente} indeterminada com relação ao \emph{acto}, ou com relação ao \emph{objecto}.\par Se a vontade é indeterminada com relação ao \emph{acto}, de modo que póde querer e deixar de querer, a sua liberdade chama-se de \emph{contradicção}. --- Se a vontade é indeterminada com relação ao \emph{objecto}, de modo que póde escolher entre objectos \emph{especificamente diversos}, ou entre objectos \emph{oppostos}, a sua liberdade é de \emph{especificação} ou de \emph{contrariedade}.\par Liberdade, pois, de \emph{contradicção} é o poder que a vontade tem de produzir ou deixar de produzir o seu acto, p. ex., de ler ou deixar de ler, etc. Chama-se liberdade de \emph{contradicção} e tambem de \emph{exercicio}: --- de \emph{contradicção}, porque entre o fazer e o deixar de fazer não ha meio termo, mas uma \emph{contradicção}, --- de \emph{exercicio}, porque consiste no poder de exercitar ou deixar de cxercitar o acto. --- Liberdade de \emph{especificação} é o poder que a vontade tem de escolher entre objectos \emph{especificamente} diversos, p. ex., entre o estudo, o passeio, ete. --- Liberdade de \emph{contrariedade} é o poder que a vontade tem de escolher entre objectos \emph{oppostos}, ex. \emph{gr}., entre o bem e o mal, etc.\par Ora, para que a vontade possa chamar-se \emph{livre}, basta que possua a liberdade de \emph{contradicção} ou d'exercicio. Porquanto, um ente é verdadeiramente livre, se tiver o poder de \emph{eleição}. Diz S. Thomaz: « Proprium liberi arbitrii \emph{est.electio} » (\emph{Sum}. \emph{Theol}., p. I, q. 88, a. 3). Ora este poder existe, desde que a vontade tenha a faculdade de operar ou deixar de operar. --- Além d'isso, é necessario advertir que a liberdade de \emph{especificação} e a de \emph{contrariedade} são apenas applicações diversas da liberdade de \emph{contradicção}.\par D'ahi se vé que o poder de escolher entre o bem e o mal, e por isso o poder de peccar, não constitue a liberdade, mas, pelo contrario, é um defeito della; e, um signal de liberdade é-o como a doença é signal de yida.\par SUN : ao signal do aiii (Cf. S. Th. \emph{De} Ver., q. 22, 8. q" \emph{ENE on} a so \emph{OD} uf dy t}.

Todavia a \emph{liberdade} é um dote da vontade, --- não quando esta tende para o bem \emph{universal}, que forma a \emph{felicidade}, --- mas só quando tende para os bens \emph{particulares}.

%%>>> t130

\pg{187} \tese{130}{A vontade não é livre na sua tendencia para o bem universal.}

2a) Um ser não é indifferente ou livre na sua tendencia para o fim, que o Auctor da natureza lhe estabeleceu; pois que, nessa tendencia, o ser não se determina a si mesmo, mas é determinado. Ora o bem \emph{universal}, que forma a felicidade, é o fim, que o Auctor da natureza estabeleceu á nossa vontade;

porque esta por si e espontaneamente deseja ser feliz. Logo a vontade não é livre na sua tendencia para o bem \emph{universal}.

b) A consciencia e a experiencia attestam que todos os homens, sempre e em toda a parte, tendem para o bem \emph{uni}- \pg{188}  \emph{E} Tov aa VA * TR T nda SR o ao, z --------- CM

188 METAPHYSICA --- PARTE TERCEIRA --- ANTHEOPO:OGIA \emph{versal}, e que essa tendencia é necessaria, e não é possivel contrarial-a, muito menos para!ysal-a ou destruil-a. Ora este \textbf{testemunho} seria inexplieavel, se á vontade fosse livre na sua tendencia para o bem \emph{universal}. Logo a vontade não é livre na sua tendencia para o bem \emph{universal}.

€) A vontade está para o bem \emph{universal} e \emph{absoluto}, que é o objecto da felicidade, como a intelligencia está para os primeiros principios de toda a sciencia; porque a felicidade é na pratica o que os primeiros principios são na theorica. Ora a intelligencia não é livre em adherir aos primeiros principios; porque a verdade, que elles contém, por ser evidente, impõe-se necessariamente á mesma faculdade. Logo a vontade não é livre na sua tendencia para o bem \emph{universal}; pois que este bem, sendo só bem e excluindo todo o mal, satisfaz completamente a capacidade da vontade, e não apresenta motivo algum, por. que deva ser rejeitado. Logo a vontade não é livre na sua tendencia para o bem \emph{universal} \footnote{São duas as faculdades, de que o homem é dotado emquanto \emph{racional}, e que o tornam immensamente superior a todos os outros seres visiveis: a \emph{intelligencia} e a \emph{vontade}. A intelligencia está orientada para a \emph{verdade}; a vontade, para o \emph{bem}. Esta orientação ou tendencia \emph{geral} é natural e necessaria, e por isso deriva do proprio Auctor da natureza. O acto da vontade não só deve ser precedido pelo da intellizencia, mas a este deve ser proporcionado. Ora, se a intelligencia fosse uma faculdade organica, e por isso só conhecesse bens sensiveis e particulares, só para estes poderia tender a vontade; e o homem não poderia nem conhecer nem amar outros bens mais elevados. Mas a intelligencia, por ser uma faculdade espiritual e desligada de todos os laços do tempo e do espaço, levanta-se, da vista dos bens particulares, á contemplação de um bem \emph{absoluto}, \emph{universal}, que concentra em si todas as perfeições, \emph{e} mostra-o á vontade como o fim ultimo d'ella, como o unico bem capaz de preencher toda a sua capacidade.\par Poderia á vontade em presença do bem \emph{universal}, da sua felicidade, despresar a posse d'ella ou pelo menos ficar indifferente? De forma alguma. Quando o objecto corresponde plenamente á capacidade do sujeito e como tal é conhecido, isto é, quando existe e se conhece uma equação perfeita entre a capacidade do sujeito e á bondade do objecto, então nenhuma razão ha para rejeitar o dito objecto, e nesse caso a liberdade seria uma desordem, supporia que o sujeito não é feito para o objecto, nem tende para elle. Este facto tem a sua razão naquella finalidade intrinseca, que governa a actividade de todos os seres: todo o ser tem uma inclinação natural, pela qual as suas faculdades e as suas operações tendem e convergem para um fim determinado. A esta lei natural, que é a base da ordem universal, não póde o homem subtrahir-se. Eis porque nos sentimos levados irresistivelmente a desejar o bem universal, a felicidade, e não podemos deixar de querer um objecto, no qual se considere só o caracter de bondade.\par A vontade, tendendo irresistivelmente para o bem \emph{universal} e \emph{perfeito},\pgn{189}não póde deixar de tender para aquelles bens \emph{particulares}, que \emph{claramente} conhecemos ter um nexo \emph{necessario} com o mesmo bem \emph{universal}, ou como meios ou como condições; pois quem quizer alcançar o fim, não póde deixar de querer aquellas coisas, que \emph{claramente} se conhecem indispensaveis para a consecução do mesmo fim. Assim não podemos deixar de amar a vida, quando esta se conhece ser o fundamento da felicidade.\par D'onde se segue que, se encontrassemos num certo ente a realisação do ideal que nos fascina \emph{érresistivelnente}, isto é, se a nossa intelligencia conhecesse \emph{claramente} que um certo ente é o bem \emph{universal}, que é só bem e exclue toda a imperfeição, a vontade teria uma tendencia necessaria para a consecução d'esse ente, e, possuindo-o, não poderia deixar de o amar. Mas, se esse conhecimento não é \emph{claro}, essa tendencia tambem não é necessaria (Cf. \emph{Sum}. \emph{Th}., 1-2, q. 10, a. 2 ad 3).\par Mas escutemos o Angelico Mestre. Diz S. Thomaz que a nossa vontade, como toda a natureza creada, está orientada por Deus para o seu fim conveniente: « Hoc est cujuslibet naturae creatae proprium, ut a Deo sit ordinata in bonum, naturaliter appetens illud. Unde et voluntati ipsi inest naturalis quidam appetitus sibi convenientis boni» (\emph{De} Ver., q. 22, a. 5). --- E esta tendencia da vontade para o bem universal é \emph{necessaria}: « Necesse est quod sicut intelleetus ex necessitate inhaeret primis principiis, ita voluntas ex necessitate inhaereat ultimo fini, qui est \emph{beatitudo} » (\emph{Sum}. \emph{Theol}., p. 1, q. 82, a. 1). --- Note-se que esta \emph{necessidade}, pela qual o homem deseja a felicidade, diz respeito --- á \emph{determinação} do acto, emquanto o homem, pensando na felicidade, não pode querer o contrario (por isso o acto é \emph{determinado}), --- e não ao \emph{exercicio} do proprio acto, porque o homem pode não pensar na felicidade, e por isso pode não deseja-la, visto que o acto d'esse pensamento e d'esse desejo é um bem \emph{particular} e por isso não se impõe necessariamente (Cf. QQ. DD., \emph{De malo}, q. 6).\par Agora podemos perceber a razão por que Deus, embora seja o Bem absoluto e universal, comtudo, \emph{sobre} « \emph{terra}, não attrahe \emph{necessariamente} a nossa vontade, nem o sen amor sacia \emph{completamente} a nossa alma. --- O bem absoluto e universal, que nesta vida conhecemos e para o qual tende necessariamente \& nossa vontade, é abstracto e indeterminado, é uma concepção formada pela consideração dos bens particulares e finitos; ao passo que Deus é o Bem absoluto e universal, mas concreto, real e determinado. E, embora se \emph{demonstre} que o bem universal e abstracto encontra a sua realisação e se concretisa em Deus, comtudo, porque a intelligencia não vê essa realisação por \emph{intuição inmediata}, \& vontade não é attrahida \emph{irresistivelmente} por Deus, mas fica sempre \emph{livre}. Sim, sempre \emph{livre}; porque, tendendo sempre e necessariamente para o bem universal, para a felicidade, pode procurar e procura os bans caducos da terra. Nem o amor de Deus sacia \emph{completamente} a nossa alma; porque esse amor, além de ser amissivel, anda sempre unido com muitos males e afílicções da alma e do \textbf{corpo}. ---\pgn{190}Dissemos --- \emph{sobre a terra}; porque no Ceu a nossa intelligencia, confortada pela luz da gloria e elevada á visão \emph{intuitiva} da Essencia divina, verá \emph{claramente} que esse bem universal e abstracto, para o qual hoje tendemos necessariamente, se concretisa em Deus; e então a vontade em Deus concentrará todo esse amor necessario, que presentemente consagra ao bem universal e abstracto, e neste amor encontrará a sua felicidade.\par Mas escutemos um trecho delicioso de S. Thomaz. Diz assim: « Sicut -intellectus naturaliter et ex necessitate inhaeret primis principiis, ita voluntas ultimo fini... Quaedam autem propositiones sunt \emph{necessariae}, quae habent connexionem \emph{necessariam}. cum primis principiis, sicut conclusiones demonstrabiles, ad quarum remotionem sequitur remotio primorum principiorum; et his intellectus \emph{ex necessitate} assentit, cognita connexione \emph{necessaria} conclusionum ad principia per demonstrationis deductionem. Non autem ex necessitate assentit, antequam hujusmodi necessitatem connexionis per demonstrationem cognoscat. \emph{Similiter} etiam ex parte \emph{voluntatis} . . . Sunt enim quaedam particularia bona, habentia \emph{zecessariam} connexionem ad beatitudinem, quibus scilicet homo Deo inhaeret, in quo solo vera beatitudo consistit. Sed tamen, antequam per certitudinem divinae visionis \emph{necessitas} hujusmodi connexionis demonstretur, voluntas \emph{non em necessitate} Dão inhaeret, neo his quae Dei sunt. Sed \emph{voluntas videntis Deum per essentium}, de \emph{necessitate inhaeret} Deo, sicut nunc ex necessitate volumus esse beati » (Sum. TÀ., p. T, q. 82, a. 2).}.

%%>>> t131

\pg{189} \tese{131}{A vontade é livre na sua tendencia para os bens particulares.}

a) Um ser é indifferente ou livre na sua tendencia para o termo, que lhe não foi determinado pelo Auctor da natureza: pois que todo o agente que, nas suas tendencias, não é \textbf{determinado} pela natureza, determina-se a si mesmo. Ora os bens \pg{190}  \emph{particulares} não constituem o \emph{fim}, que o Auctor da natureza estabeleceu á nossa vontade; pois que esse fim, como vimos, é o bem \emph{universal}. Logo a vontade é livre na sua tendencia para os bens \emph{particulares} \footnote{Os bens \emph{particulares}, de que fallamos neste numero, são os que não estão \emph{necessariamente} ligados com o bem \emph{universal}. --- Todos os bens, que existem, são \emph{particulares} ou \emph{em si} ou \emph{para nós}, e não estão ou não se conhecem \emph{necessariamente} ligados com o bem \emph{universal}. ;}.

b) O consenso unanime dos povos confirma a these. Todos os povos, ainda os mais barbaros, sentem a força do \emph{dever}, têm leis que os governam, deeretam recompensas para a virtude e castigos para o vicio, empregam exhortações e conselhos. Ora, se a nossa vontade, na sua tendencia para os bens \emph{particulares}, não fosse inteiramente livre, e por isso não podesse deixar de fazer o que faz, o sentimento do \emph{dever} seria inexplicavel, como inexplicaveis seriam as leis; as recompensas e os castigos seriam uma injustiça ou uma tyrannia, os conselhos um absurdo. Logo a vontade é livre na sua tendencia para os bens \emph{particulares} \footnote{A conclusão é evidente. Diz S. Thomaz: « Homo est liberi arbitrii" alioquin frustra essent consilia, exhortationes, praecepta, prohibitiones, praemia et poenae » (\emph{Sum}. \emph{Theol}, p. Ll, q. 89, a. 1). --- Mas insistamos na idea de \emph{dever}, que é como que uma luz a esclarecer o interior da nossa alma e uma voz a fallar ao nosso coração. O \emph{dever} impõe-se á razão e á\par \emph{D}\pgn{191}consciencia em todos os tempos e em todos os logares, aos povos civilisados e aos selvagens, como norma suprema e absoluta, á qual devemos subordinar os nossos actos deliberados, sob pena de nos sentirmos reos de culpa e dignos de castigo. Tudo se tem contestado no campo da \emph{theoria}; mas, na \emph{pratica}, ninguem sequer suspeitou que a grande lei do \emph{deter}, regulador da vida humana, fosse uma illusão e não uma realidade indiscutivel. O proprio \emph{Kant}, que negou o valor \emph{objectivo} de todas as nossas percepções, não pos em x a existencia do \emph{dever}, e sobre o \emph{dever}, que elle chama \emph{pedra angular} da Religião e da moral, tentou reconstruir o edificio. da Metaphysica. Existe, pois, no fundo das nossas almas á grande lei do \emph{dever}.\par Mas se o homem está sujeito á grande lei do \emph{dever}, se se sente obrigado a praticar o bem e a evitar o mal, é porque \emph{pode} conformar-se com o \emph{dever}, e pode tambem afastar-se d'elle, é porque possue o dominio dos seus actos, é porque é \emph{livre}. Não se dà uma ordem, ou uma prohibição, se não. aos que podem obadecer. CN desobedecer. Portanto a existencia do \emph{dever} exige e \textbf{suppõe} a existencia do \emph{livre arbitrio}. ---\par Os deterministas modernos admittem a existencia do \emph{dever}, e comtudo negam ao homem o \emph{livre arbitrio}. Mas de que modo podem elles conciliar estas duas coisas: a existencia do \emph{dever} e a \emph{não-existencia} do \emph{livre arbitrio}? De um modo muito simples: desvirtuando ou falsificando a idea do dever. --- \emph{Guyan} diz que o \emph{dever} é o \emph{attractivo do bem}. Mas isto é falso. O \emph{altraclivo}, se pode ser uma \emph{consequencia secundaria} do \emph{dever}, não pode ser o proprio \emph{dever}. Por quanto, se fosse um \emph{attractivo}, o \emph{dever} apresentaria sempre um aspecto agradavel e sorridente, e nunca tomaria, como muitas vezes toma, as austeras apparencias da immolação e do sacrificio. Além d'isso, o \emph{altractivo} é mutavel, pessoal, caprichoso; ao passo que o \emph{dever} é immutavel e independente das conveniencias pessoaes, é um imperativo categorico incondicional. --- Um outro determinista, \emph{Fouillée}, escreve que o \emph{dever} é o ideal, descoberto pela razão, e por isso pode, como todo o ideal, existir independentemente da liberdade. Esta opinião não é menos absurda do que a primeira. O dever impõe-se, mas o ideal não é imposto, nem pela consciencia, nem pela leis. Quem poderá convencer-se de que está sujeito no dever de uma figura formosa? Poderiam acerescentar que só querem fallar de um \emph{ideal moral}. Mas, nem assim correcta, a explicação pode ser admittida. O \emph{ideal moral} só se torna \emph{dever}, \emph{quando é imposto}. Mas, embora 8e admittisse a identidade do \emph{dever} com o \emph{ideal moral}, seguir-se-hia a negação\pgn{192} de \emph{livre arbitrio}? Não: antes o contrario. Nenhum ideal pode ser imposto ao homem, se este não é livre de o copiar e de o reproduzir em seus actos. --- Outros deterministas dizem que o \emph{dever} é uma obrigação, que o homem impõe a si mesmo. Se isso fosse verdade, nenhum criminoso seria lacerado, como realmente o é, pelos remorsos do mal commettido!\par Concluimos. O \emph{dever} existe. Mas o dever não pode existir sem o \emph{livre arbitrio}. Logo o livre arbitrio existo. --- E' esta uma prova \emph{indirecta} da existencia do livre arbitrio; mas será sempre a mais popular, porque se refere a toda a nossa vida moral. :\par Objectam: O argumento, tirado da sancção dos castigos destinados aos criminosos, não tem \emph{força}; pois que applicam-se castigos aos animaes, que todavia não são dotados de \emph{livre arbitrio}.\par Respondemos que os casos são differentes. --- Os castigos, que se applicam aos animaes, não são verdadeiros \emph{castigos}, pois que os animaes não têm merito nem demerito; e applicam-se-lhes, para que estes, determinados pela imagem da dor e da morte, não tornem a commelter certos actos. --- Pelo contrario, os castigos, que se applicam ao homem, são verdadeiros \emph{castigos}, pois que ás nossas acções têm a razão de merito e de demerito; e applicam-se-nos, para que nos abstenhamos de certas acções, não determinados pelo instincto, mas levados pela livre eleição da vontade.}.

0) A consciencia, \emph{antes}, no \emph{momento} e \emph{depois} da operação, attesta que a nossa vontade, tendendo para um bem \emph{particular}, é livre. --- Antes dà operação; porque sentimos que podemos fazer ou deixar de fazer uma coisa, e isto sem violencia ou eonstrangimento, mas espontaneamente. --- No \emph{momento} da operação; pois conhecemos que podemos deixar e muitas vezes deixamos uma obra incompleta. --- \emph{Depois} da operação; quando a» mesma consciencia, por meio de uma intima consolação, approva uma obra bóa, ou, por meio de remorsos, condemna uma obra má, --- approvacao e condemnação que não teriam motivo sufficiente, se a vontade não tivesse sido livre no seu \pg{192}  acto. Logo a vontade é livre na sua tendencia para os bens

\emph{particulares} \footnote{Se o \emph{livre arbitrio} foi concedido ao homem como meio indispensavel para a vida moral e por isso para a consecução do fim ultimo, a sua existencia é um facto psychologico, e, como tal, deve ser percebido pela experiencia \emph{interna}, pela \emph{consciencia}. Não seria extranho que o homem não tivesse consciencia da sua \emph{liberdade}? --- Dizem os neo-eriticos que a experiencia, se pode conhecer os actos, os \emph{phenomenos}, não pode attingir o \emph{agente}, o. sujeito das faculdades, o »\emph{owmeno}. Mas isto é uma utopia. Porque, se é verdade que, só pela experiencia interna e sem Oo raciocinio e a inducção, não podemos conhecer a natureza do acto e do agente, é inteiramente falso que não possamos conhecer a \emph{existencia} do agente, que se manifesta no proprio acto. Deveremos nós fazer uma serie de raciocinios para nos convencermos de que, se operamos, existimos? O acto e o agente, se na razão são entidades distinctas, na realidade são unidas e inseparaveis, uma não se pode attingir sem a outra, e por isso a \emph{emistencia} do agente é um facto de constatação immediata e evidente. Logo o sujeito com as suas faculdades é uma coisa \emph{que se} vê, é um \emph{phenomeno}; e por isso a existencia da \emph{liberdade} pode ser conhecida pela experiencia \emph{interna}.\par \emph{A liberdade}, portanto, pode e devo ser conhecida pela experiencia \emph{interna}, pela \emph{consciencia}. --- Ora, conforme attesta a \emph{consciencia}, --- nós praticamos actos, de que somos verladeiramente senhores, e que, ficando identicas todas as condições antecedentes, podemos produzir ou deixar de produzir, --- e \textbf{tambem} escolhemcs este ou aquelle bem, e este bem ou este mal. --- Em\pgn{193}qual dos nossos actos se revela, pela consciencia, e consiste propriamente a \emph{liberdade} ? Analysemos um acto, que chamamos \emph{livre}. --- Antes de tudo, representam-se ao nosso espirito dois objectos ou bens \emph{particulares}. Neste acto não encontramos a \emph{liberdade}; a representação dos bens não depende necessariamente da nossa vontade; o primeiro acto da vida racional é um acto da intelligencia, e não da vontade. --- Mas essa representação desperta na nossa vontade umas \emph{disposições affectivas}, uma attracção ou uma repulsão, as quaes provocam alguns juizos \emph{espontaneos}, relativos \& bondade ou utilidade dos objectos. Nisto tambem não se vê ainda vestigio de liberdade. --- Em seguida, começamos a considerar esses juizos espontaneos e a reflectir sobre os objectos das representações (reflexão \emph{objectiva} ou \emph{ontologica}), e sobre os sentimentos de sympathia ou antipathia, que elles nos inspiraram (reflexão \emph{subjectiva} ou \emph{psychologica}), apreciamos esses objectos em si e nas suas consequencias; numa palavra, \emph{deliberamos}. A consciencia diz-nos que, em todo este trabalho, podemos, conforme nos apraz, --- determinar e dirigir as reflexões, concentrando a nossa attenção sobre um objecto e desviando-a de um outro, --- como tambem pôr em evidencia as boas qualidades do objecto e esconder os seus defeitos, e, viceversa, pôr em evidencia os defeitos de um outro objecto e esconder as suas perfeições, de modo que nos deixamos attrahir por um objecto, e não por outro. Nestes actos começa a sentir-se o influxo da \emph{liberdade}, mas não é na \emph{deliberação} que ella propriamente consiste. As comparações e as apreciações poderiam continuar e continuariam ao infinito (porque naquella serie de juizos \emph{reflexos}, relativos ao valor d'este ou d'aquelle objecto, não ha nenhum que possa ter sobre a vontade um imperio predominante), se não interviesse a vontade a interromper a \emph{deliberação}, levando a intelligencia ao acto da \emph{decisão}, pelo qual se julga que o objecto que, neste momento, convém, e deve ser escolhido, é este, e não é aquelle. Tal \emph{decisão} é obra, por excellencia, da \emph{liberdade}, e chama-se \emph{ultimo juizo pratico}. Este juizo --- não é \emph{necessario} (porque o bem \emph{particular} não está necessariamente ligado com o bem \emph{universal}), --- nem é \emph{evidente} (porque o desejo do bem particular não deriva \emph{necessariamente} do principio, que é o amor da felicidade), e por isso deve ser determinado pela vontade. --- Da \emph{decisão} brota espontaneamente a \emph{volição}, chamada commummente \emph{eleição}, de um bem \emph{util} á consecução do fim; e a \emph{eleição}, isto é, a acceitação de uma coisa a preferencia de uma outra, é o acto \emph{proprio} da vontade emquanto \emph{livre}. Diz S. Thomaz: « \emph{Liberum arbitrium nihil aliud est quam vis electiva} » (\emph{Sum}. \emph{Theol}., p. T, q. 83, a. 4). --- A \emph{eleição é livre}: dil-o claramente a consciencia. E' livre --- não só nos seus antecedentes, que são a \emph{deliberação} e a \emph{decisão}, --- mas tambem \emph{em si mesma}, porque a consciencia nos attesta que o bem, por nós escolhido, não o queremos \emph{necessariamente}, e que poderiamos querer outro objecto diverso, e até contrario. --- Finalmente, a \emph{eleição} exeeuta-se, e esta execução, ainda que não seja um elemento constitutivo da liberdade do nosso acto, comtudo nol-a manifesta, --- não no processo physiologico, em que se desenvolve e que é determinado, --- mas no seu principio e na sua direcção. --- São duas, pois, as causas, que concorrem para a producção do acto \emph{livre}, como facto \emph{psychologico}: uma que \emph{esclarece}, outra que \emph{move}. A causa que esclarece é a \emph{intelligencia}, a que move é a \emph{vontade}.\par t\par 18 Vor. 1L\pgn{194}Mas qual é o factor \emph{principal} do acto \emph{livre}: o juizo da intelligencia, ou o consenso da vontade? Do que acabamos de dizer se vê que o factor \emph{principal} do acto \emph{livre} é o \emph{consenso da vontade}: porque a \emph{eleição} consiste num movimento para o bem, e esse movimento é da vontade. Por isso o acto \emph{livre} pode definir-se: um acto da vontade esclarecida pela razão. Diz S. Thomaz: « Proprium liberi arbitrii est \emph{electio}. Ex hoc enim liberi arbitrii esse dicimur, quod possumus unum recipere, alio recusato, quod est eligere; et ideo naturam liberi arbitrii ex \emph{electione} considerare oportet. Ad electionem autem concurrit aliquid ex parte \emph{cognitivae virtutis}, et aliquid ex parte \emph{appetitivae}. Ex parte quidem \emph{cognitivae} requiritur \emph{consilium}, per quod dijudicatur, quid sit alteri praeferendum. Ex partó autem \emph{aypetitivae} requiritur quod appetendo \emph{acceptetur} id quod per consilium dijudicatur... Proprium objectum eleotionis est illud, \emph{quod est ad finem}. Hoc autem, inquantum hujusmodi, habet rationem boni, quod dicitur \emph{utile}. Unde cum \emph{bonum}, \emph{inquantum hujusmodi}, \emph{sit objectum appetitus}, sequitur quod \emph{electio} sit \emph{principaliter} actus \emph{appetitivae virtutis} » (\emph{Sum}. \emph{Th}., p. I, q. 83, a. 3; Cf. 1-2, q. 18; \emph{De} Ver., qq. 22, 24).\par Ouvimos dizer a S. Thomaz que o \emph{livre arbitrio} consiste na propria \emph{força de eleger}. Comtudo \emph{livre arbitrio} costuma tambem significar o principio productor do acto, que se chama \emph{eleição}, e por isso \emph{livre arbitrio} é a faculdade de \emph{escolher}, ou é a \emph{vontade emquanto escolhe livremente}. Ora dá-se uma certa proporção entre as faculdades perceptivas e as appetitivas. Assim como, na esphera da percepção \emph{intellectual}, encontramos a \emph{intelligencia} e a \emph{razão}, que não são duas faculdades distinctas, mas representam uma só faculdade com duas funcções distinctas; assim tambem, na esphera do appetite \emph{intellectual}, encontramos a \emph{vontade} e o \emph{livre arbitrio}, ou \emph{poder de eleger}, que não são duas faculdades distinctas, mas duas funcções da mesma faculdade (\emph{Sum}. \emph{Theol}., p. T, q. 88, a. 4). Com effeito, assim como o \emph{entender} significa a simples apprehens\&o da coisa (assim \emph{entendem-se} propriamente os principios de evidencia immediata) e o \emph{raciocinar} exprime a passagem de uma coisa ou verdade conhecida para uma outra coisa desconhecida (e por isso \emph{raciocinamos} Acerca das \emph{conclusões}, que se deduzem dos principios); assim tambem o \emph{querer} significa o simples appetite de uma coisa (e por isso a \emph{vontade} refere-se ao \emph{fim}, que se deseja por si mesmo) e o \emph{eleger} exprime a acceitação de uma coisa \emph{util}, isto é, que serve para alcançar uma outra coisa que nos é \emph{conveniente} (por isso a \emph{eleição} refere-se não ao \emph{fim}, que se impõe e não está sujeito á escolha, mas aos \emph{meios}, que servem para esse fim). Portanto, a vontade está para o livre arbitrio, como \& intelligencia está para a razão; porque, na esphera do appetite, o fim está para os meios, que, se desejam em razão do fim, como, na esphera da percepção, os principios estão para a conclusão, \& qual adherimos por causa dos mesmos principios. D'onde se . Segue que, assim como é impossivel o discurso ou o movimento da razão, se não o precede alguma coisa immovel, que é o principio, percebido sem raciocinio mas por simples apprehensao, assim é impossivel a \emph{eleição}, ou o movimento do appetite para os meios, se não o precede o simples appetite ou vontade do fim; porque todo e qualquer moyimento deve ter e tem o seu começo numa coisa immovel. ! \textbar{}\par Muitos escriptores negam a veracidade da consciencia, quando nos attesta\pgn{195} a liberdade da vontade. Sobresahem \emph{Bayle}, \emph{Spinosa}, \emph{Maudsley}, \emph{Stuart Mill}, \emph{Beawnis}. Examinemos o que estes dizem, e depois daremos a resposta. --- 4) \emph{Bayle} objecta: O testemunho da \emph{consciencia} em favor do \emph{livre arbitrio} não póde nem deve ser attendido; por quanto, embora não sejamos livres, como effectivamente não somos, teremos sempre o sentimento da liberdade; pois, para termos esse sentimento, basta que a vontade tenha \emph{desejo} de produzir um effeito, embora a produceüo d'esse effeito dependa de um agente exterior; assim, por exemplo (diz elle), um catavento, que tivesse desejo de virar-se para o norte e no mesmo instante fosse impellido pelo vento para aquella direcção, julgar-se-hia causa \emph{livre} do proprio movimento.\par A objeecao de \emph{Bayle} funda-se em duas falsas supposições.\par 1.9) Suppõe que o \emph{desejo} seja uma e a mesma coisa que o \emph{acto livre}. --- Ora esta supposição é falsa. O \emph{desejo} não é o \emph{acto livre}, nem a consciencia o attesta como \emph{livre}; pois o \emph{desejo} é anterior ao acto da vontade, e muitos desejos ha, a que a vontade é contraria. Por isso o catavento do apologo de \emph{Bayle}, que só tivesse \emph{desejo} de virar-se para o norte, não produziria um \emph{acto livre}, nem o sentiria como \emph{tal}; pois que, dado o \emph{desejo}, \emph{a} vontade póde deixar de consentir.\par 2.9) Suppõe que a consciencia possa attestar que um simples \emph{desejo} é causa do movimento, embora o não seja. --- Tambem esta supposição é falsa. O simples \emph{desejo} nem é causa de movimento, nem a consciencia o dá \textbf{como} tal; porque sentimos intimamente que o \emph{desejo} é inefficaz, não produzindo movimento algum, emquanto não empregarmos outros meios. Por isso o caiavento, que só tivesse desejo de se mover e que realmente se movesse, não poderia considerar-se causa \emph{livre} do movimento.\par \emph{b}) \emph{Spinoza}, seguindo \emph{Socrates}, affirmou que o homem se julga e se diz livre, porque ignora as causas que o determinaram ao acto. Assim a consciencia da nossa independencia, no entender d'este escriptor, não é outra coisa senão a inconsciencia da nossa dependencia.\par Respondemos que os factos desmentem o asserto de \emph{Spinoza}. O \textbf{sentimento} da nossa liberdade não só não provém da ignorancia das causas, que nos levam a operar, mas está precisamente na razão \emph{directa} do conhecimento d'essas causas. Porque é que nos sentimos livres? Porque, tendo considerado e examinado attentamente os bens particulares, a sua natureza, a sua relação com o nosso ultimo fim, encontramos nelles o caracter de contingencia e de limitação e nos reconhecemos desligados de todos elles; e se adherimos a um e rejeitamos outro, isso não effeito de uma attracção irresistivel, que o bem particular tenha exercido sobre a nossa vontade, mas porque assim o queremos, assim o decidimos, assim o escolhemos.\par \emph{c}) \emph{Maudsley} tambem contesta .a veracidade do testemunho da \emph{consciencia} em favor do \emph{livre arbitrio}. O raciocinio do materialista inglez é o seguinte: « Os loucos julgam-se senhores da propria vontade e dos proprios actos, quando o não são. Ora a mesma illusão póde dar-se e dá-se nos outros homens. Logo o testemunho da consciencia não deve ser attendido n.\par Resposta. O argumento de \emph{Maudsley} refuta-se, negando a egualdade dos casos. Os loucos enganam-se, quando se julgam senhores da propria vontade. E porque se enganam? porque são loucos. Mas os outros homens, que não\pgn{196}perderam o dom da intelligencia, quando se julgam senhores da propria vontade, não podem enganar-se. Ej porque? porque não são loucos; porque, como provámos na \emph{Logica}, \emph{consciencia é por si} um criterio segurissimo de verdade. --- O materialista inglez sabe perfeitamente que a excepção não destroe mas confirma a regra. Por isso, assim como das illusões, a que póde estar sujeito um sentido no estado anormal e pathologico, não podemos concluir que todos os sentidos de todos os homens soffrem as mesmas illusões; assim tambem de um engano, a que póde estar sujeita a \emph{consciencia} de um louco, não podemos concluir que a \emph{consciencia} de todos os homens soffre o mesmo engano. --- De resto, não é verdade que os loucos sintam sempre que são senhores da propria vontade; pois, em muitos casos, aquelles infelizes advertem que estão sob o dominio de uma força irresistivel, e nos intervallos lucidos deploram a propria desgraça e pedem perdão dos actos maus, praticados sob o impulso d'aquella força.\par \emph{d}) \emph{Stuart Mill} diz que a \emph{consciencia} do livre arbitrio é uma coisa impos- \&ivel. Eis as suas palavras: « Ter consciencia do livre arbitrio significa ter, antes da decisão, consciencia de se poder decidir de outro modo. Ora isso é impossivel. A consciencia não me diz se não o que eu faço ou o que sinto. \emph{Quanto ao que eu posso fazer}, isso não cae sob o dominio da consciencia. Se nós podemos ter consciencia de uma força, se podemos sentir uma aptidão \emph{independentemente de todo} o \emph{exercicio presente ou passado d'essa aptidão}, tal consciencia é um facto unico sem analogia alguma em toda a natureza ». (\emph{Exame da philosophia de Hamilton}, c. 26).\par Resposta. A diffieuldade de \emph{Stuart Mill} é tão falsa como superficial. E' falso que a consciencia do livre arbitrio seja independente do seu exercicio. Temos consciencia da nossa liberdade, por isso mesmo que praticamos alguns actos só porque assim o queremos, --- só porque a nossa vontade transforma uma \emph{sollicitacüo} em \emph{determinação}. --- E, embora a \emph{capacidade} de fazer ou deixar de fazer não possa \emph{directamente} cahir sob o alcance da consciencia, cae \emph{indirectamente}, emquanto, reflectindo sobre um acto praticado, percebemos que esse acto foi praticado \emph{livremente}, e que \emph{por isso} tinhamos a \emph{capacidade} de o praticar ou de o não praticar.\par e) \emph{Beaunis} deixa em paz os loucos, recorre aos \emph{hypnotisados}, e diz: « Os hypnotisados julgam-se \emph{livres}, e comtudo operam machinalmente. Dá-se o mesmo com o resto dos homens; julgam-se livres, mas não o sãon. E, em prova do seu asserto, adduz muitos factos de hypnotisados, que dizfim ser\par livres, quando o não eram.\par Uma breve resposta. --- Antes de tudo, muitos sabios, especialmente \emph{Delbeuf}, contestam os factos apresentados por \emph{Beaunis} e dizem não ser exacto que os hypnotisados se julguem perfeitamente livres. --- Além d'isso, uma profunda observação dos phenomenos hypnoticos, relativos a este assumpto, deu os seguintes resultados, geralmente admittidos: --- 1.º) na maior parte dos casos, os hypnotisados affirmam ignorar o motivo que os levou a operar, e não sabem dizer se o seu acto foi deliberado, ou não; --- 2.º) muitas vezes os hypnotisados sentem-se possuidos e como que sitiados pela idea fixa da suggestão, que os determina a um certo acto, não se sentem senhores de si mesmos e a sua vontade parece impotente para resistir; --- 9.9) ás vezes\pgn{197}dizem que são \emph{livres} e \emph{não-livres}, e esta, affirmação é verdadeira, porque sentem-se livres no \emph{principio} da suggestão da idea fixa, \emph{e} não se sentem e não são livres no \emph{ultimo periodo}, por não terem empregado os meios necessarios para resistir ao ascendente imperioso da idea, e recuperar a propria independencia. --- Finalmente, admittindo embora a verdade dos factos apontados por \emph{Beaunis}, qual é á consequencia legitima, que d'elles devemos deduzir? A unica consequencia é que a consciencia humana, no seu estado pathologico e anormal, como é a dos hypnotisados, pode estar sujeita a illusões. Mas consentirá o senso commum que a uma faculdade, considerada no seu estado normal e são, se applique o que só lhe acontece num estado anormal e pathologico, e em virtude d'este mesmo estado? Não, por certo. A. illusao dos hypnotisados é uma \emph{excepção}. Se os doentes se enganam, é porque se.\par t não enganam os sãos. Logo a illusao dos iypnotisddos não só não destroe, mas confirma \& veracidade da consciencia.\par Concluimos que a prova da existencia da \emph{liberdade} pelo testemunho da consciencia é apodictica, convincente. Ella completa a prova da consciencia \emph{moral}, já exposta, e é-lhe immensamente superior; porque, se a prova \emph{oral}. demonstra que a liberdade \emph{deve existir}, \& prova pela \emph{consciencia} constata que a liberdade \emph{existe} realmente. Por isso, a prova da \emph{liberdade} pela consciencia foi sempre tida como a \emph{pedra angular} de todo o edificio moral.}. \textbar{}

d) A vontade está para o bem, como a. intelligencia está para a verdade, Ora á intelliveneia, se adhere necessariamente aos primeiros principios da sciencia, é livre em aeceitar as conclusões particulares e contingentes, que não estão \emph{necessariamente} \pg{193}  ligadas com aquelles principios; porque essas conclusões não se impõem á mesma intelligencia. Logo a vontade, \pg{194}  se tende necessariamente para o bem \emph{universal}, é comtudo livre em seguir ou recusar os bens \emph{particulares}; pois que

------ - --- \emph{TRIN} » ds \emph{UNT}

\emph{SECÇÃO} II --- CAPITULO III --- FACULDADES INTELLECTUAES DO HOMEM 195 esses bens, por não satisfazerem completamente a capacidade immensa da vontade, e por não estarem ou por se não \textbf{conheeerem} \pg{196}  \emph{necessariamente} ligados com o bem \emph{universal}, não attrahem de um modo irresistivel a nossa vontade, mas deixam-na \pg{197}  sempre indifferente. Logo a vontade é livre na sua tendencia para os bens \emph{particulares} \footnote{Os dados da consciencia \emph{moral} e o testemunho da consciencia \emph{psychologica} demonstram a \emph{necessidade} e o \emph{facto} da existencia da \emph{liberdade}. Não basta: é preciso investigar a \emph{razão intrinseca} d'esse facto. Tal trabalho pertence á \emph{metaphysica}.\par A prova \emph{metaphysica} não gosa da sympathia dos modernos escriptores, mesmo dos mais ardentes defensores da \emph{liberdade}. \emph{Renowvier} diz que n\&o é possivel encontrar uma prova \emph{a priori} em favor da liberdade e accrescenta que « demonstrar pelo syllogismo a \emph{necessidade da liberdade} seria o mesmo que negal-a ». \emph{l'onsegrive}, outro habil espiritualista, acha justissima a observação de \emph{Renowvier}. Mas esta observação não é seria; funda-se num jogo de palavras. Demonstrar a \emph{necessidade da liberdade} não é dizer que a \emph{liberdade é necessidade}. --- A prova \emph{metaphysica} existe, e o nosso Angelico Doutor tem-na exposta com tal profundidade e evidencia, que só não illumina os que não querem ver (Cf. \emph{Sum}. \emph{Th}., p. I, q. 82, a. 2; p. I-II, qq. 9-10; QQ. DD. \emph{De malo}, q. 6, a. un.; \emph{De verit}, q. 22).\par Como vimos no numero anterior, a vontade não é livre na sua tendencia para o bem \emph{universal}, para a \emph{felicidade}. Essa tendencia é natural, e por isso é \emph{necessaria}. Mas esta mesma tendencia \emph{necessaria} da vontade para o bem \emph{universal}, para a \emph{felicidade}, é o ponto de apoio, á razão ultima dos nossos actos livres, como o absoluto é a razão do relativo, e a causa immovel é a razão ultima do movimento. Diz S. Thomaz: « Finis se habet in operativis, sicut principium in speculativis. Oportet enim quod illud, quod naturaliter alicui convenit, et immobiliter, sit fundamentum et principium omnium aliorum; quia natura rei est primum in unoquoque, et omnis motus procedit ab aliquo immobili ». (\emph{Sum}. \emph{Th}., p. T, q. 82, a. 1). O principio immovel, de que falla S. Thomaz, é a felicidade, que é a causa de todas as nossas volições: « Omnis homo naturaliter vult beatitudinem, et ex hac naturali voluntate\pgn{198} (in finem ultimum seu beatitudinem) \emph{causantur omnes aliae volitiones}, cum quidquid homo vult, velit propter finem. Dilectio igitur boni, quod \textbf{homo} naturaliter vult sicut finem, est dilectio \emph{naturalis} (seu \emph{necessaria}); dilectio autem ab hac derivata, quae est boni quod diligitur propter finem, est dilectio \emph{electiva} (seu \emph{libera}) ». (\emph{Sum}. \emph{Th}., p. I, q. 60, a. 2).\par Tendendo irresistivelmente para o bem universal, a vontade sente-se inclinada para aquellas entidades particulares, que contém um reflexo, uma participação do bem universal. A intelligencia distingue immediatamente em todos estes bens particulares a contingencia, a imperfeição; todavia, postos em relação com o fim, alguns d'elles têm com o proprio fim um nexo necessario, outros um nexo apenas contingente. E a vontade, seguindo o juizo da intelligencia, em virtude da sua tendencia necessaria para o fim, abracará tambem necessariamente os bens particulares que têm, e nós conhecemos terem, um nexo necessario com o proprio fim; mas em virtude da mesma tendencia abragará ou rejeitará, a seu bello prazer, aquelles bens particulares que não têm, ou não conhecemos terem, um nexo necessario com a felicidade. --- Os bens particulares são \emph{bens imperfeitos}. Emquanto sao \emph{bens}, e podem conduzir ao fim, podem attrahir a vontade; emquanto são \emph{imperfeitos}, e não conduzem necessariamente ao fim, não a attrahem \emph{necessariamente}. O movimento só então se segue necessariamente, quando a força do movente é superior á potencia do movel, de modo que toda a capacidade d'este seja dominada pela força d'aquelle; mas jamais uma força inferior poderá mover uma superior, se esta não cede livremente. Ora, como todos os bens particulares possuem uma força immensamente inferior á capacidade da vontade humana, nenhum d'elles que não tenha, ou não se reconheça ter, connexao necessaria com o fim ultimo, póde mover a vontade, se esta não se rende livremente por si mesma. E, quando se rende, não se póde dizer que opére sem motivo, porque dar-se-hia um acto sem raz\&o sufficiente, e até sem o seu objecto; mas nem se deve dizer, que o motivo seja \emph{causa} efficiente, e muito menos \emph{determinante}. Em poucas palavras: a vontade tende necessariamente para o fim ultimo, e este é a razão, o motivo unico necessario de toda a eleição. Ora o fim ultimo póde subsistir sem cada um dos varios bens particulares. Portanto, com relação a cada um dos varios bens particulares, a vontade é indifferente e livre.\par E'esta a doutrina do nosso Angelico Mestre, Primeiramente S. Thomaz prova que os bens particulares não podem encher e saciar a capacidade da vontade, e por isso não a determinam necessariamente: « Movens tuno ex necessitate causat motum in mobili, quando potestas moventis excedit mobile, ita quod tota ejus possibilitas moventi subdatur. Cum autem possibilitas voluntatis sit respectu boni universalis et perfecti, non subjicitur ejus possibilitas tota alicui particulari bono, et ideo. non ex necessitate movetur ab illo » (\emph{Sum}. \emph{Theol}., p. I, q. 82, a. 2 ad 2). --- Depois o S. Doutor prova a liberdade da vontade humana com relação aos bens \emph{particulares}. Diz assim: \& Sicut intellectus naturaliter et ex necessitate inhaeret primis principiis, ita voluntas ultimo fini. Sunt autem quaedam intelligibilia, quae non habent necessariam connexionem ad prima principia, sieut contingentes propcsitiones\pgn{199}, ad quarum remotionem non sequitur remotio primorum principiorum; et talibus non ex necessitate assentit intellectus... Similiter etiam ex parte voluntatis. Sunt enim quaedam particularia bona quae non habent necessariam connexionem ad beatitudinem, quia sine his potest aliquis esse beatus; et hujusmodi. bonis voluntas \emph{non em necessitate} inhaeret » (\emph{Sum}. \emph{Theol}., p. I, q. 82, a. 2). E noutro logar diz: \& In omnibus particularibus bonis potest (homo) considerare rationem boni alicujus, et defectum alicujus boni quod habet rationem mali; et secundum hoc potest unumquodque hujusmodi bonorum apprehendere ut eligibile, vel fugibile. Solum autem perfectum bonum, quod est beatitudo, non potest ratio apprehendere sub ratione mali aut alicujus defectus; et ideo \emph{em necessitate} beatitudinem homo vult, nec potest velle non esse beatus, aut esse miser. Electio autem, cum non sit de fine, sed de his quae sunt ad finem, non est perfecti boni, quod est beatitudo, sed aliorum particularium bonorum. Et ideo homo \emph{non e}« \emph{necessitate}, \emph{sed libere eligit} ». (\emph{Sum}. \emph{Th}., 1-9, q. 18, a. 6; cf. p. I, q. 88, a. 1; 1-2, q. 10, a. 2).\par Da doutrina exposta se deduz claramente que a verdadeira \emph{rais} da liberdade é a \emph{razão}, emquanto percebe o bem \emph{universal}. Só o homem entre todas as creaturas visiveis é dotado de livre arbitrio, porque só elle é \emph{racional}. Diz S. Thomaz: « Necesse est quod homo sit liberi arbitrii ex hoc ipso quod \emph{rationalis} est » (\emph{Sum}. \emph{Th}., p. I, q. 88, a. 1). De facto, como observa o mesmo Santo Doutor (\emph{l}. \emph{c}.), ha tres ordens de agentes. --- Uns são destituidos de toda a especie de conhecimento, operam só em virtude da sua forma natural, e a sua inclinação é determinada a uma coisa só, e não podem operar d'outro modo: taes são as plantas. --- Outros têm um conhecimento exclusivamente \emph{sensitivo} e são dotados de juizos \emph{instinctivos}, operam em virtude de uma forma, que foi apprehendida pelos sentidos e representa coisas singulares e materiaes, e por isso, como essa forma é singular e esse juizo é fixo e determinado pela natureza, a sua inclinação não pode tender senão para um certo e determinado acto; taes são os animaes irracionaes. O animal não é capaz de conhecer nem de desejar o bem \emph{universal} (conhecimento e desejo, que constituem a base de toda e qualquer \emph{eleição}), --- por isso, nem é capaz de escolher entre os proprios bens creados e materiaes, que apprehende pelos sentidos. Os seus conhecimentos e as suas tendencias estão todas circumscriptas na esphera das coisas materiaes e concretas, e os bens particulares, que convém actualmente ás suas disposições, attrahem-no necessariamente e satisfazem a sua capacidade subjectiva. --- Outros são dotados de conhecimento \emph{intellectual}, possuem ideas \emph{universaes}, \emph{e} operam em virtude d'essas formas \emph{universaes}; e, como a forma \emph{universal} pode tornar-se concreta e traduzir-se na realidade por muitos e differentes modos \emph{particulares}, a sua inclinação não está obrigada e determinada a um modo \emph{particular}, mas, entre os muitos e varios modos, pode escolher o que mais lhe apraz (assim o architecto, que possue a forma \emph{universal} de um palacio, a qual contém diversas figuras de palacios, pode escolher a figura redonda a preferencia de outra); taes são os homens. O homem concebe, pela razão, a idea do bem \emph{universal}, e só por este bem é attrahido \emph{necessariamente}, porque só elle sacia toda a sua capacidade. Se só o bem \emph{universal} attrahe \emph{necessariamente} o homem, tudo o que não é esse bem, ou a elle não se refere\pgn{200}necessariamente, deixa a vontade indifferente e livre na escolha d'este ou d'aquelle bem \emph{particular}; porque em todos esses bens particulares a razão encontrará sempre o caracter de contingencia e de limitação, e por isso a vontade sentir-se-ha sempre desligada de cada um d'elles, e se esta faculdade adhere a um e rejeita outro, tal maneira de operar não é effeito de necessidade, mas de livre escolha. --- E quanto mais penetrante fôr a razão e mais claramente descobrir a contingencia e os defeitos dos bens creados ou particulares, e melhor conhecer que os attractivos d'elles não têm nada de necessitante, tanto mais independente e livre se sentirá a vontade em presença d'esses attractivos. Podemos, pois, afirmar que a \emph{razão} não só é a \emph{condição} ou a \emph{raiz} da liberdade, mas é tambem a sua \emph{medida}; porque a vontade é tanto mais livre, quanto mais esclarecida é a escolha (relativamente aos bens finitos). D'aqui se vê que ha mais liberdade no homem justo do que no peccador; mais no Anjo do que no justo; mais, infinitamente mais, em Deus do que em todos os Anjos e Santos. Em Deus, a liberdade é infinita, porque infinito é o seu entender e o seu querer. Elle pode querer e não querer os bens creados, e esta liberdade não deriva da imperfeição ou contingencia da sua vontade, mas sim da imperfeição e contingencia das coisas creadas, das quaes nenhuma é necessaria para completar a sua absoluta e infinita felicidade. --- Concluimos que a razão é á \emph{rais} ou a \emph{causa} da liberdade. Diz a este respeito S. Thomaz que, se a \emph{vontade} pode chamar-se a \emph{raiz} da liberdade emquanto é \emph{sujeito} d'ella, \emph{a rasdo} deve dizer-se \emph{raiz} da liberdade emquanto é \emph{causa} d'ella. « \emph{Radix} libertatis est \emph{voluntas} sicut \emph{subjectum}; sed sicut \emph{causa}, est \emph{ratio}; ex hoc enim voluntas libere potest ad diversa ferri, quia ratio potest habere diversas conceptiones boni » (\emph{Sum}. \emph{Th}., 1-9, q. 17, a. 1 .ad 2; cf. C. \emph{Gentes}, 1. III, c. 78; \emph{De verit}., \emph{q}. 24, aa. 1-2).\par Completemos agora o estudo ácerca das analogias, que existem entre o operar da intelligencia e o operar da vontade; pois só algumas foram expostas nesta e na these anterior. --- A vontade está para o seu ultimo fim, para a felicidade, como a intelligencia está para os primeiros principios da sciencia: a felicidade é na pratica aquillo que os primeiros principios são na theoria. --- Assim como a intelligencia, de sua natureza, necessariamente e sem discurso adhere aos primeiros principios, porque a verdade evidente satisfaz toda a capacidade da intelligencia humana; assim tambem, por sua natureza, necessariamente e sem eleição, a vontade adhere ao seu ultimo fim. --- Assim como a intelligencia, se fossem supprimidos todos os principios, não poderia de forma alguma chegar ao acto de qualquer conclusão; do mesmo modo seria impossivel á vontade fazer qualquer eleição, quando cessasse para ella o ultimo fim; porque a eleição não se verifica, se falta aquillo para que se escolhe. --- Assim como uma intelligencia a necessidade de admittir as conclusões depende da connexàüo necessaria e evidente que estas têm com os principios, de sorte que não se podem negar as conclusões, sem negar ao mesmo tempo os principios; assim na vontade a necessidade de querer um bem particular depende do nexo necessario e evidente entre aquelle bem e o ultimo fim, de modo que não póde deixal-o sem renunciar simultaneamente ao proprio fim. --- Pelo contrario, assim como a intelligencia, sem renunciar aos primeiroes principios, é livre em acceitar ou rejeitar\pgn{201}as conclusões particulares que não têm um nexo necessario ou evidente com os mesmos principios e póde escolher d'entre ellas uma meramente provavel, deixando de lado a mais provavel, porque nenhuma se lhe impõe necessariamente; assim tambem a vontade, adherindo sempre e necessariamente ao seu ultimo fim, é livre em querer os bens particulares, quando a connexão d'estes com o ultimo fim, não é necessaria, ou não é evidente, e póde determinar-se por um motivo menos forte em presença d'um outro mais forte, porque nenhum dos dois exerce sobre ella um attractivo necessario.}.

%%>>> t132

\pg{191} \pg{195} \pg{198} \tese{132}{A vontade é movida por Deus para o bem universal, mas move-se a si mesma para os bens particulares.}

a) A \emph{vontade é movida por Deus para o bem universal}. --- \pg{199}  A vontade tende, por um movimento necessario e irresistivel, para o bem \emph{universal}, que é a felicidade. Ora a tendencia da \pg{200}  vontade para o bem \emph{universal} não pode derivar senão do \textbf{impulso} do Bem \emph{universal}, que é \textsc{Deus}; pois cada agente produz \pg{201}  um movimento proporcionado á sua natural energia. Logo a vontade é movida por Deus para o bem \emph{universal} \footnote{E' necessario, para a exacta comprehensao da these, distinguir um duplice movimento: \emph{objectivo} e \emph{suljectivo}. \emph{O objectivo} é necessario para a \emph{especificação} do acto; o \emph{subjectivo}, para o \emph{exercicio} do proprio acto. --- O movimento, que \emph{especifica} o acto, deriva do proprio \emph{objecto} (e é por isso que se diz \emph{objectivo}), e assim o acto da vontade é \emph{necessario} ou \emph{contingente}, conforme \emph{necessario} ou \emph{contingente} (\emph{particular}) é o objecto, para o qual ella tende. --- O movimento, que determina o \emph{sujeito} ao \emph{exercicio}, isto é, á \emph{pratica} do acto, quer este acto seja \emph{necessario} quer \emph{contingente} ou \emph{livre}, de qual principio deriva? Esta é a questão.\par E' innegavel que a nossa vontade tende necessariamente para o bem \emph{universal}. Esta tendencia, que é um effeito, deve ter uma causa, um principio proporcionado. Qual é este principio? --- Será a propria alma? Não; porque a alma está \emph{em potencia} para o bem \emph{universal}, \emph{e} um ente que está \emph{em potencia} para um bem não pode passar ao \emph{acto}, isto é, não pode \emph{actualmente} tender para esse bem, senão pela força de outro ente, que esteja \emph{em acto}, isto é, de outro ente que seja esse \emph{acto}, que seja esse bem, ou que possua esse \emph{acto}, que tenda \emph{actualmente} para esse bem. --- Será um bem \emph{creado}, \emph{particular} ? Tambem não; porque nenhum \emph{agente} pode imprimir um movimento \emph{superior} á propria energia (pois a força e a qualidade de um movimento não pode exceder a força e a qualidade do agente, que o produziu), e tal aconteceria, se o bem \emph{particular} imprimisse á vontade uma tendencia para o bem \emph{universal}. Logo só \textsc{Deus}, que, além de ser a Causa creadora da vontade humana, é o Bam \emph{universal} (não abstracto mas concreto), pode mover a nossa vontade para o bem \emph{universal} (como principio \emph{extrinseco}). Santo Thomaz à questão: « utrum voluntas moveatur a Deo solo ut ab exteriori principio » responde: \& Voluntatis causa nihil aliud esse potest quam Deus. Et hoc patet dupliciter. Primo quidem ex hoc quod voluntas est potentia animae rationalis, quae a solo Deo causatur per creationem ... Secundo vero ex hoc quod voluntas habet ordinem ad \emph{wniversale bonwm}; unde nihil aliud potest esse voluntatis causa, nisi ipse \emph{Deus}, \emph{qui est universale bonum}. Omne autem. aliud bonum per participationem dicitur, et est quoddam \emph{particulare} bonum; \emph{particularis autem causa non dat inclinationem}. \emph{universalem} » (S. \emph{T}., 1-2, q. 9, a. 6).}.

\emph{b}) A \emph{vontade move-se a si mesma para os bens particulares}. --- Uma faculdade, que tende \emph{actualmente} para um objecto \emph{universal}, pode mover-se a si mesma para os objectos \emph{particulares}, que estão de algum modo contidos na extensão d'esse objecto \emph{universal}; assim a intelligencia passa naturalmente, pela sua propria energia, de um principio conhecido para uma conclusão \pg{202}  desconhecida, mas virtualmente contida naquelle principio. Ora a vontade tende \emph{actualmente}, em virtude do movimento recebido de \textsc{Deus}, para o bem \emph{universal}, que contém \emph{eminentemente} todos os bens \emph{particulares}. Logo á vontade pode \textbf{mover-se} e move-se a si mesma para os bens \emph{particulares} \footnote{A vontade do homem é movida, \emph{no principio} e \emph{necessariamente}, pelo proprio Auctor da natureza para o bem universal, que é a felicidade. Sem esta moção primitiva, o homem nao poderia ter nenhum desejo; porque a vontade nunca poderia, por si só, passar da potencia para o acto. Posta porém uma vez em acto, ou em movimento, para o bem universal, já não repugna que ella se mova e se determine a si mesma, sob a direcção da intelligencia, para este ou aquelle bem particular, real ou apparente. Ella tem em si mesma o seu motor: aquelle primeiro acto \emph{necessario} excita a faculdade a produzir outros actos \emph{livres}. Diz S. Thomaz: « Deus movet voluntatem hominis, sicut universalis motor ad universale objectum voluntatis, quod est bonum, et sine hac universali motione homo non potest aliquid velle; sed homo per rationem determinat se ad volendum hoc vel illud, quod est vere bonum, vel apparens bonum » (\emph{Sum}. \emph{Theol}., 1-2, q. 9, a. 6 ad 3). D'este modo o homem é movido por Deus como um instrumento, que, todayia, se moye a si mesmo pelo livre arbitrio: porque um ente pode ser movido por outro, e comtudo pode mover-se a si mesmo. --- Portanto, o processo da vontade não consiste na passagem espontanea da potencia para o acto (o que é impossivel), mas na passagem do primeiro acto para o segundo, para o terceiro, e por isso de egual para egual, ou, melhor, do mais para o menos, porque a volição do bem \emph{universal} contém \emph{eminentemente} a volição dos bens \emph{particulares}. --- Por isso, assim como a intelligencia não se poderia mover a Bi mesma de um primeiro principio para as conclusões, se não conhecesse actualmente esse principio, mas, conhecido o primeiro principio, pode por si mesma passar ao conhecimento das conclusões, virtualmente contidas no principio; assim tambem a vontade não poderia mover-se nem tender para este ou aquelle bem particular, se não tendesse «\emph{actualmente} para o bem \emph{universal}, mas, posta esta tendencia para o bem \emph{universal}, pode, por si mesma, mover-se e tender para os bens \emph{particulares}, contidos na comprehensão do bem \emph{universal}. E' o que diz S. Thomaz: « Hoc modo se habet finis in appcetibilibus, sicut principium in intelligibilibus. Manifestum est autem quod intellectus per hoc quod cognoscit principium, reducit seipsum de potentia in actum, quantum ad cognitionem conclusionum, et hoc modo movet seipsum; et similiter voluntas per hoc quod vult finem, movet seipsam ad volendum ea quae sunt ad finem » (\emph{Sum}. \emph{Theol}. 1-9, q. 9, a. 8).\par E' certo que tambem esta tendencia \emph{livre} do homem para os bens \emph{particulares}, por ser acto de uma causa \emph{segunda}, exige o concurso da Causa \emph{primeira}. Mas esse concurso não tira á vontade o poder de se mover a si mesma, nem torna \emph{necessario} o acto \emph{por si livre}. Deus accomoda o seu concurso á natureza dos agentes creados, e assim como move os agentes \emph{necessarios} por modo que elles produzam \emph{necessariamente} os seus actos, assim \textbf{move} os agentes \emph{livres} por modo que elles produzam os seus actos \emph{livremente}. Sendo, pois, á nossa vontade um agente \emph{livre}, isto é, não determinado para um acto só, mas indifferente para muitos, é claro que Deus a\pgn{203}move, não emquanto a determina a produzir necessariamente um acto, mas emquanto auxilia aquella indifferença, isto é, o livre arbitrio do homem para que possa livremente produzir o seu acto. Diz S. Thomaz: « Ad Providentiam divinam non pertinet naturam rerum corrumpere, sed seryare. Unde omnia movet secundum eorum conditionem; ita quod ex causis \emph{necessariis} per motionem divinam sequuntur effectus \emph{ew necessitate}, ex causis autem \emph{contingentibus} sequuntur effectus \emph{contingentes}. Quia igitur voluntas est activum principium \emph{non determinatum ad wnwm}, sed \emph{indifferenter} se habens ad multa; sic Deus ipsam movet, quod non ex necessitate ad unum determinat, sed remanet motus ejus \emph{contingens}, et non necessarius, nisi in his ad quae naturaliter movetur » (\emph{Sum}. \emph{Theol}, 1-2, q. 10, a. 4).\par D'atjui se vê a superficialidade, o equivoco e a falsidade do argumento, com que os sequazes de \emph{Kant} julgam destruir o \emph{livre arbitrio}. Dizem, mais ou menos, assim. Todo o ente deve ter a sua rasdo \emph{sufficiente}. Ora o acto \emph{livre}, se fosse possivel, não teria \emph{razão sufficiente}; porque a vontade, sendo por si indifferente e não tendo recebido uma moção precedente, passaria \emph{espontaneamente} da potencia para o acto, e esta passagem repugna. Logo o acto \emph{livre} é um absurdo. --- E' claro que estes escriptores negam o movimento \emph{necessario} impresso por Deus á vontade para o bem \emph{universal}. Se o admittissem, como devem admittir, encontrariam immediatamente a \emph{razão sufficiente} da passagem, que a vontade faz do bem \emph{universal} para os bens \emph{particulares}. Como tambem é evidente, elles confundem a indifferenca \emph{passiva}, que é a \emph{capacidade} de receber a determinação de um agente, com a indifferença \emph{activa}, que é a força, de que um ente é dotado para, se determinar a si mesmo, sahindo da indifferenca e produzindo o acto. Ora a liberdade consiste na indifferenga \emph{activa}, que não cria nem communica a si mesma o movimento, mas utilisa o movimento recebido, e o continua a seu modo.\par Objectam: Se a vontade se movesse a si mesma, estaria, ao mesmo tempo, em \emph{potencia} e \emph{em acto}. Ora isto repugna. Logo.\par Resposta. Movendo-se a si mesma, a vontade esta \emph{em potencia} e \emph{em acto} ao mesmo tempo, mas \emph{sob diversos aspectos}. Está \emph{em acto} com relação ao bem \emph{universal}; está \emph{em potencia} com relação aos bens \emph{particulares}. Estando \emph{em acto}, pode por si passar do acto primeiro para o segundo, para o terceiro, isto é, pode passar do \emph{actual} desejo do bem \emph{universal} para o desejo dos bens \emph{particulares}. D'este modo a vontade, emquanto está \emph{em acto}, move-se a si mesma, emquanto está \emph{em potencia}. Isto não repugna, porque, para haver repugnaneia é preciso que uma coisa esteja \emph{em potencia} e \emph{em acto} ao mesmo tempo e \emph{sob o mesmo} «\emph{aspecto}. Diz S. Thomaz: « Voluntas non secundum idem movet et movetur, unde nec secundum idem est in actu et in potentia; sed inquantum \emph{actu} vult finem, reducit se de \emph{potentia} in\pgn{204}actum respectu eorum quae sunt ad finem, ut scilicet actu ea velit » (\emph{Sum}. \emph{Theol}., 1-2, q. 9, a. 8 ad 1).\par Insistem: A vontade humana está \emph{em potencia} com relação ao desejo dos bens \emph{particulares}, \emph{e} por isso não se pode mover por si mesma a desejar \textbf{esses} bens, mas deve ser movida por outra causa: porque o ente em \emph{potencia} não pode ser reduzido \emph{ao acto} senão pelo ente \emph{em acto}. Logo a vontade não pode mover-se a si mesma para os bens \emph{particulares}.\par Resposta. Esta difficuldade resolve-se como a precedente. A nossa vontade, com relação aos bens \emph{particulares}, não está inteiramente \emph{em potencia}, mas em parte está em \emph{acto}, \emph{e} em parte está \emph{em potencia}, e em quanto está em \emph{acto} move-se a si mesma emquanto está \emph{em potencia}. Está em \emph{acto} com relação ao bem \emph{universal}, porque deseja \emph{actualmente o} bem \emph{universal},. que enche toda a sua capacidade subjectiva, e está em potencia com relação aos bens \emph{particulares}, porque \emph{pode} desejar os bens \emph{particulares}, que não enchem \& sua capacidade; e por isso pode mover-se a si mesma, e o seu movimento fica sempre \emph{contingente} e \emph{livre}. Diz S. Thomaz: « Um ente não pode ser motor e movido sob o mesmo respeito, mas pode-o ser sob respeitos differentes; e assim como a intelligencia, emquanto percebe actualmente os principios, se reduz a si mesma ao acto emquanto ás conclusões, assim tambem a vontade, emquanto quer o fim, se reduz a si mesma ao acto emquanto aos meios, que levam a esse fim » (QQ. DD., \emph{De malo}, q. 6). Por isso a nossa vontade, emquanto se move por si mesma da potencia para o acto, distingue-se --- tanto dos agentes \emph{necessarios}, que são meramente \emph{passivos}, como da Causa primeira, que é acto \emph{purissimo}. (QQ. DD. \emph{De verit}., q. 92, a. 4).\par Concluindo, agradecamos a Deus o dom do \emph{livre arbitrio}, que é o maximo entre todos os dons, porque mais semelhantes nos torna á Bondade infinita. O \emph{livre arbitrio} esta fundamentalmente orientado, não para o \emph{mal}, mas para o \emph{bem}, --- para o bem \emph{universal}, que se deve alcançar pelo uso dos bens \emph{particulares}. Se, ás vezes, segue um bem \emph{apparente}, que na realidade é um \emph{mal}, a culpa não é d'elle (pois elle soffre com isso violencia), mas é das da nossa \emph{ignorancia}, da nossa \emph{maldade}, ou, melhor, da \emph{natural imperfeição}, que acompanha e penetra as nossas faculdades e o nosso proprio ser. Oitamos, a este respeito, um magnifico trecho de S. Thomaz, que resume, em parte, a doutrina exposta neste artigo. Diz o S. Doutor: « Ad rationem liberi arbitrii non pertinot ut indeterminate se habeat ad bonum vel ad malum, quia liberum arbitrium per se in bonum ordinatum est, cum bonum sit objectum voluntatis, nec in malum tendat nisi propter aliquem defectum, quia apprehenditur ut bonum, cum non sit voluntas aut electio nisi boni, aut apparentis boni; et ideo ubi perfectissimum est liberum arbitrium in malum tendere non potest, quia imperfectum esse non potest, sed hoc ad libertadem arbitrii pertinet, ut actionem aliquam facere vel non facere possit » (\emph{In Sent}., dist. 25, q. 1, a. 1 ad 2; \emph{Of}. \emph{Sum}. \emph{Th}., p. 1, q. 69, a. 8 ad 8).}.

%%>>> art-19

\section*{Artigo III.}
\addcontentsline{toc}{section}{Artigo III.}

\pg{203} A \textbf{liberdade da vontade humana}

e o \textbf{determinismo}.

%%>>> t133

\tese{133}{Determinismo.}
\emph{Determinismo}. 6 o systema, que nega a liberdade da vontade humana e sustenta que \emph{todos} OS nossos actos, mesmo os chamados \emph{livres}, são o resultado fatal dos \pg{204}  seus antecedentes. --- Segundo este systema, pois, o nexo entre os nossos actos e os seus antecedentes é tão necessario, que não pode ser nem quebrado nem alterado \footnote{\emph{Determinismo} é uma palayra ambigua. Significa, ás vezes, o conjuncto das causas sufficientes para a producção de um phenomeno; e assim se estabelece o \emph{determinismo} de um facto physiologico, quando se indicam todos os antecedentes necessarios d'esse facto, de modo que se pode predizer\pgn{205} com certeza a realisação do facto, desde que estejam realisados os antecedentes. --- Mas, no sentido rigoroso, \emph{determinismo} é um systema philosophico relativo á liberdade humana.}.

%%>>> t134

\pg{205} \tese{134}{Especies de determinismo.}
O \emph{determinismo} divide-se, principalmente, em \emph{psychologico}, \emph{mechanico}, \emph{physiologico} e \emph{theologico}, conforme os antecedentes, que teriam o poder de determinar, por um modo irresistivel, á vontade, pertencem á ordem \emph{psychologica}, ou \emph{mechanica}, ou \emph{physiologica}, ou \emph{theologica}.

a) \emph{Determinismo psychologico}. --- O determinismo \emph{psychologico} ensina que os antecedentes dos nossos actos são constituidos pelo \emph{motivo mais forte}, isto é, pelo \emph{bem melhor}; de modo que a nossa vontade deve necessariamente obedecer ao \emph{motivo mais forte}, isto é, seguir o \emph{bem melhor} \footnote{O determinismo \emph{psychologico} foi inventado por \emph{Leibnitz} e é seguido por \emph{Littré}, \emph{Maudsley}, \emph{Herzen}, \emph{Herbert Spencer}, e por todos os positivistas.}.

\emph{b}) \emph{Determinismo mechanico}. --- O determinismo \emph{mechanico} sustenta que os antecedentes dos actos da nossa vontade são os movimentos \emph{mechanicos}, e que por isso a nossa vontade, como qualquer outra força material, deve estar sujeita, na producção de seus actos, ás leis fataes da mechanica \footnote{O determinismo \emph{mechanico} é defendido por \emph{Fouillé}, \emph{Du Bois-Reymond}, \emph{Vogt}, \emph{Büchner}, \emph{Helmolts}, \emph{Molleschott}, \emph{Huxley}, e por todos os materialistas.}.

c) \emph{Determinismo physiologico}. --- O determinismo \emph{physiologico} admitte que os antecedentes dos actos da nossa vontade consistem numa sensação \emph{involuntaria}, ou numa modalidade \emph{organica}, transmittida pela hereditariedade; de modo que todos os actos da nossa vontade, mesmo os chamados \emph{livres}, se reduzem a movimentos \emph{reflexos}, ou são a resultante fatal d'aquella modalidade \emph{organica} \footnote{O principal defensor deste determinismo é \emph{Ribot}, seguido por varios physiologistas. --- Devemos advertir que o determinismo \emph{phiysiologico}, de que nos occupamos, é um systema \emph{phülosophico}, e por isso diverso do determinismo \emph{pliysiologico}, de que falla a physiologia e que consiste no modo uniforme e necessario, por que operam as forças do nosso organismo. A'liberdade oppõe-se o primeiro, não o segundo. Este pode e deve coexistir com a liberdade; porque, como diz \emph{Cl}. \emph{Bernard}, á liberdade refere-se á \emph{direcção} do phenomeno, que se diz \emph{livre}, e o determinismo, emquanto é objecto da \emph{physiologia}, refere-se á \emph{execução} do proprio phenomeno. --- O determinismo \emph{physiologico}, emquanto é um systema \emph{philosophico}, subdivide-se em varias especies, conforme a natureza ou qualidade do antecedente, que determinaria fatalmente todos os nossos actos. As principaes são as que sustentam que o antecedente dos nossos actos é --- ou uma sensação \emph{involuntaria}, e por isso os nossos actos não passam de movimentos \emph{reflexos}, --- ou uma modalidade \emph{organica}, e por isso os nossos actos são apenas o producto \emph{necessario}, a resultante \emph{fatal} d'essa modalidade.}.

\emph{d}) \emph{Determinismo theologico}. --- O determinismo \emph{theologico} ensina que os antecedentes dos nossos actos são os actos da presciencia ou da providencia de \textsc{Deus}; de modo que \emph{todos} os actos da nossa vontade são necessaria e fatalmente fixados por um ou outro desses attributos divinos \footnote{O determinismo \emph{theologico} não encontra geralmente defensores nos modernos escriptores de materias scientificas, porque estes prescindem de Deus, se não o negam; mas foi e é seguido --- pelos sequazes de \emph{Mahomet}, que reconhecem a vontade de \emph{Allahh} como a causa unica de todos os actos humanos, --- e por alguns theistas, que affirmam a incompatibilidade da liberdade humana com a presciencia e providencia de Deus.\par Advertimos que \emph{determinismo} se emprega, ás vezes, como synonimo de \emph{fatalismo}. Mas, rigorosamente, \emph{determinismo} e \emph{fatalismo} são coisas bem distinctas; porque o primeiro designa uma necessidade \emph{exterior}, ao passo que o segundo indica uma necessidade \emph{interior}. --- Por isso, se ás primeiras tres especies de \emph{determinismo} convém o nome que lhe demos, á ultima especie, isto é, ao determinismo \emph{theologico} convém o nome de \emph{fatalismo}.}.

185. --- O determinismo, em geral, é falso.

a) Todo o systema, que contradiz os dados da consciencia moral e psyehologiea, e os principios da razão, é falso. Ora o \emph{determinismo}, pelo facto de negar o livre arbitrio ao homem, contradiz os dados da consciencia moral e psyehologiea, e os principios da razão. Logo o \emph{determinismo} é falso.

b) O \emph{determinismo} suppõe que, se o homem fosse dotado de livre arbitrio, haveria um desarranjo na ordem cosmiea, a qual consiste no nexo fatal e necessario entre os phenomenos naturaes e os seus antecedentes. Ora esta supposição é falsa; porque não se podem applicar ao espirito as leis, que regulam a materia, nem se verificam desarranjos na ordem cosmica pelo facto da vontade humana ser livre. Logo o \emph{determinismo}, em geral, é falso \footnote{São muitos e varios os argumentos, que os sequazes do \emph{determinismo} adduzem para combater o \emph{livre arbitrio}. Uns são \emph{communs} a todas as classes ou especies mencionadas; outros são \emph{proprios} de cada uma. Os principaes argumentos \emph{communs} são dois: um pretende provar que o livro ar- -bitrio \emph{não pode existir}, o outro que \emph{realmente não existe}.\par O \emph{primeiro} argumento foi exposto no texto, mas podemos desenvolvel-o mais. Dizem os deterministas que o \emph{livre arbitrio não pode existir}. E porque não pode existir? Porque, respondem, todos os reinos da natureza estão sujeitos a leis constantes, necessarias, fataes; todos os phenomenos naturaes dependem dos phenomenos precedentes; e por isso não ha nem pode haver logar para seres dotados de livre arbitrio, porque, se o houvesse, se podessem existir causas capazes de produzir effeitos variaveis e não ligados necessariamente aos antecedentes, os phenomenos seriam sempre sujeitos a desvios anormaes sob o influxo do livre arbitrio, seria destruida a ordem immutavel de natureza, e destruida tambem a sciencia, que não\pgn{207}poderia alcançar a necessaria certeza no conhecimento e na predição dos factos. Tal é o raciocinio dos deterministas. --- Mas será racional a sua inducção? Não. Por quanto, se elles interrogaram os seres corporeos, não interrogaram a propria alma, e por isso a inducção é incompleta. Se tivessem investigado a \emph{natureza} da alma humana, teriam concluido que esta é \emph{espiritual}, \emph{e} que, sendo espiritual, não pode estar sujeita ás leis reguladoras dos seres materiaes. Na verdade, as leis, que regulam os seres, não são coisas \emph{abstractas}, mas \emph{concretas}; são os proprios seres, emquanto operam de um modo e não de outro; e por isso as leis são diversas, conforme a diversa natureza dos seres. Se, pois, existem seres espirituaes, substancialmente diversos dos seres corporeos, as leis, que regulam estes, não se podem reduzir ás leis, que regulam aquelles, por isso mesmo que a materia não se pode identificar com o espirito: --- Quando se diz que tudo esta \emph{encadeado} no mundo, não se entende que todo o ser deva estar sujeito a uma: unica lei --- \& da \emph{necessidade}, mas entende-se que todo o ser é ordenado, e que os dois mundos --- o dos corpos e o dos espiritos --- formam um unico plano, e se reunem harmonicamente numa synthese grandiosa. ---- Mas a liberdade, dizem, é uma desordem no mundo. E porque será ella uma desordem ? Acaso o movimento dos planetas cessa, unicamente porque eu assim o quero; --- ou os tres angulos de um triangulo deixam de ser eguaes a dois rectos, só porque eu sou dotado de livre arbitrio? --- Mas as previsóes da sciencia são incompativeis com o livre arbitrio. E porque serão ellas incompativeis? Se uma lei physica não se pode dizer e não se diz destruida, quando um agente imprevisto e necessario suspende o effeito della, poderá dizer-se destruida, quando esse agente é livre e só porque é livre? Não, por certo. O resultado, isto é, a suspensão do effeito, é o mesmo em ambos os casos, e a \textbar{} liberdade ou a necessidade do acto do agente é indifferente para a proprio \textbar{} effeito. --- Mas, se a vontade é livre, dar-se-ha uma anomalia, qual é a exi- \textbar{} stencia de uma força creada não sujeita a uma lei propria. Respondemos \textbar{}\par que tambem a nossa vontade está sujeita a essa lei, \emph{a} qual consiste em que esta faculdade tenha uma tendencia \emph{necessaria} para o bem \emph{universal}, e uma tendencia \emph{livre} para os bens \emph{particulares}.\par O \emph{segundo} argumento, com que os \emph{deterministas} pretendem demonstrar h a não existencia do livre arbitrio, é tirado da \emph{estatistica}, e é o seguinte: A \emph{estatistica} mostra que um grande numero de factos, que se consideram como livres, se reproduzem com espantosa fixidez e regularidade; em todos os annos e em todos os paizes dá-se \emph{o} mesmo numero de desgraças, de crimes, de suicídios, de casamentos, etc. Ora isto indica que esses phenomenos estão \textbar{} sujeitos a leis fixas e inovitaveis, e que por isso a vontade humana não é livre. --- Em poucas palavras: Os phenomenos fixos são necessarios; ora os phenomenos sociaes são fixos; logo são necessarios (Buckle, \emph{Histoire de la civilisation en Angleterre}, trad. Baillot, t. I, pag. 36).\par Este argumento dos \emph{deterministas}, se parece verdadeiro quanto á \emph{forma}, é falso quanto á \emph{materia}, como vamos demonstrar. Com effeito:\par a) A premissa. \emph{maior}, á qual diz que os \emph{p/tenomenos fixos são necessarios},\pgn{208}é falsa. E' certo que todos os phenomenos necessarios são fixos; pois que, como diz Newton (\emph{Principes mathématiques de la Philosophie naturelle}, Scolie gónéral) a necessidade metaphysica, que sempre e em toda a parte é a mesma, não póde produzir effeitos diversos; mas é falso que todos os phenomenos fixos sejam necessarios, pois que os phenomenos naturaes são fixos, e todavia, como demonstrámos na Cosmologia, não são necessarios.\par \emph{b}) A premissa \emph{menor} --- \emph{os phenomenos sociaes são ficos} --- é falsa. Os verdadeiros dados da \emph{estatistica} não assignalam a fixidez dos phenomenos sociaes, mas sim uma grande variação. Esta variação dá-se --- não só relativamente aos diversos paizes, pois Legoyt (\emph{Le swicide ancien et moderne}) enumera, para cada milhão de habitantes, 261 suicidios na Allemanha e só 17 na Hespanha, --- mas tambem relativamente ao mesmo paiz, assim em Franca, no anno 1827 deram-se 1739 casos de suicidio, e no anno 1878 subiram a 6485. Logo os phenomenos sociaes não são fixos.\par \emph{Buchle} (ob. cit.) objecta: Se. os phenomenos mudam, é porque o meio social muda; mas nas mesmas circumstancias os phenomenos são os mesmos.\par Resposta. E' falso que a mudança dos phenomenos deva exclusivamente attribuir-se á mudança do meio social. Em Franca, nestes ultimos annos, não houve propriamente mudança de meio, e todavia houve mudança nos phenomenos, pois o numero dos suicidios nesse paiz, foi --- em 1875 de 5472, --- em 1876 de 5804, --- em 1877 de 59292, --- em 1818 de 6485. --- Mas admittamos que os phenomenos mudem com a mudança das circumstancias e do meio, que se seguirá d'ahi? Seguir-se-ha por ventura que os phenomenos são necessarios? Não. A historia das nações manifesta que a mudança das cirecumstanoias \emph{depende principalmente do livre arbitrio dos homens}, e que por isso os phenomenos sociaes são livres.\par c) A \emph{conclusão}, pois, que diz que \emph{os phenomenos sociaes são necessarios}, é falsa; porque uma proposição, que se apoia em principios falsos, não póde deixar de ser falsa. --- Mas, embora os phenomenos sociaes fossem \emph{necessarios}, seguir-se-hia que o homem não é dotado de \emph{livre arbitrio}? Nào, certamente. Se o crime é determinado, não o é o criminoso. Entre muitos homens de bem, ha um criminoso. E porque é que um homem commette um crime, e os outros fioam probos e honestos? A unica resposta racional é que uns se determinaram a fazer o bem, e outros a fazer o mal. Ora ninguem poderia determinar-se para o bem nem para o mal, se não tivesse o \emph{livre arbitrio}.\par Como se vé, a difficuldade confunde a falta de \emph{variabilidade} com a falta de \emph{liberdade}. Mas os dois conceitos são muito diversos. Deus é \emph{immutaxvel}, e, comtudo, é eminentemente \emph{livre}. A liberdade, longe de ser um principio de perturbação e de variabilidade, é um principio de ordem e de constancia ; porque, s»b a luz da razão, nos faz conitrahir disposições e habitos, que se tornam uma segunda natureza.}. .

%%>>> t136

\pg{206} \pg{207} \tese{136}{O determinismo psychologico é falso.}

a) Os defensores d'esta especie de \emph{determinismo} suppõem que os \emph{motivos} exercem sobre a nossa vontade uma pressão \emph{material}, impellindo-a \emph{necessariamente} para um bem certo e \pg{208}  individuo. Ora esta supposição é falsa; porque o \emph{motivo}, por ser o conhecimento do bem, contido num objecto, é uma coisa \emph{espiritual}, como espiritual é a nossa vontade, e por isso não se pode dar essa supposta pressão \emph{material}. Logo o determinismo \emph{psychologico} é falso.

b) O determinismo \emph{psychologico} affirma que á nossa vontade não pode deixar de obedecer ao \emph{motivo mais forte}, ou de \pg{209}  seguir o \emph{bem melhor}. Ora esta affirmação é falsa; porque, como provámos, a nossa vontade, se n\&o pode deixar de seguir a attracção do bem \emph{universal}, que enche toda a sua capacidade, é completamente indifferente e livre com relação aos bens \emph{particulares}, seja qual fór o seu grau de intensidade ou perfeição. Logo o determinismo \emph{psychologico} é falso \footnote{E' a famosa theoria dos \emph{motivos-forças}. Teve por auctor o celebre \emph{Leibnitz}, que raciocinava assim: « O acto livre ou tem um motivo, ou não o tem. Se o não tem, é cego, irracional, sem razão sufficiente. Se o tem, deve seguir o mais forte, aliás seria um acto sem razão sufficiente ». Como, porém, não quiz este philosopho negar o livre arbitrio, por isso inventou uma distincção entre a \emph{determinação infallivel} e a \emph{necessidade absoluta}, e disse que a vontade, solicitada por varios motivos, obedecerá \emph{infallivelmente} ao mais forte, ficando comtudo livre, em todos os outros casos, nos quaes não se conhece o \emph{bem melhor}. --- Os sequazes de \emph{Leibnitz} foram mais radicaes, mas mais logicos; defenderam a theoria do \emph{motivo mais forte} até á ultima conclusão, que é a negação do livre arbitrio. \emph{Maudsley} diz: « Existe um laco morale uniforme entre os motivos e a operação. A vontade obedece sempre ao motivo mais forte. O livre arbitrio, pois, é um absurdo sem sentido ».\par Para que a falsidade do determinismo \emph{psychologico} seja mais evidente, procuraremos explicar com exactidão a \emph{natureza} do motivo, e a \emph{parte}, que elle desempenha no acto livre da nossa vontade.\par A vontade tem por objecto o bem. O bem é, ao mesmo tempo, \emph{absoluto} e \emph{relativo}. E' \emph{absoluto}, emquanto deve possuir alguma realidade ou perfeição; 6 \emph{relativo}, emquanto se diz \emph{bem} o que \emph{convém} \& nós ou aos outros. O \emph{ser} ou a \emph{perfeição} é o elemento \emph{material} ou \emph{fundamental} do bem; á \emph{conveniencia} é o seu elemento \emph{formal}. Esta \emph{conveniencia} consiste na aptidão, que o ser possue, de aperfeiçoar os outros, e é o \emph{motivo} ou a razão \emph{sufficiente}, pela qual a vontade quer o objecto. Assim como não se pode querer um objecto sem que se descubra nelle uma razão de conveniencia, assim tambem não se pode querersem um motivo sufficiente. Preferir um objecto a outro sem motivo, não seria uma escolha, não seria um acto de uma faculdade intellectual,mas um acto cego, um simples movimento mechanico. --- Ora o motivo, manifestando-se á intelligencia, \emph{excita} a vontade, a qual está sempre em \emph{acto} com relação ao bem \emph{universal}, e produz nella um amor \emph{espontaneo} pelo objecto, que se lhe apresentou como \emph{conveniente}. --- Mas é claro que esta acção do motivo não é \emph{necessaria}, de modo que arraste \emph{érresistivelmente} a vontade para o acto. Por quanto, se assim fosse, isto é, se o motivo exercesse sobre a vontade uma força irresistivel, não teriamos o poder de examinar os motivos, de analysal-os e comparal-os, nem de preferir um ao outro. Todavia temos este poder; comparamos os motivos entre si; temos consciencia de que poderiamos rejeital-os todos, e todavia escolhemos um, não attendendo aos outros. Ora isto demonstra que não é o motivo que torna efficaz a vontade, mas a vontade torna efficaz o motivo. Logo o motivo é a rasdo \emph{sufficiente}, á causa \emph{final} do \emph{acto livre}, mas não é a sua causa \emph{efficiente} ou \emph{determinante}. A causa \emph{determinante} é unicamente a vontade. --- Nem se diga que o motivo, pelo facto de excitar o amor \emph{espontaneo} na vontade, determina neces-. \emph{sariamente} esta faculdade ao acto. O amor espontaneo, que os motivos des-\par 14 Vor. II.\pgn{210}pertam na vontade, é uma \emph{condição}, sem a qual a vontade não tomaria nenhuma deliberação, --- mas não é a \emph{eleição}; porque a vontade, recebida a impressão, fica sempre na sua indifferença \emph{activa}, \emph{e}, depois do exame dos - varios motivos, escolhe aquelle, ao qual quer ceder, e assim o motivo escolhido recebe da vontade o caracter \emph{pratico} de \emph{predominante}.\par Estes principios, deduzidos da analyse dos nossos actos livres, não só vigoram quando se trata de motivos de egual peso, mas tambem quando um motivo é \emph{mais forte} que os outros. Com effeito, a força maior ou menor de um motivo é mais subjectiva, do que objectiva. E' sempre a vontade que inclina a intelligencia a considerar um objecto, e não um outro, e a considerar, no mesmo objecto, as qualidades boas e favoraveis e não as más e as contrarias. D'este modo a inclinação, que no principio era mais fraca, torna-se mais forte, e o objecto que, á primeira vista, nos parecia pouco amavel, depois parece-nos o melhor, --- o melhor não tanto pelo que tinha de per si, quanto pelo que lhe accrescentou a vontade. Por isso, quando cedemos ao motivo mais forte, cedemos a nós mesmos, obedecemos á nossa livre escolha. --- Mas supponhamos que um motivo é mais forte que os outros \emph{objectivamente}, por si mesmo, independentemente da nossa vontade. Que se seguirá d'ahi? Seguir-se-ha por ventura que elle ha-de exercer uma acção irresistivel sobre a nossa vontade? Não, por certo. Attesta-nos a consciencia, que se, ás vezes, nos determinamos pelo motivo mais forte, fazemol-o livremente e que, se quizessemos, podiamos abandonar o motivo mais forte e seguir o mais fraco. Na verdade, quem poderá sustentar que a vontade escolhe sempre o bem, que a intelligencia apontou como o melhor? Quantas vezes, apezar de todos os protestos da intelligencia, a vontade prefere um bem mesquinho, mas presente, a um bem superior, mas futuro! Quantas pessoas podem repetir com Ovidio: « \emph{Video meliora}, \emph{proboque}, \emph{deteriora sequor} »?! --- E qual será a razão, por que um motivo, embora mais forte que outro, não determina necessariamente a vontade? A razão é evidente. Todo o motivo, que se não refere ao bem \emph{universal} e \emph{absoluto}, refere-se a um bem \emph{particular} e \emph{relativo}, Referindo-se a um bem \emph{relativo} e \emph{particular}, nenhum motivo, embora seja mais forte que outro, exerce uma pressão necessaria e irresistivel sobre a nossa vontade, mas deixa-a na indifferenca; porque a vontade, como provámos, só pode ser attrahida irresistivelmente por um motivo, que se refira ao bem \emph{universal} e \emph{absoluto}.\par O erro d'estes deterministas consiste em quo consideram os motivos como outras tantas causas \emph{e/ficientes} ou \emph{determinantes}, e a vontade como um simples \emph{movel}, sujeito a varias forças contrarias. E, assim como, dizem elles, no conflicto de duas forças ou causas efficientes, --- se ellas são eguaes ou oppostas, o effeito é nullo, --- e, se são deseguaes, o effeito segue necessariamente a direcção da mais forte; assim tambem a vontade, sujeita ao\pgn{211}influxo de varios motivos, não poderá escolher, se os motivos são absolutamente eguaes (e um faminto, collocado entre dois pães de todo eguaes, morrerá de fome), ou escolherá necessariamente o mais forte, se os motivos são deseguaes. --- Mas, se é absurdo identificar os motivos com as causas \emph{efficientes}, quando elles não passam de \emph{razões sufficientes} ou \emph{causas finaes}, não é menos absurdo comparar a vontade, que delibera, com um \emph{novel}, sujeito ao contraste de varias forças. Por quanto, no contraste das forças, estas são a causa \emph{efficiente} do acto e o movel é o sujeito \emph{passivo}, em que o proprio acto se realisa; ao passo que, na deliberação da vontade, os motivos são apenas causas \emph{finaes} ou razões \emph{sufficientes}, e a propria vontade é, ao mesmo tempo, \emph{sujeito} e \emph{principio elicitivo} do acto. Além d'isso, no conflicto das forças, o movel encontra-se num estado de indifferenca \emph{passiva}, e por isso não pode actuar nas forças, não as fixa, não as oppõe umas ás outras, mas limita-se a experimentar os effeitos d'ellas, se d'isso é capaz; ao passo que a vontade encontra-se na indifferenca \emph{activa}, e, portanto, se não é senhora da primeira idea, pode todavia, depois de a ter recebido, fixal-a, examinal-a, comparal-a com as outras, e depois tomar uma resolução ou fazer a eleição. --- D'ahi se segue que --- quando os motivos ou os bens são \emph{eguaes}, a vontade escolherá um, porque é bem, e deixará o outro, porque não é um bem necessario; --- mas, quando os motivos ou os bens são \emph{deseguaes}, a vontade pode ou escolher o bem maior, porque é bem e porque é maior; ou preferir-]he o bem menor, porque é um bem sufficiente, embora inferior. E), pois, falso que a preferencia do bem menor ao maior seja \emph{sempre} um acto irracional, ou destituido de razão sufficiente. Dizemos \emph{sempre}; porque essa preferencia é irracional, quando um certo bem é um meio \emph{necessario} para se conseguir mais efficazmente um fim determinado, e quando a vontade fez o proposito de conseguir um fim pelo \emph{melhor modo possivel}; visto que, em tal caso, a escolha de um bem \emph{inferior} seria uma contradicção manifesta.}.

%%>>> t137

\pg{210} \tese{137}{O determinismo mechanico é falso.}

a) Todo o systema, que se funda num supposto falso, é falso. Ora o determinismo \emph{mechanico}, reduzindo a energia da vontade humana a uma força \emph{material}, sujeita, como qualquer outra força do mesmo genero, ás leis da mechanica, funda-se num supposto falso; porque a energia da nossa vontade é inteiramente \emph{espiritual}. Logo o determinismo \emph{mechanico} é falso.

b) Se a nossa vontade estivesse sujeita ás leis da mechaniea, não poderia operar de um modo contrario a essas leis. Ora a nossa vontade opera de um modo contrario ás leis da mechaniea. Com effeito, segundo as leis da mechaniea, uma substancia em movimento deve tocar todos os pontos do espaço, que percorre; mas a nossa vontade passa de repente de uma coisa a outra contraria, sem tocar os pontos intermedios. Logo o determinismo \emph{mecAhanico} é falso \footnote{O argumento contra o \emph{livre arbitrio}, que é apresentado pelos \emph{deterministas mechanicos}, funda-se na lei da \emph{conservação da energia}, e é o seguinte: Todo o movimento do nosso corpo é produzido por um outro movimento material ou mechanico, e não pela alma. Na verdade, se fosse produzido pela alma, admittir-se-hia uma nova força, causada pela propria alma e distincta da força material; o que é impossivel, porque a \emph{força} ou \emph{energia} é invariavel, não é capaz de augmento nem de diminuição. Logo \emph{todos} os actos do homem, mesmo os chamados \emph{livres}, dependem necessariamente de um previo movimento mechanico, e por isso o homem não é dotado de livre arbitrio.\par Este argumento --- \emph{a}) é falso na premissa, que estabelece, --- 5) é inconcludente na prova, que adduz, --- \emph{c}) é illogico na consequencia, que tira.\par a) \emph{E falso na premissa}, \emph{que estabelece}. A premissa é esta: « Todo o movimento\pgn{212} material do nosso corpo é produzido por outro movimento material, e não pela alma ». --- Esta asserção é falsa.\par Na verdade, nenhum movimento material é produzido por outro movimento material. Por quanto, como observa \emph{Zr}», fundador da \emph{thermodynamica}, todo o moyimento material é precedido de um momento de quietação; porque todo o movimento ou é \emph{vibratorio}, ou reduz-se ao \emph{vibratorio}, e as vibrações estão comprehendidas entre dois momentos, em que o movimento é nullo. Ora, no momento de quietação, o movimento antecedente já não existe, e por isso não póde determinar o movimento seguinte. O movimento material pois, não póde ser produzido por outro movimento material.\par Todavia, o movimento, por ser um effeito real exige uma causa real. Esta causa será a materia? Não; a materia é por si inerte, indifferente tanto para a quietacáo como para o moyimento, e ninguem dá o que não possue. Deve ser, pois, uma força distincta da materia. Esta força é a alma.\par A alma, por estar intimamente unida ao systema nervoso, apodera-se das forças brutas, e dirige-as para a plena realisação da sua ordem vital. A sua acção desenvolve-se toda no interior do nosso organismo, mas nós vemol-a na vasta esphera de seus admiraveis effeitos. Não devemos, comtudo, Julgar que a alma se ponha em antagonismo com as forças brutas. Não: ella revela-se sempre como potencia jerarchicamente superior e reguladora d'essas energias. Por isso, no emprego das forças ou energias organicas, a acção da alma acha-se restricta entre dois limites: um determinado pelas proprias forças disponiveis, que só podem produzir certos effeitos, --- outro determinado pelas leis mechanicas, physicas e chimicas, que a alma pode utilisar em beneficio proprio, mas não pode alterar. --- Nem se opponha o principio dos mechanicos: uma força material não pode passar do estado potencial para o actual senão pela acção de outra força material. Esse principio refere-se unicamente ao dominio da mechanica, e por isso, se pode applicar-se ás mutuas operações \emph{transeuntes} dos seres materiaes, não se pode applicar ás operações \emph{immanentes} dos seres vivos. --- E\textasciicircum{}, portanto, falsa a premissa dos deterministas \emph{mechanicos}, que diz: Todo o movimento do nosso corpo é produzido por outro movimento material ou mechanico, e não pela alma.\par b) \emph{E}' \emph{inconcludente na prova}. Dizem estes deterministas que, se o movimento do nosso corpo fosse produzido pela alma espiritual, deveria admittir-se uma nova força, distincta da força material; o que é impossivel, porque contrario á \emph{lei da conservação da energia}, segundo a qual a força é invariavel, incapaz de augmento e de diminuição. Ora esta prova é inconcludente. Na verdade, a \emph{lei da conservação da energia} é expressa por \emph{Helmholts} na seguinte formula: « A somma das forças vivas e das energias potenciaes é constante n. (\emph{Mêm}: sur \emph{la conserv}. \emph{de la force}). O que é que ensina essa lei, segundo a opinião dos naturalistas? Essa lei ensina que em todas as transformações, que se realisam no mundo \emph{material} e pelos agentes \emph{materiaes}, fica constante e inalteravel a quantidade da energia actual e potencial, por meio da qual se desenrola o processo cosmico. Por isso esta lei só se pode applicar aos\pgn{213}systemas \emph{materiaes}, onde só actuam agentes materiaes. E', pois, um erro querer excluir do mundo um agente \emph{immaterial}, com o pretexto de que a sua intervenção vem augmentar a energia cosmica, destruindo assim a constancia da força; porque a força de um agente \emph{imunaterial}, como é a nossa alma, é \emph{immaterial}, e a força dos agentes corporeos é \emph{material}, e por isso a existencia da força immaterial não importa augmento nem diminuição na força material, e fica sempre inviolada a \emph{lei} da \emph{conservação da energia}, a qual, como dissemos, se refere unicamente aos systemas materiaes, onde actüam agentes \emph{materiaes}. --- Este pensamento merece ser desenvolvido. Supponhamos que a \emph{lei da conservação} se possa e deva applicar a \emph{todos} os agentes, \emph{materiaes} e \emph{immateriaes}, o que se seguirá d'ahi? Seguir-se-ha por ventura que se deve negar á nossa alma o poder de determinar ao movimento o nosso corpo? Não, certamente. A lei da conservação da força diz respeito á \emph{quantidade}, e não á \emph{qualidade}. Todos concordam em que a \emph{mesma} força pode encontrar-se no estado potencial ou actual, tomar uma certa direcção ou uma outra differente, passar por varios e diversos estados, sem que por isso augmente a sua \emph{quantidade}. Ora a actividade \emph{immanente} dos seres vivos, e por isso tambem da nossa alma, não é \emph{quantitativa}, isto é, não acocrescenta nada ao deposito das energias physicas, mas é \emph{qualitativa}, isto é, actua nas energias potenciaes e actuaes, transformando-as umas nas outras, e, debaixo de formas multiplas e variadas, restitue fielmente a mesma \emph{quantidade}, que recebeu. --- Logo a prova dos deterministas é, pelo menos, \emph{inconcludente}.\par €) \emph{E}' \emph{illogico na consequencia}. Se a \emph{lei da conservação da energia} não se refere á \emph{quantidade}, mas á \emph{qualidade}, é evidente que ella, ainda que se applique ao mundo vivo, não pode destruir o livre arbitrio do homem. \emph{Liberdade} e \emph{necessidade} são dois modos, por que a nossa vontade opéra, são duas \emph{qualidades} differentes, e não constituem uma nova \emph{quantidade}, que venha perturbar as equações dos mechanicos. Que o meu braço levante \emph{livre} ou \emph{fatalmente} um peso de cinco kilos, a \emph{quantidade} de energia é sempre a mesma, a differenca entre o modo livre e necessario é inteiramente psychica e moral; --- numa palavra, só a \emph{qualidade} é diversa. Logo a lei \emph{da conservação da energia} nada tem que ver com o livre arbitrio do homem, e por isso é, pelo menos, illogico fundar sobre essa lei a negação d'esta excelsa prerogativa da nossa vontade. --- Concluimos com as palavras de \emph{Raoul Pictet}: « A liberdade opéra sem prejuizo de nenhuma lei ensinada na physica experimental ». (\emph{Etude critique du materialismo}, p. 449).}.

%%>>> t138

\pg{211} \pg{212} \tese{138}{O determinismo physiologico é falso.}

a) E' falso o determinismo \emph{physiologico}, emquanto reduz todos os actos da nossa vontade a movimentos \emph{reflexos}. Com effeito, movimento ou acto \emph{reflero} é como que uma resposta ou reacção a uma excitação peripherica ou exterior, sem intervenção \pg{213}  da vontade nem do appetite sensitivo. Ora o acto \emph{livre} deriva da vontade e presuppõe necessariamente o conhecimento da intelligencia. Logo é absurdo reduzir os actos da vontade a movimentos \emph{reflexos} \footnote{O acto \emph{reflexo}, sendo uma reacção involuntaria e muitas vezes in- consciente, provocada por uma excitação peripherica ou \emph{exterior}, distingue-se --- tanto do movimento \emph{espontaneo}, que é provocado pelo appetite sensitivo e, no homem, em ultima analyse, pela vontade, --- como do movimento \emph{automatico}, que é determinado por um estimulo \emph{interno}, ainda que independentemente do appetite sensitivo e da vontade. Ora, reduzir os actos \emph{voluntarios}, e sobretudo os \emph{livres}, a movimentos \emph{reflexos}, é confundir e identificar a luz com as trevas, porque o acto \emph{livre} procede da razão, ao passo que o movimento \emph{reflexo} exclue todo e qualquer conhecimento,}.

b) E' falso o determinismo \emph{pAjysiologico}, emquanto ensina que todos os actos da nossa vontade são o resultado fatal de uma modalidade \emph{organica}, transmittida pela hereditariedade. Por quanto, pela hereditariedade, não se transmitte ao filho o estado pathologieo dos paes, nem as suas qualidades, boas ou más, mas o mais que se transmitte é uma certa predisposição para esse estado ou para essas qualidades. Ora esta predisposição não só não é fatal, nem destroe o livre arbitrio, mas pode ser e muitas vezes é completamente subjugada e destruida. Logo o determinismo \emph{physiologico} é falso \footnote{Segundo esta forma de determinismo, o homem nasce bom ou mau, como nasce branco ou preto. D'onde a celebre theoria do \emph{criminoso nato}, de que mais tarde nos occuparemos. --- Para não haver equivocos, repetimos uma observação. Não se pode negar que, no homem, ha uma grande dependencia entre o physico e o moral. A razão é clara. O homem, embora composto de duas substancias diversas, tem comtudo uma só natureza e é uma só pessoa. Todas as nossas faculdades brotam da mesma unica essencia da alma, e por isso não só o exercicio de uma depende do exercicio da outra, mas tambem as suas energias desenvolvem-se na razão inversa, isto é, quanto mais intensa é a energia de uma, tanto mais fraca é a energia da outra, ficando esta por vezes completamente neutralisada. Ora a nossa vontade, embora seja, pelo seu livre \emph{arbitrio}, a rainha de todas as potencias, comtudo, no seu exercicio, depende das outras, e está sujeita ás influencias d'ellas. --- Mas, se ó racional admittir uma dependencia entre o physico e o moral, é illegitimo identificar um com o outro; entre os dois ha uma distancia immensa. --- Como tambem, se é certo que as paixões, as condições physicas ou moraes, as disposições naturaes e hereditarias, exercem uma influencia notavel no habitual exercicio do livre arbitrio, é comtudo falso que essa influencia seja \emph{ordinariamente} necessaria, fatal, irresistivel. Dizemos \emph{ordinariamente}, porque, dependendo do exercicio da intelligencia o exercicio da vontade, se a intelligencia soffre um desarranjo nas suas funcções, este desarranjo terá uma repercussão na vontade, e esta não poderá produzir um acto plenamente deliberado. Voltaremos a este assumpto com maior desenvolvimento,}.

%%>>> t139

\pg{214} \tese{139}{0 determinismo theologico é falso.}

a) E' falso o determinismo \emph{theologico}, emquanto ensina que todos os nossos actos, pelo facto de terem sido \emph{previstos} por \textsc{Deus}, devem infallivelmente realisar-se, e por isso são \emph{necessarios}. Na verdade, \textsc{Deus} não \emph{prevê}, mas vê todos os nossos actos, porque tudo é \emph{presente} aos seus olhos. Ora a \emph{visão} de um facto, embora infallivel, não altera a \emph{natureza} d'esse facto, mas deixa-o ficar como é; de modo que se o facto é necessario, permanece necessario, e se é livre, permanece livre. Logo é falso que todos os actos da nossa vontade, pelo facto de terem sido \emph{vistos} por \textsc{Deus}, sejam necessarios,

kb) E' falso o determinismo \emph{theologico}, emquanto sustenta que todos os actos da nossa vontade são necessarios, porque todos exigem a moção divina, á qual não pode a creatura resistir. Com effeito, a nossa vontade é movida por \textsc{Deus} não como um instrumento qualquer, mas como um instrumento racional e livre, que é movido e, ao mesmo tempo, se move a si mesmo. Ora, aecomodando-se a Causa primeira á indole e ás condições das causas segundas, quando Ella move a vontade humana, não só não lhe tira o livre arbitrio, mas, antes, produz e conserva nella tão admiravel prerogativa. Logo é falso que todos os actos da nossa vontade sejam necessarios, pelo facto de esta potencia ser movida por \textsc{Deus} \footnote{Diz S. Thomaz: « Deus é a Causa primeira, que move as causas naturaes e as voluntarias. E assim como, movendo as causas naturaes, não obsta a que os actos d'ellas sejam naturaes, assim tambem, movendo as causas voluntarias, não só não obsta a que os actos d'ellas sejam voluntarios, mas até faz que sejam taes; porque Deus opéra nas creaturas em conformidade com as propriedades, de que ellas são dotadas ». (\emph{Sum}. \emph{Theol}, p. I, q. 88, a. 1 ad 3). --- A difficuldade de conciliar a sciencia e a providencia de Deus com o livre arbitrio do homem nasce da limitação e da ignorancia da nossa intelligencia, a qual não conhece bem o modo por que Deus opera. --- Na \emph{theodicea} voltaremos a este assumpto, e procuraremos dar uma solução mais satisfactoria a esta difficuldade.}.

%%>>> t140

\pg{215} \tese{140}{Consequencias da negação da liberdade humana.}
A negação da \emph{liberdade humana}, além de ser um erro contrario aos principios da razão e á consciencia individual e social, dá origem ás mais perniciosas consequencias, pois que tende a destruição da \emph{moralidade}, da \emph{sociedade} e da \emph{religião}.

a) Destroe a \emph{moralidade}; pois que, se todas as acções fossem necessarias e inevitaveis, desappareceria o fundamento da differenca entre o bem e o mal moral, e assim todas as acções seriam boas e honestas, e nenhum erime seria imputavel.

b) Destroe a \emph{sociedade}; visto que, sem o \emph{livre arbitrio}, não teriam razão de existir as leis e os magistrados, os direitos e os deveres reciprocos, os premios e os castigos, e, sem estes elementos, não haveria sociedade.

c) Destroe a \emph{religião}; porque os homens, se não fossem dotados de \emph{liberdade}, não poderiam ser obrigados a cumprir certos deveres para com \textsc{Deus}, nem a dirigir os pensamentos, as palavras e as obras para aquella Bondade infinita, que os creou e com tão desvelada providencia os conserva \footnote{A. vista das funestas consequencias, que derivam da negação do livre arbitrio, tem embaracado os deterministas. Mas estes nem por isso rejeitaram\pgn{216}os seus principios. --- \emph{Diderot}, \emph{Lamettrie}, \emph{Helvetius}, no seculo 18, reconheceram o rigor logico d'aquellas consequencias, admittiram-nas cynicamente, e 8Ó accrescentaram que, na vida pratica, os homens devem comportar-se diversamente, isto é, como se fossem dotados de livre arbitrio. --- Os deterministas dos ultimos tempos mudaram de tactica, negaram essas consequencias, e disseram que a moral não só póde mas até deve prescindir do livre arbitrio; porque este, longe de ser o fundamento d'aquella, é, pelo contrario, a sua destruição. Por isso procuraram demonstrar que as noções de \emph{merito} e de \emph{demerito}, de \emph{virtude} e de \emph{vicio}, de \emph{recompensa} e de \emph{castigo}, não só não ficam destruidas no seu systema, mas d'elle recebem uma força efficacissima, --- mais efficaz do que recebiam da velha theoria do livre arbitrio. A affirmação era demasiado atrevida. Mas o que haviam de fazer? Como já indicámos, alteraram o sentido de todas as noções moraes, substituindo á legitima significação, reconhecida pelo genero humano, uma significação arbitraria, em harmonia com o seu systema, e assim, e só assim, poderam proclamar que a sciencia e a vida moral tinham progredido immensamente, depois do abandono da theoria do livre arbitrio. --- A estas falsificações nos referiremos ainda, e com o devido desenvolvimento, na parte \emph{moral} d'este curso.}.

%%>>> art-20

\section*{Artigo IV.\\ Relações entre a intelligencia e a vontade.}
\addcontentsline{toc}{section}{Artigo IV. --- Relações entre a intelligencia e a vontade.}

%%>>> t141

\pg{216} \tese{141}{A intelligencia move a vontade.}
Uma faculdade move outra, quando aquella apresenta o objecto, para o qual esta tende. Ora a \emph{intelligencia} apresenta o objecto para o qual a \emph{vontade} tende; pois a vontade não poderia tender para o bem, que não fosse já conhecido pela intelligencia, visto que ninguem deseja o que ignora. Logo a \emph{intelligencia} move a vontade.\footnote{Cf. \emph{Sum}. \emph{Theol}., p. T, q. 82, a. 4; I-II, q. 9, a. 1, a. 8 ad 8; QQ. DD. \emph{De verit}., q. 22, a. 12. --- Os philosophos não estão de accordo ácerca do influxo da intelligencia na vontade, mas emittem tres diversas opiniões.\par 1.) Uns, como \emph{ZJoocke}, dizem que a vontade não depende da intelligencia, podendo tender para um objecto, sem previo juizo ácerca da bondade ou malicia d'elle, --- e accrescentam que ella não é boa ou má, pelo facto de tender para um objecto bom ou mau, mas sim que um objecto é bom, só pelo facto de ser o termo da vontade.\par 2.) Outros, como \emph{Scoto}, \emph{Suarez}, sustentam que, embora a nossa vontade não possa operar sem um previo juizo da intelligencia, todavia não é necessario que esse juizo seja \emph{pratico}, mas basta que seja \emph{especulativo}. --- Advertimos que o juizo, com que a intelligencia decide se uma coisa deve ser escolhida ou rejeitada pela vontade, se divide em \emph{especulutivo} e \emph{pratico}. E' \emph{especulativo}, quando affirma que uma coisa é bóa ou má, e por isso merece ser escolhida ou rejeitada. E' \emph{pratico}, quando, depois do exame dos motivos, decide que uma coisa, \emph{nestas determinadas circumstancias}, deve ser escolhida ou rejeitada pela vontade.\par 9.) Outros, como \emph{Santo Thomas}, seguido por. \emph{Leibnita} e muitos outros sabios, ensinam não só que a vontade não póde operar sem um previo juizo\pgn{217}\emph{pratico}, mas tambem que o juizo \emph{pratico}, seguido pela vontade, é o \emph{ultimo} dos que á razão \emph{pratica} emitte. Façamos uma breve critica d'estas opiniões.\par a) A \emph{primeira} opinião é falsa. Com effeito, a nossa vontade, sendo uma faculdade \emph{racional}, deve necessariamente operar por algum motivo, apresentado pela intelligencia. Logo é falso que a vontade possa determinar-se sem motivo. --- De mais, é impossivel que uma faculdade forme ou produza o seu proprio objecto; pois o objecto de uma faculdade deve existir antes que ella opere. Ora o objecto da vontade é o bem. Logo o bem não é produzido pela vontade, como afirmam os adversarios, mas existe antes e independentemente da mesma operação. --- Finalmente, esta opinião é \emph{impia} e \emph{perniciosa}. \emph{Impia}; porque a bondade das creaturas só póde derivar da bondade do Creador. \emph{Perniciosa}; porque nenhuma coisa seria em si bóa ou má, ou, melhor, toda a acção, embora má, seria bóa, pelo facto de ser objecto da vontade.\par \emph{b}) \emph{A segunda} opinião tambem é falsa. Por quanto, a vontade não póde determinar-se á operação senão por um juizo, que se refira á \emph{pratica} e que indique o que \emph{nestas circumstancias} se ha de fazer ou deixar de fazer. Ora o juizo meramente \emph{especulativo} não se refere á \emph{pratica}; pois, embora diga que uma coisa é bóa ou má, comtudo não decide sobre o que se ha de fazer ou deixar de fazer \emph{nesta determinada circumstancia}. Logo é falso que a vontade possa determinar-se por um juizo meramente \emph{especulativo}.\par \emph{c}) A \emph{terceira} opinião é verdadeira. Com effeito, a vontade não póde deixar de seguir um juizo \emph{pratico}, \emph{e} o \emph{ultimo} juizo \emph{pratico}, isto é, o juizo que precede immediatamente o assentimento da vontade. --- Não póde deixar de seguir um juizo \emph{pratico}. Por quanto, como demonstrámos contra os defensores da segunda opinião, a vontade não póde determinar-se a operar por um juizo meramente \emph{especulativo}. --- Não póde deixar de seguir o \emph{ultimo} juizo pratico. Com effeito, \emph{a} vontade não póde operar sem motivo. Ora, se operasse contra o \emph{ultimo} juizo da razão \emph{pratica}, operaria sem motivo; pois que o motivo, que determina uma acção, é o que a razão \emph{pratica} dicta por ultimo, --- é o ultimo a mover a vontade, e não pode separar-se da acceitação da vontade.\par Objectam: Se a vontade não podesse deixar de seguir o \emph{ultimo} juizo da razão \emph{pratica}, não seria livre. Ora é livre. Logo póde deixar de seguir o \emph{ultimo} juizo da razão pratica.\par Respondemos, negando a premissa \emph{maior}. Por quanto, o \emph{ultimo} juizo \emph{pratico}, que a vontade não póde deixar de seguir, é um acto, que a mesma vontade \emph{livremente} determinou e regeu, movendo e inclinando a razão a pensar mais numa certa coisa do que noutra, e por isso a produzir este juizo e não outro, e fazendo com que este juizo fosse o \emph{ultimo}. Se a vontade segue o juizo, que ella mesma assim determinou e regeu, não pode deixar de ser livre; porque então, longe de seguir o acto da \emph{razão}, a \emph{vontade} subordina ao acto proprio o acto d'aquella faculdade. Diz S. Thomaz: « Iudicium\pgn{218}de hoc particulari operabili, \emph{ut nunc} (nestas determinadas cireumstancias), numquam potest esse contrarium appetitui » (\emph{De} Ver., q. 24, a. 2). --- Este \emph{ultimo}. juizo \emph{pratico}, que é determinado \emph{livremente} pela vontade e que deve preceder o acto da eleição, é essencialmente \emph{livre}, é \& razão \emph{adequada} do livre arbitrio, é o proprio \emph{livre arbitrio}, como diz S. Thomaz: u \emph{Arbitrium liberum nihil aliud}. \emph{est}, \emph{quam liberum Judicium} » (\emph{Sum}. \emph{Th}., p. T, q. 88, a. 2). --- Lembramos que o juizo \emph{especulativo é necessario}, mas quando se refere a verdades \emph{evidentes} e não precede a eleição da vontade.}.

D'onde se segue que um acto, para ser \emph{voluntario}, deve derivar não só de uma deliberação da vontade, mas tambem de um previo conhecimento da intelligencia; de modo que o acto da vontade contém tanta bondade ou malicia, quanta foi conhecida pela intelligencia. --- Por isso, se a intelligencia, por qualquer causa ou accidente, fôr perturbada a ponto de não poder apreciar a bondade ou a malicia do acto, este não \pg{218}  é \emph{voluntario}, \emph{e} quem o praticou não é responsavel por elle, nem deve ser punido \footnote{São dois, portanto, os elementos necessarios e indispensaveis para constituir o acto \emph{voluntario}, --- a deliberação da vontade e o previo conhecimento da intelligencia. Se faltar um d'estes elementos, \emph{o} acto não é \emph{voluntario}, e quem o praticou não é responsavel por elle, nem merece castigo.\par A questão da \emph{voluntariedade} dos actos humanos prende-se com uma outra questão de altissima importancia, que nos ultimos tempos tem agitado as escholas criminalistas, e que é a da \emph{responsabilidade moral}. Por quanto, Be 08 nossos actos sao \emph{voluntarios}, isto é, se procedem da vontade deliberada com o previo conhecimento da intelligencia, somos \emph{moralmente} responsaveis por elles e merecemos recompensa ou castigo, --- e se, pelo contrario, os nossos actos não são \emph{voluntarios}, não temos responsabilidade do que praticamos ou deixamos de praticar, e não merecemos premio nem castigo.\par E' evidente que os adversarios do \emph{livre arbitrio} não reconhecem, nem podem reconhecer nenhum acto propriamente \emph{voluntario}; pois que, se a vontade não é dotada de \emph{livre arbitrio}, não póde \emph{deliberar}, e a \emph{deliberação} da vontade é, como dissemos, um elemento essencial do acto \emph{voluntario}.\par Não reconhecendo no homem nenhum acto propriamente \emph{voluntario}, os \emph{deterministas} não podem admittir, como effectivamente não admittem, a \emph{vesponsabilidade moral} dos nossos actos; visto que, se á nossa vontade está sujeita a leis fataes e indeclinaveis, e se todas as nossas acções são consequencia necessaria da constituição organica, do meio physico e social, de doenças hereditarias, não se lhe podem imputar os actos que pratica.\par Que pretendem, pois, os adversarios do \emph{livre arbitrio}?\par A pretensão é excessiva. Pretendem que se fechem as prisões e se abram hospitaes, que a sciencia do direito seja substituida pelo estudo da anatomia e da pathologia, que o jury de magistrados ceda o logar ao jury de medicos.\par Um dos mais modernos defensores d'esta theoria, \emph{Garofalo}, ensina que se devem applicar á ceriminologia os dois principios fundamentaes do darwinismo, --- a \emph{selecção} e a \emph{adaptação}. u E' necessario, diz elle, escolher o elemento que foi organisado para o bem e que mais se adapta ás condições da sociedade, e eliminar todo o elemento que possa ser causa de desordem e de perturbação. Esta eliminação não é um castigo, pois a, responsabilidade \emph{moral} é uma chimera, mas é um preservativo, um meio de defeza, de que\pgn{219}a sociedade se serve para a sua conservação. A segurança publica! Eis o unico motivo que justifica o emprego dos meios coercitivos n».\par E' verdade que muitos criminalistas recuaram deante das perniciosas consequencias, que resultam da negação absoluta do \emph{livre arbitrio}, e, seguindo uma theoria moderada, admittiram que a vontade humana é livre, que os seus actos são \emph{voluntarios}, e por isso merecedores de premio ou de castigo.\par Mas depois fizeram taes e tantas excepções á responsabilidade moral, que, embora admittissem o \emph{livre arbitrio} em theoria, negaram-no de facto.\par D'ahi o enfraquecimento da repressão e o augmento de criminosos, que muitas vezes revestem o crime de horrorosas circunstancias, com o fim de basear no excesso da sua malyadez o argumento dà propria defeza.\par Ha, pois, duas opiniões ácerca da \emph{voluntariedade} dos actos da vontade, que reclamam uma critica imparcial e severa.\par A primeira opinião, que se funda na negação absoluta do \emph{livre arbitrio}, não precisa ser aqui refutada; pois que já o foi no artigo antecedente.\par A segunda opinião, porém, que admitte o \emph{livre arbitrio} em theoria, mas que o nega de facto, merece uma exposição mais ampla; pois, como dissemos, está sendo discutida com muito empenho e calor nas escholas criminalistas, Para isso, investigaremos se, quando e até que ponto, o \emph{livre arbitrio}, e por isso a \emph{voluntariedade} dos nossos actos, póde admittir excepções.\par E' principio assente entre todos os criminalistas que o homem não é responsavel pelo acto que pratica, se este não fôr \emph{voluntario}. Um acto é \emph{voluntario}, se deriva da vontade deliberada com previo conhecimento da intelligencia. Por isso, como dizemos no texto, se a intelligencia, por qualquer causa ou accidente, fôr perturbada a ponto de não poder apreciar a bondade ou a malicia do acto, a vontade, por não receber as indicações da intelligencia, fica tambem perturbada no seu exercicio, não \emph{delibera} ácerca da bondade ou malicia dos objectos, e por isso não póde regular nem reprimir as tendencias das outras faculdades, especialmente do \emph{appetite sensitivo}, e o acto d'ella não é \emph{voluntario}, \emph{e} quem o praticou não é responsavel por elle.\par Mas poderá a intelligencia soffrer uma perturbação, sendo, como é, espiritual, independente da materia? Sim, certamente.\par A intelligencia é uma faculdade espiritual e independente da materia, mas está unida á materia. E' a intelligencia, e não o orgão material, que entende; mas ella não póde operar sem o concurso e o auxilio das faculdades sensitivas, e principalmente da imaginação. A intelligencia, pois, como provámos, depende das faculdades sensitivas e especialmente da \emph{imaginação}.\par D'onde se segue que, se as faculdades sensitivas, e especialmente a imaginação, soffrorem um desarranjo, uma anomalia, este desarranjo, esta anomalia reflectir-se-ha necessariamente na intelligencia, e, por meio da intelligencia;\pgn{220}na vontade; não porque a intelligencia e a vontade fiquem intrinsecamente alteradas, mas porque a intelligencia não percebe o objecto, ou não o percebe como é, e a vontade não póde tomar uma \emph{deliberação}.\par Mas as faculdades sensitivas estão ligadas a certos e determinados pontos do organismo, ao systema nervoso, e por conseguinte as suas operações dependem necessariamente das disposições physicas do sujeito, da sua constituição, do seu estado de saude, e das doenças hereditarias ou adquiridas, que o acommettem, Se, pois, os actos da intelligencia e da vontade exigem o concurso do organismo, é claro que uma perturbação no cerebro, ou um desarranjo no systema nervoso, deve reflectir-se na intelligencia, e, por meio d'esta, na propria vontade.\par Agora podemos comprehender a relação que existe entre o estado do systema nervoso e o acto humano. O systema nervoso influe no exercicio da sensibilidade, esta na intelligencia, e esta na vontade.\par Por isso, se se demonstrar que o systema neryoso soffreu um desarranjo, por fórma que a inteligencia não pôde considerar a qualidade do acto, devemos concluir que este acto não é voluntario, e que por isso o sujeito não é responsavel pelo crime que pôde commetter durante um estado anormal.\par Dão-se, pois, excepções na \emph{voluntariedade} dos actos da nossa vontade, e por isso na responsabilidade moral. Mas quando? e até que ponto?\par Uma resposta cabal e definitiva a estas questões é muito difficil, --- Procuraremos expor alguns criterios, que se encontram nos tratados de \emph{antliropologia criminal}, publicados nestes ultimos annos.\par *\par Os facto principaes, que podem occasionar a \emph{énvoluntariedade} dos nossos actos, dividem-se em \emph{pathologicos} e \emph{physiologicos}, segundo revelam um estado \emph{pathologico} ou um estado \emph{physiologico} defeituoso.\par A) Aos factos \emph{pathologicos} costumam reduzir a \emph{loucura}, a \emph{epilepsia}, O \emph{hysterismo}, as \emph{paixões}, o \emph{alcoolismo}, e o \emph{somnambulismo}.\par a) A \emph{loucura} é um desarranjo da intelligencia, o qual deriva da alteração na \emph{imaginação}. --- Como já dissemos, dão-se alterações na \emph{imaginação}, --- ou quando a substancia do cerebro se acha affectada por inflammação, pela acção do alcool, etc.; --- ou quando a circulação do sangue no cerebro soffre um augmento demasiado, ou uma demasiada, diminuição, --- ou quando as funcções do systema nervoso estão extraordinariamente perturbadas, como nas paixões violentas.\par A \emph{loucura} divide-se em \emph{geral}, chamada \emph{mania}, e em \emph{parcial}, conhecida 80b o nome de \emph{monomania}. --- A \emph{mania} é um desarranjo \emph{completo} da intelligencia. A \emph{monomania} é um desarranjo \emph{parcial} da intelligencia.\par As principaes formas da \emph{monomania} são quatro: 13) monomania de\pgn{221}\emph{grandezas}, quando predominam ideas grandes; --- 2.º) monomania \emph{religiosa}, quando predominam ideas relativas á religião; --- 8.º) monomania de \emph{perseguição}, quando a pessoa tem a illusio de que é perseguida; --- 4.º) monomania \emph{moral} ou \emph{impulsiva}, quando se sente um impulso cego e irresistivel para o crime. \textbar{} Posto isto, pergunta-se: O acto do \emph{louco} ser\& \emph{voluntario}?\par Quando a loucura é \emph{geral}, é certo que o acto do \emph{louco} não é \emph{voluntario}.\par Quando a loucura é \emph{parcial}, o caso é mais complicado. O \emph{monomantaco} num só ponto mostra desarranjo mental; em tudo o mais é rasoavel, opéra com deliberação, sabe o que faz. Serão \emph{voluntarios} os seus actos?\par Os sabios apresentam duas soluções d'este problema.\par Uns dizem que nos \emph{monomaniacos} são \emph{voluntarios} todos os actos, que estão fóra da esphera das ideas delirantes; mas que são \emph{énvoluntarios} todos os que são consequencia do delirio.\par Outros sustentam que \emph{todos} os actos dos \emph{monomaniacos} são \emph{involuntarios}.\par Ambas estas soluções têm contra si serias difficuldades.\par A primeira solução é excellente em theoria, mas é difficil na pratica. Na verdade, quem poderá estabelecer com certeza o termo, onde acaba a zona do estado são, e onde principia a zona do delirio? O cerebro dos \emph{monomaniacos} esta fóra das condições normaes. Com que fundamento, pois, poderá abtribuir-se uma operação á intelligencia intacta, e outra ao delirio? « Nos \emph{monomaniacos}, diz \emph{Tardieu} no seu \emph{Estudo imedico-legal sobre a loucura}, a idea predominante destaca-se sobre um fundo geral e primitivamente alterado, e o delirio parcial é a nota mais aguda do profundo desaccordo, que existe entre as differentes funcções intellectuaes e moraes ».\par A segunda solução, que admitte a \emph{Znvoluntariedade} em \emph{todos} os actos dos \emph{qmonomaniacos}, toca outro extremo. « Quantos criminosos, diz \emph{Vacherot}, que a physiologia deu por \emph{monomaniacos}, tinham uma comprehensão do plano, um calculo reflectido dos meios, e um grande sangue frio na execução: o que mostra que sabiam o que faziam .\par Diremos que é necessario estndar a vida e os actos do \emph{monomaniaco}; pois que este estudo nos dará o criterio sufficiente para julgarmos do seu estado mental, e por isso da \emph{voluntariedade} dos seus actos.\par Alguns criminalistas, além da loucura \emph{intellectual}, de que temos fallado até agora, admittem uma outra especie de loucura --- a loucura \emph{moral}, que \emph{Pinel} chamou \emph{mania sem delirio}, e \emph{Esquirol} denominou \emph{monomania instinctiva}. A loucura \emph{moral} deixaria sã a intelligencia, e só atacaria a vontade. \emph{T'relat} diz que os loucos \emph{moraes} têm a consciencia de si mesmos, mas são arrastados necessariamente para o acto.\par Deverá admittir-se esta loucura \emph{moral}, esta \emph{mania sem delirio}, esta, monomania \emph{instinctiva}? --- Não; porque os factos não nos habilitam a isso. Quem nos póde dizer com certeza que a intelligencia de taes loucos, se existem\pgn{222}, ficou intacta? E poderá haver uma perturbação na vontade, sem: uma previa perturbação na intelligencia? Certamente que não. Diz \emph{Moreau}, de Tours: « A lesão na vontade, a irresistibilidade nas determinações instinetivas, sem uma lesão na intelligencia, é uma chimera ».\par Diga-se o mesmo com relação a dois outros factos, que foram inventados para negar a \emph{voluntariedade} dos actos humanos, e que são --- a \emph{incubação da loucura} e a \emph{força irresistivel}, Certas asserções refutam-se com uma simples negação. Se fossem verdadeiras semelhantes theorias, não haveria nenhum acto \emph{voluntario} e por isso nenhum criminoso seria responsavel, porque todos poderiam appellar para a \emph{loucura em estado de incubação}, para a \emph{força irresistivel}, sem receio de serem desmentidos.\par b) A \emph{epilepsia} é uma nevrose convulsiva. Não tem caracteres proprios e definidos; apresenta innumeras formas e variedades.\par Os ceriminalistas discutem se a \emph{epilepsia} é distincta, ou não, da \emph{loucura}. \emph{Trousseau} diz que é distincta; \emph{Lombroso} affirma que não; \emph{Ferri} admitte a identidade entre a \emph{epilepsia e} \& monomania \emph{instinctiva}.\par Serão \emph{voluntarios} os actos do \emph{epileptico} ? 3\par Respondemos que é necessario distinguir. Ou o \emph{epileptico} soffre tambem de loucura, ou não. --- No primeiro caso, devemos julgar da \emph{voluntariedade} de seus actos com aquelles mesmos criterios, com que se julga da \emph{voluntariedade} dos actos dos loucos. --- No segundo caso, são \emph{voluntarios} os actos que pratica nos intervallos dos ataques, em que tem consciencia de si; mas são \emph{involuntarios} os actos que pratica durante o ataque, em que não possue o uso da razão. --- Advertimos com \emph{Legrand}. \emph{du Saulle} que o epileptico, durante o ataque, está insensivel ao mundo exterior, e é incapaz de commetter um crime.\par A \emph{epilepsia} não se apresenta sempre com signaes manifestos, com ataques periodicos, que todos possam observar. Ás Vezes, as suas crises dão-se durante a noite, \emph{e} nem sempre são advertidas pelo proprio doente; manifesta-se só por certos desarranjos mentaes, por certos symptomas extranhos com caracteres vagos e indefinidos. E' a epilepsia \emph{larvada}.\par Serão \emph{voluntarios} os actos do epileptico \emph{larvado} ?\par A \emph{voluntariedade} dos actos d'este doente não póde precisar-se, se não se conhecer o estado da sua intelligencia. Mas tal conhecimento é muito difficil se não impossivel; porque, não existindo o typo do epileptico \emph{larvado}, não póde determinar-se um criterio absoluto e invariavel para todos os casos. Diremos com o dr. \emph{Motet} que, na pratica, é necessario usar de muita cautela, e analysar severa e conscenciosamente os actos de quem se diz epileptico \emph{lareado}, para se julgar do grau da sua \emph{voluntariedade}.\par c) O \emph{Aysterismo}, segundo \emph{Legrand du Saulle}, admitte tres graduações.\par Umas vezes, esta doença não impede de modo algum o uso da razão, nem o exercicio do livre arbitrio, e nestes casos não tira a \emph{voluntariedade} dos actos; outras vezes, tolda a luz da razão, mas não a supprime inteiramente,\pgn{223}e então diminue a \emph{voluntariedade}; --- algumas vezes, porém, chega a tirar completamente o uso da razão, e nestes casos os actos da Pessoa \emph{Aysterica} não são \emph{voluntarios}.\par Toda a difficuldade, pois, consiste em classificar um caso concreto de \emph{hysterismo}, Para essa classificação não existe um criterio seguro, uma formula geral. E' necessario examinar cada um dos Casos, em si e nas suas cireumstancias, e em seguida formar o juizo do estado mental e da voluntariedade e responsabilidade dos actos da pessoa hysterica.\par d) Às \emph{paixües} são vohementes tendencias do appetite sensitivo para o bem sensivel e material.\par E' certo que as \emph{paixões} podem levar e muitas vezes levam ao crime; mas exercerao ellas uma pressão fatal e irresistivel sobre a nossa vontade, de modo que os actos d'esta faculdade deixem de ser \emph{voluntarios}?\par \emph{Erlenmeyer} diz que o medo e o desespero tém o caracter fundamental de um ataque subitaneo de loucura. \emph{Caspier} sustenta que sob a acção de uma paixão violenta não ha liberdade moral.\par A verdade é que o acto praticado sob o impulso da \emph{paixão é geralmente} mais ou menos \emph{voluntario}, pois a vontade póde e deve resistir ás paixões.\par Dissemos --- \emph{geralmente}; porque, se a paixão for tão violenta, que chegue a produzir uma perturbação mental e a impedir o exercicio regular da intelligencia, então o acto da vontade não é \emph{voluntario}.\par e) O \emph{alcoolismo}, segundo demonstram as estatisticas, exerce uma grande influencia no crime e na loucura. Os crimes abundam notavelmente nos paizes, onde é maior o consumo de bebidas alcoolicas. Na França, sobre 100 casos de loucura, 18 são consequencia do abuso do alcool; na Inglaterra, metade dos alienados eram dados ao alcoolismo.\par Será \emph{voluntario} o acto do \emph{alcoolico} ?\par E' necessario saber se o acto foi praticado --- por uma pessoa em estado de embriaguez, --- ou por um \emph{alcoolico chronico}, mas fôra d'aquelle estado.\par Na primeira hypothese, so á embriaguez tirou completamente o uso da razão e o aleoolico não previu o effeito do seu acto, este não é \emph{voluntario}. --- Be, porém, a embriaguez não supprimiu inteiramente o uso da razão, ou o aleoolico previu o effeito do seu acto, então este é \emph{voluntario,e} a sua \emph{voluntariedade} é proporcionada ao grau do conhecimento ou da previsao.\par Na segunda hypothese, a solução é mais difficil e complicada. \emph{Roussel} diz que a alteração da liberdade nos \emph{alcoolicos chronicos} é continua e progressiva. \emph{Bergeron} affirma que o \emph{alcoolico chronico} é um doente de corpo e de espirito, um louco que perdeu a liberdade e a responsabilidade.\par Estas asserções tem um certo fundo de verdade. Os vinhos o os licores artificiaes, como provam \emph{JMagnan} e \emph{Laborde}, são na maxima parte bebidas toxiecadas. Os aleoolicos envenenam-se lentamente. Comecam por experimentar movimentos convulsivos, ataques epilepticos, desarranjos no apparelho respiratorio, congestóes cerebraes, até ao \emph{delirium tremens}. Depois de algum tempo, variavel segundo a qualidade e a quantidade dos liquidos absorvidos e segundo a, constituição e as forças do sujeito, os desarranjos accentuam-se,\pgn{224}e o alcoolismo pôde degenerar, e muitas vezes degenera em loucura ou em paralysia geral. Nestas circumstancias, em que a intelligencia do \emph{alcoolico} não póde attender á malicia do acto, este não é \emph{voluntario}.\par £) O \emph{somnambulismo} tambem é adduzido como excepção á \emph{voluntariedade} dos nossos actos e por isso á responsabilidade moral.\par Será \emph{voluntario} o acto do \emph{somnambulo}?\par Ha duas especies de \emph{somnambulismo}, --- um \emph{natural}, e outro \emph{artificial}. --- O \emph{somnambulismo artificial} chama-se commummente \emph{hypnotismo}.\par Quando o somnambulismo é \emph{natural}, se os factos demonstrarem que o somuaimbulo praticou um acto, achando-se sob o impulso do automatismo somnambolieo, o seu acto não é \emph{voluntario}; pois que, sob a influencia do somno magnetico, como dizem \emph{Devergie}, \emph{Tardieu}, \emph{Laségne}, \emph{Charcot} e \emph{Vulpian}, o exercicio da vontade não é livre. ;\par Quando, porém, o somnambulismo é \emph{artificial}, quando se trata do \emph{hypnotismo}, a coisa muda de aspecto. O hypnotizado, em certos casos, póde ser culpado por se ter prestado a soffrer a influencia de um agente extranho, a qual impede o uso do livre arbitrio, e póde tambem ser responsavel pelos actos que praticou durante aquelle estado, se os previu.\par Todavia, nem sempre, como adverte \emph{Figuier}, deve o hypnotizado considerar-se um automato, que obedece a um impulso exterior; aliás nenhum criminoso poderia ser condemnado, e cada um poderia desculpar-se, dizendo que foi victima do odio de outra pessoa, a qual quiz vingar-se d'elle, hypnotizando-o e arrastando-o assim necessariamente para o crime.\par B) Aos actos \emph{plysiologicos} subordinam a \emph{hereditariedade}, a \emph{influencia} do «\emph{ele social} e da \emph{educação}, o \emph{temperamento} e o \emph{criminoso-nato}.\par g) A \emph{hereditariedade} é um dos factos \emph{physiologicos}, em que muitas vezes se baséam os argumentos da defeza.\par E' certo que a gotta, á tuberculose, a nevrose, a loucura, a epilepsia, e, em geral, as doenças \emph{diatAesicas} se transmittem, embora, ás vezes, modificadas e transformadas.\par Mas poderá a \emph{hereditariedade} ser causa de \emph{involuntariedade} ?\par Antes de darmos uma resposta categorica, devemos advertir que a \emph{hereditariedade} não transmitte ao filho o proprio estado pathologico dos paes, mas sim uma certa disposição ou propensão para reproduzir aquelle estado. Tal disposição ou propensão póde ser combatida e attenuada, e muitas vezes é inteiramente vencida. --- De mais, a doença dos paes póde ter principiado depois do nascimento do filho, ou póde ter sido effeito de causas fortuitas o \&ccidentaes, por exemplo, de uma lesão nos centros nervosos, motivada por uma quéda ou por um golpe recebido. .\par Posto isto; diremos que, para a \emph{Hereditariedade} poder ser tida como causa da \emph{énvoluntariedade} dos actos de um criminoso, não basta provar que os seus paes soffreram de loucura, ou de qualquer outra doença (porque, como dissemos, a hereditariedade não é uma lei geral, que pese fatalmente sobre os d.scendentes) mas é necessario provar que o criminoso está sujeito\par ,\pgn{225}ao mesmo estado pathologico, a que estiveram sujeitos seus paes. --- Se não. se provar que o criminoso está sujeito ao estado pathologico, a que estiveram sujeitos os paes, a \emph{hereditariedade} não pôde invocar-se como causa da \emph{involuntariedade} dos actos do mesmo criminoso.\par Uma manifestação da hereditariedade é o \emph{atavismo}. Quando uma criança reproduz as feições e o caracter dos antepassados mais remotos, este retrocesso ao typo antigo chama-se \emph{atavismo}.\par E' certo que ha casos de \emph{atavismo}; mas esses casos são raros, muito irregulares e inconstantes, e não ha uma regra geral, uma formula fixa, para julgar da \emph{voluntariedade} ou \emph{involuntariedade} dos actos dos \emph{atavicos}. O criterio mais seguro e prudente é estudar o caso particular, e resolyel-o á luz dos principios anteriormente expostos.\par Nào podemos deixar de advertir que, para os materialistas, o \emph{atavismo} é o retrocesso aos instinctos bestiaes e ferozes d'aquelles animaes que, segundo elles, teriam sido os antepassados do homem. --- Não insistiremos na refutação de tão absurda como perniciosa doutrina, que faz do homem um descendente de uma especie inferior; pois, na \emph{Cosmologia especial} e na \emph{Secção quarta} d'este tratado, procuràámos e procuraremos demonstrar que aquella doutrina é contraria á razão e á experiencia.\par h) A \emph{influencia} do \emph{meio social} e da \emph{educação} é outro facto \emph{physiologico}, invocado muitas vezes como causa da \emph{involuntariedade}.\par O dr. \emph{Bordier} diz claramente que a \emph{educação} exerce uma pressão fatal sobre a vontade. --- \emph{Tine} affirma o mesmo relativamente ao \emph{meio social}.\par O erro d'estes escriptores é manifesto. Concordamos em que o \emph{meio social} e a \emph{educação} exercem uma grande influencia sobre a nossa vontade, principalmente por causa dos respeitos humanos; mas, apoiados no testemunho da consciencia, negamos que esta influencia seja fatal. Quem quer, resiste a todas as influencias do meio em que vive, e da educação que recebeu.\par i) O \emph{temperamento} será causa da \emph{énvoluntariedade} dos nossos actos?\par \emph{Temperamento} é o predominio de um certo grupo de paixões nos homens, proveniente da disposição do organismo, e, sobretudo, do systema nervoso.\par Os physiologistas não estão de accordo com relação ás especies de \emph{temperamento}. Wundt (\emph{Principios de psychologia physiologica}), seguindo \emph{Galeno}, admitte quatro especies: \emph{sanguineo}, \emph{bilioso}, \emph{melancholico}, \emph{fleugmatico}.\par \emph{a}) No temperamento \emph{sanguineo}, proveniente da rapida e abundante circulação do sangue, predominam as paixões do appetite \emph{concupiscivel}, que são o amor, o \emph{desejo} e a \emph{alegria}. --- Este temperamento, conforme a sua direcção, ás vezes produz um caracter franco, amoravel e expansivo; mas outras \textbf{vezes} dá um caracter inclinado ás desordens.\par \emph{b}) No temperamento \emph{Dilioso} predominam as paixões do appetite \emph{érascivel}, que levantam o animo, como são a \emph{esperança}, a \emph{audacia}, \& \emph{ira}. --- Este temperamento produz um caracter altivo ou iracundo.\par \emph{c}) No temperamento \emph{melancholico} predominam as paixões \emph{repulsivas} do appetite \emph{concupiscivel}, e as paixões \emph{deprimentes} do appetite \emph{érascivel}, que são\par 15 Vor. II.\pgn{226}o \emph{desespero} e o \emph{medo}. --- Este temperamento produz um caracter doce, manso e respeitoso, ou um caracter sombrio, intratavel, mau.\par \emph{d}) No temperamento /\emph{lewgiatico}, as paixões não são intensas, mas acham-se equilibradas. Este temperamento transforma-se ou num caracter reflexivo, prudente e commedido, ou num caracter apathico, sceptico, egoista.\par A maior parte dos modernos dividem os \emph{temperamentos} em \emph{sanguineo}, \emph{nervoso}, \emph{lymphatico} e \emph{mixto}, conforme predonima nos individuos uma riqueza de desenvolvimento de vasos sanguineos e sangue, --- ou o desenvolvimento do systema nervoso ou do systema lymphatico, ou se associam alguns d'estes caracteres para determinar formas mixtas, que correspondem geralmente aos que poderemos designar \emph{equilibrados de temperamento}.\par Ora, quer com a divisão acima desenvolvida quer com esta outra, poderemos sempre dizer que o \emph{temperamento}, embora ás vezes torne mais facil o exercicio de certas virtudes e mais diffieil o de outras, comtudo, geralmente, não tira o \emph{livre arbitrio} e por isso a \emph{voluntamiedade} dos nossos actos; pois a vontade póde corrigir o \emph{temperamento}, subordinando-o á razão. - ' Dissemos --- \emph{geralmente}; porque o \emph{temperamento}, se produzir uma perturbação profunda ma intelligencia, impedirá tambem o exercicio do \emph{livre arbitrio}, e os actos da vontade não serão \emph{voluntarios}.\par jd) O \emph{eriminoso-nato} é o facto mais saliente, a razão mais poderosa que hoje se adduz em favor da \emph{énvoluntariedade} de certos actos. Diz-se que o homem, que \emph{nasceu criminoso}, não delibera, não escolhe, mas é arrastado irresistivelmente para o mal, e por isso não é responsavel pelos seus actos.\par \emph{Topinard}, \emph{Bordier}, \emph{Lacassagne}, \emph{Magitot}, na França, e \emph{Lombroso}, \emph{Virgilio}, \emph{Ferri}, na Italia, procuraram arranjar um typo do \emph{criminoso-nato}, para abaterem a \emph{gothica fortaleza do livre arbitrio} (Letourneau). Para isso, recorreram á \emph{anthropologia} e á. \emph{craneologia}, embora esta ultima tivesse sido abandonada depois das inuteis tentativas de \emph{Gall}. E arranjaram esse typo; mas sahiu-Ihes tão exotico e contrafeito, que os sabios acabaram por se convencer de que elle nunca existiu nem póde existir.\par E, se não, vejamos.\par \emph{Lowmbroso} enumera um conjuncto de caracteres, que seriam proprios do \emph{eriminoso-nato}, e distinguil-o-hiam do homem de bem, --- caracteres \emph{anatomicos}, \emph{plysiologicos}, \emph{pathologicos} e \emph{psycologicos}. --- \emph{Garofalo}, pelo contrario, pensa que não existe nenhum caracter physico, que distinga o criminoso de quem o não é; porque certas anomalias, que se encontram no criminoso, encontram-se no homem de bem. --- A opinião de \emph{Garofalo} é seguida por \emph{Benedickt}, que affirma ser impossivel encontrar no criminoso symptomas anthropologicos certos de criminalidade.\par A divergencia entre os materialistas \&ecentua-se, quando se trata de indicar os caracteres \emph{p/ysicos} do \emph{eriminoso-nato}.\par Em relação á capacidade do craneo do eriminoso, \emph{Lombroso} diz que essa capacidade é inferior á media, pois é de 1,4150e.; \emph{Bordier}, examinando trinta e seis craneos de criminosos guilhotinados, diz que a sua capacidade é superior\pgn{227} á media, pois é de 1,541e.; \emph{Ranche} sustenta que é egual á media; \emph{Magitot} afirma que o craneo do criminoso, em capacidade, não differe do craneo de qualquer homen de bem.\par A mesma divergencia existe, emquanto ao desenvolvimento do craneo, emquanto ao peso, estatura e cór do criminoso. --- Segundo \emph{Zeger}, \emph{Dallemagne}, \emph{Bordier}, \emph{Bagenoff}, no criminoso a circumferencia craneana posterior é mais desenvolvida; segundo \emph{JMarro}, dá-se o contrario. --- Para \emph{Lombroso}, o criminoso é grosso e alto; para \emph{Virgilio} e \emph{Towpson}, não é nem grosso nem alto. --- Segundo a eschola italiana, o criminoso é mais trigueiro do que loiro; segundo a eschola allemá e sueca, é mais loiro do que trigueiro.\par \emph{Lombroso} assignala aos criminosos certos modos especiaes de escrever; \emph{Laurent} não os encontrou. --- \emph{Lombroso} diz que os criminosos apresentam uma conformação caracteristica nas orelhas; \emph{Lannois} diz o contrario.\par \emph{Lombroso} nota uma falta de symetria no cerebro dos criminosos; mas havia falta de symetria nos cerebros de \emph{Adolpho Thiers} e do naturalista \emph{Bichat}, e todavia estes dois sabios não eram criminosos.\par Se grande é a divergencia entre os escriptores materialistas, quando \textbf{apresentam} os caracteres \emph{phiysicos} do \emph{criminoso-nato}, não é menor, quando tratam de assignar os seus caracteres \emph{psychicos}.\par Para uns, o \emph{criminoso-nato} é timido, para outros corajoso; para uns é devoto, para outros impio; para uns é circumspecto, para outros arrojado.\par \emph{Lombroso} enumera, entre outros caracteres, á \emph{éncorrigibilidade}. --- Mas foram sempre empregados todos os meios possiveis para a regeneração do delinquente? Por certo que não. E, se o não foram, com que razão poderá concluir-se que o criminoso é \emph{incorrigivel}?\par Não pôde negar-se que alguns criminosos apresentam certos caracteres maus; mas estes caracteres, longe de serem causa do crime, são geralmente uma consequencia d'elle.\par Poderiamos continuar esta exposição comparada, mas basta o que fica dicto.\par Concluimos que a moderna theoria ácerca do \emph{criminoso-nato}, baseando-se no incerto, no hypothetico e no contradictorio, deve ser rejeitada não só pelos espiritos de criterio são, mas tambem pelos proprios materialistas, se quizerem ser coherentes com o systema que professam, e que só admitte o \emph{positivo}. --- E aqui não podemos deixar de applaudir a decisão do Congresso anthropologico de \emph{Bruxelles}, celebrado em agosto de 1892, que declarou \emph{gratuita} e \emph{antiscientifica} \& doutrina de \emph{Lombroso}.\par Quando, pois, o advogado da defeza affirmar que o seu cliente apresenta os caracteres que constituem o typo do \emph{criminoso-nato}, esta affirmacüo não deve ser attendida; pois, como vimos, o typo do \emph{criminoso-nato} não existe.\par Os defensores do \emph{criméinoso-nato}, como tambem os adversarios do \emph{livre arbitrio}, dizem: A sciencia deve assentar unica e exclusivamente nos factos, prescindindo da sua interpretação. Ora os factos dizem que o homem não é dotado do \emph{livre arbitrio} e que o typo do \emph{eriminoso-nato} existe.\par Esta pretensão é \emph{contradictoria} e \emph{falsa}.\par a) E' \emph{contradictoria}. Porquanto, esses escriptores, ao passo que censuram os defensores do \emph{livre arbitrio}, por estes interpretarem os factos e \textbf{fazerem}\pgn{228}questões de \emph{metaphysica}, não deixam pela sua parte de fazer taes interpretações e questões de \emph{metaphysica}.\par \emph{b}) E' \emph{falsa}. Com effeito, o methodo scientifico, como diz \emph{Cl}. \emph{Bernard}, não consiste na simples enumeração dos factos, mas na legitima interpretação d'elles, pois que as sciencias progridem não pela accumulação dos factos, mas pela \emph{luz da idea}, que \emph{contém} a explicação dos proprios \emph{factos}.\par Citaremos ainda a este respeito o testemunho de um outro escriptor, o mais insuspeito que se pode imaginar ---- \emph{Haeckel}. E' o seguinte:\par \& Os naturalistas da eschola exclusivamente empirica, dominados por uma extranha illusao, lisonjéam-se de poder construir o edificio inteiro da historia natural com factos isolados, sem os ligar philosophicamente entre Si, e com simples noções isoladas, sem lhes perceber o sentido.\par « Sem duvida, todo o systema puramente especulativo, absolutamente philosophico, que não se basêa no fundamento inabalavel dos factos empiricos, é um mero eastello de cartas, que a primeira experiencia lançará por terra; mas tambem toda a obra scientifica, puramente empirica, composta unicamente de factos, será, quando muito, um montão de pedras, nunca um edificio. Os factos são por si \emph{arídos}, e, como os apresenta a experiencia, são apenas materiaes grosseiros; \emph{se não forem fecundados pelo pensamento e ligados pela philosophia}, \emph{nunca poderão constituir una sciencia} . . .\par u Um lamentavel antagonismo entre as sciencias naturaes o a philosophia, um empirismo grosseiro, que a maior parte dos naturalistas contemporaneos desgraoadamente proclama como \emph{sciencia exacta}, --- aos são as causas de tantas e tão grandes faltas de logica, e da impotencia absoluta para tirar dos factos as conclusões mais elementares...\par \& Eis a que leva o desleixo da cultura philosophica, e da educação do espirito » (\emph{Histoire de la création} --- \emph{Inconvénients de l'empirisme}, pag. 694).\par Mas, embora concedamos que a sciencia reside exclusivamente nos factos, todavia os adversarios nunca poderão encontrar factos, capazes de destruir o \emph{livre arbitrio} e de provar a existencia do \emph{criminoso-nato}.\par * *\par Resumimos estes apontamentos nos termos seguintes.\par Em theoria: O homem é dotado de livre arbitrio, e por isso os seus actos são \emph{voluntarios}. Todavia ha excepções, que ás vezes tiram toda a \emph{voluntariedade}, e outras vezes apenas a diminuem.\par Na pratica: Quando uma pessoa commette um crime, e não consta que houvesse uma d'estas circumstancias excepoionaes, deve imputar-se-lhe esse acto e deve ser punido, porque é a regra que se presume, e é a excepção que cumpre provar-se. --- Se constar que houve excepção, que se deu um dos factos \emph{pathologicos} ou \emph{physiologicos}, que enumerámos, deve averiguar-se se o exercicio do livre arbitrio ficou completamente destruido, ou apenas attenuado. No primeiro caso, o criminoso, por ser irresponsavel, não merece castigo; no segundo, por ser responsavel em parte, devem reconhecer-se-lhes as \emph{circumstancias altenuantes}.}.

%%>>> t142

\pg{217} \tese{142}{A vontade move a intelligencia.}
Uma faculdade move outra, quando aquella determina esta a produzir ou a deixar de produzir um certo acto. Ora a vontade determina a intelligencia a considerar ou a deixar de considerar esta ou aquella coisa. Logo a vontade move a intelligencia.

Todavia, se o exercicio da vontade presuppõe sempre o da intelligencia, o exercicio d'esta não presuppõe sempre o d'aquella faculdade (aliás admittir-se-hia um processo indefinido, o que é absurdo); mas a intelligencia é movida, em ultima analyse, pela Causa Primeira de todo o movimento, emquanto é a Verdade.

A mutua dependencia, pois, que existe entre a intelligencia e a vontade, é de natureza diversa. Com effeito, a intelligencia move a vontade emquanto á \emph{especificação}, pois que apresenta o objecto, para o qual a vontade deve ou pode tender, e o objecto \emph{especifica} as faculdades, --- ao passo que a vontade move a intelligencia emquanto ao \emph{exercicio}, porque a determina a produzir ou a deixar de produzir o proprio acto.

%%>>> t143

\pg{219} \pg{220} \pg{221} \tese{143}{A intelligenocia é mais nobre que a vontade.}

\&) Uma faculdade é mais nobre que outra, quando esta tem a razão proxima da sua existencia naquella. Ora a \emph{vontade} tem a razão proxima da sua existencia na \emph{intelligencia}; pois não teriamos tendencia para um objecto espiritual, se este não fosse anteriormente conhecido pela \emph{intelligencia}. Logo \pg{222}  a \emph{intelligencia} é mais nobre que a \emph{vontade}. --- Por isso a \emph{von}-

\emph{Z tade} diz-se faculdade \emph{intellectual}, porque nasce \emph{immediatamente} da intelligencia, e só \emph{mediatamente} da \emph{essencia} da alma.

b) Uma faculdade é mais nobre que outra, quando a operação d'aquella é mais nobre que a d'esta; pois que as faculdades \pg{223}  especificam-se pelas suas operações. A operação da \emph{intelligencia} é mais nobre que a da \emph{vontade}. Por quanto, a intelligencia, quando entende, attrahe a si um objecto exterior, communieando-lhe uma existencia ideal ou um ser intelligivel \pg{224} ; ao passo que a vontade, quando tende para o objecto, sahe de si mesma, tornando-se semelhante ao proprio objecto. Ora uma faculdade que attrahe a si os objectos exteriores, \pg{225}  tornando-os semelhantes á propria essencia, é, em geral, mais nobre e mais activa do que a faculdade que é attrahida por objectos exteriores, tornando-se semelhante a elles. \textbf{Logo} a \emph{intelligencia} é mais nobre que a \emph{vontade}.

c) Uma faculdade é mais nobre que outra, se o objecto d'aquella fôr mais abstracto e universal que o d'esta; porque as. faculdades especifieam-se pelos objectos, e um objecto é \textbf{mais} ou menos nobre, segundo o seu maior ou menor grau \pg{227}  de abstracção. O objecto da \emph{ntelligencia} é mais abstracto e universal do que o da \emph{vontade}. Com effeito, a \emph{intelligencia}

\emph{Z} apprehende o ente, emquanto é dotado de uma certa \emph{essencia} \pg{228}  (pois o objecto \emph{proprio} da \emph{intelligencia} é a \emph{essencia}), --- ao passo que a \emph{vontade} tende para o objecto, emquanto é dotado de \emph{existencia real} (pois um objecto, emquanto \emph{existe} na realidade, é bem, e por isso é termo da \emph{vontade}). Ora o ente considerado \pg{229}  na sua \emph{essencia} é mais abstracto e universal que o ente considerado na sua \emph{existencia real}; pois que a \emph{essencia} não inclue a \emph{existencia real}, mas esta inclue aquella. Logo a \emph{intelligencia} é mais nobre que a \emph{vontade} \footnote{A intelligencia, \emph{absolutamente} fallando, é mais nobre que a vontade; porque o objecto \emph{formal} d'aquella é mais universal que o d'esta.\par Mas, se attendermos aos objectos \emph{materiaes} das duas faculdades, umas vezes a intelligencia é superior á vontade, e outras esta é superior áquella.\par Quando o objecto é inferior ao espirito humano, a intelligencia é superior á vontade; assim conhecer uma pedra é mais perfeito do que amal-a.\par Quando, pelo contrario, o objecto é superior ao espirito humano, então a vontade é superior á intelligencia; assim o amor, que nos une a Deus, é mais perfeito e mais nobre do que o simples conhecimento do Ente Infinito. \emph{A} razão é clara. Na primeira hypothese, o objecto, que é inferior ao espirito humano, recebe da inteligencia uma forma mais nobre do que a que existe no proprio objecto e que attrahe a vontade; assim o ser ideal, que a pedra adquire no nosso espirito, é mais nobre que o ser material que é proprio d'ella: --- como, pelo contrario, o objecto, que é superior ao nosso espirito, existe em si mesmo de um modo mais nobre do que existe ha nossa intelligencia; assim a Bondade de \textsc{Deus}, que é o objecto do nosso amor, é infinitamente mais nobre e elevada na Essencia Divina, do que no nosso espirito, que a conhece (Cf. \emph{Sum}. \emph{Theol}, p. I, q. 82, a. 3).\par Concluimos que, \emph{absolutamente} fallando, a intelligencia é mais nobre que a vontade; porque as faculdades especificam-se pelos seus objectos \emph{formaes}, e não pelos seus objectos \emph{materiaes}.\par Todavia, se a inteligencia, pela sua elevação e extensão, é mais nobre que a vontade, esta faculdade é mais intima, mais central, mais profunda e por isso mais secreta e mysteriosa do que aquella. D'ahi a expressão commum a quasi todas as linguas: eu te amo \emph{do fundo do coração}. Por isso é mais facil conhecer os pensamentos da intelligencia, do que os affectos do coração. \emph{Larochefowcauld} dizia: « Nem todos os que conhecem o seu espirito, conhecem o seu coração ».}.

\pg{226}
